% ==== XIV ====
\yoschapter{XIV}{The Integration of the Equation of Evolution*}
\label{chap:XIV}
\yosstarfootnotetext{See also Supplementary Notes, p.~\pageref{notes:IX-XIV}.}

The ordinary exponential function solves the initial value problem 
\[
\dd y / \dd x = \alpha y, y (0) = 1.
\]

We consider the diffusion equation 
\[
\partial u/\partial t=\Delta u,
\]
where $\Delta=\sum_{j=1}^m\partial^2/\partial x_j^2$ is the Laplacian in $\mathbb R^m$.

We wish to find a solution $u = u(x, t)$, $t > 0$, of this equation satisfying the initial condition $u(x, 0) = f(x)$, where $f(x) = f(x_{1}, \ldots, x_{m})$ is a given function of $x$. We shall also study the wave equation 
\[
\partial^{2} u / \partial t ^{2} = \Delta u, - \infty < t < \infty ,
\]
with the initial data 
\[
u (x, 0) = f (x) \quad \text{ and } \quad (\partial u / \partial t) _{t=0} = g (x),
\]
$f$ and $g$ being given functions. This may be written in vector form as follows: 
\[
\frac {\partial}{\partial t} \begin{pmatrix}u \\ v\end{pmatrix} = \begin{pmatrix}0 & I \\ \Delta & 0\end{pmatrix} \begin{pmatrix}u \\ v\end{pmatrix}, v = \frac {\partial u}{\partial t}
\]
with the initial condition 
\[
\begin{pmatrix}u (x, 0) \\ v (x, 0)\end{pmatrix} = \begin{pmatrix}f (x) \\ g (x)\end{pmatrix}.
\]

So in a suitable function space, the wave equation is of the same form as the diffusion (or heat) equation---differentiation with respect to the time parameter on the left and another operator on the right---or again similar to the equation $\dd y/\dd t = \alpha y$. Since the solution in the last case is the exponential function, it is suggested that the heat equation and the wave equation may be solved by properly defining the exponential functions of the operators 
\[
\Delta \text{ and } \begin{pmatrix}0 & I \\ \Delta & 0\end{pmatrix}
\]
in suitable function spaces. This is the motivation for the application of the semi-group theory to Cauchy's problem. It is to be noted that the 

Schrödinger equation 
$i^{-1}\partial u / \partial t = Hu = (\Delta +U(x))u$, where $U(x)$ is a given function, 
gives another example of the equation of evolution of the form 
\begin{equation}\label{eq:XIV.0.1}
\partial u / \partial t = A u, t > 0,
\end{equation}
where $A$ is a not necessarily continuous, linear operator in a function space. 

The equation of the form \eqref{eq:XIV.0.1} may be called a temporally homogeneous equation of evolution. We may integrate such an equation by the semi-group theory. In the following three sections, we shall give typical examples of such an integration. We shall then expound the integration theory of the temporally inhomogeneous equation of evolution 
\begin{equation}\label{eq:XIV.0.2}
\partial u / \partial t = A (t) u, a < t < b.
\end{equation}


\yossection{1}{Integration of Diffusion Equations in \texorpdfstring{$L^{2}(\mathbb{R}^{m})$}{L2(Rm)}}
\label{sec:XIV.1}

Consider a diffusion equation 
\begin{equation}\label{eq:XIV.1.1}
\partial u / \partial t = A u, t > 0,
\end{equation}
where the differential operator 
\begin{equation}\label{eq:XIV.1.2}
A = a ^{ij} (x) \frac {\partial^{2}}{\partial x _{i} \partial x _{j}} + b ^{i} (x) \frac {\partial}{\partial x _{i}} + c (x) \left(a ^{ij} (x) = a ^{ji} (x)\right)
\end{equation}
is strictly elliptic in an $m$-dimensional euclidean space $\mathbb{R}^{m}$. We assume that the real-valued coefficients $a, b$ and $c$ are $C^{\infty}(\mathbb{R}^{m})$ functions and that 
\begin{equation}\label{eq:XIV.1.3}
\begin{gathered}
\max \left(\sup\limits _{x} | a ^{ij} (x) |, \sup\limits _{x} | b ^{i} (x) |, \sup\limits _{x} | c (x) |, \sup\limits _{x} | a _ {x _{k}} ^{ij} (x) |, \right. \left. \sup\limits _{x} | b _ {x _{k}} ^{i} (x) |, \sup\limits _{x} | a _ {x _{k} x _{8}} ^{ij} (x) |\right) = \eta < \infty .
\end{gathered}
\end{equation}

The strictly elliptic hypothesis concerning $A$ means that positive constants $\lambda_{0}$ and $\mu_{0}$ exist such that 
\begin{equation}\mu_{0} \sum\limits_{j=1} ^{m} \xi_{j} ^{2} \geqq a ^{ij} (x) \xi_{i} \xi_{j} \geqq \lambda_{0} \sum\limits_{i=1} ^{m} \xi_{i} ^{2} \text{ on } \mathbb{R}^{m} \text{ for any real } \text{ vector } \xi = (\xi_{1}, \xi_{2}, \dots , \xi_{m}) \label{eq:XIV.1.4}\end{equation}

Let $\hat{H}_0^1$ be the space of all real-valued $C_0^\infty (\mathbb{R}^{m})$ functions $f(x)$ normed by 
\begin{equation}\label{eq:XIV.1.5}
\left\| f \right\| _{1} = \left(\int_ {\mathbb{R}^{m}} f ^{2} \dd x + \sum\limits_{j=1} ^{m} \int_ {\mathbb{R}^{m}} f _ {x _{j}} ^{2} \dd x\right) ^{1/2},
\end{equation}
and let $H_0^1$ be the completion of $\hat{H}_0^1$ with respect to the norm $\|f\|_1$. Let similarly $H_0^0$ be the completion of $\hat{H}_0^1$ with respect to the norm 
\begin{equation}\label{eq:XIV.1.6}
\left\| f \right\| _{0} = \left(\int_ {\mathbb{R}^{m}} f ^{2} \dd x\right) ^{1/2}.
\end{equation}


We have thus introduced two real Hilbert spaces $H_0^1$ and $H_0^0$, and $H_0^1$ and $\hat{H}_0^1$ are $\|\|_0$-dense in $H_0^0$. We know, from the Proposition in Chapter~\ref{chap:I}, \ref{sec:I.10}, that $H_0^1$ coincides with the real Sobolev space $W^1(\mathbb{R}^{m})$; we know also that $H_0^0$ coincides with the real Hilbert space $L^2(\mathbb{R}^{m})$. We shall denote the scalar product in the Hilbert space $H_0^1$ (or in the Hilbert space $H_0^0$) by $(f,g)_1$ (or by $(f,g)_0$). 

To integrate equation \eqref{eq:XIV.1.1} in the complex Hilbert space $L^{2}(\mathbb{R}^{m})$ under condition \eqref{eq:XIV.1.3} and \eqref{eq:XIV.1.4}, we shall prepare some lemmas. These lemmas will also play important roles in the following sections. 

\begin{lemma}[concerning partial integration]\label{lem:XIV.1.Lemma.1} Let $f, g \in \tilde{H}_0^1$. We have then 
\begin{equation}\label{eq:XIV.1.7}
\begin{aligned}
&(A f, g) _{0} =  - \int_ {\mathbb{R}^{m}} a ^{ij} f _ {x _{i}} g _ {x _{j}} \dd x - \int_ {\mathbb{R}^{m}} a _ {x _{j}} ^{ij} f _ {x _{i}} g \dd x + \int_ {\mathbb{R}^{m}} b ^{i} f _ {x _{i}} g \dd x + \int_ {\mathbb{R}^{m}} c f g \dd x,
\end{aligned}
\end{equation}
that is, we may partially integrate, in $(Af, g)_{0}$, the terms containing the second order derivatives as if the integrated terms are zero. 
\end{lemma}

\begin{proof}
 By \eqref{eq:XIV.1.3} and the fact that $f$ and $g$ both belong to $H_0^1$, we see that $a^{ij}f_{x_ix_j}g$ is integrable over $\mathbb{R}^{m}$. Hence, by the fact that $f$ and $g$ are both of compact support, 
\[
\int_ {\mathbb{R}^{m}} a ^{ij} f _ {x _{i} x _{j}} g \dd x = - \int_ {\mathbb{R}^{m}} a ^{ij} f _ {x _{i}} g _ {x _{j}} \dd x - \int_ {\mathbb{R}^{m}} a _ {x _{j}} ^{ij} f _ {x _{i}} g \dd x. \]
\end{proof}

\begin{remarku} The formal adjoint $A^{*}$ of $A$ is defined by 
\begin{equation}\label{eq:XIV.1.8}
(A ^{*} f) (x) = \frac {\partial^{2}}{\partial x _{i} \partial x _{j}} \left(a ^{ij} (x) f (x)\right) - \frac {\partial}{\partial x _{i}} \left(b ^{i} (x) f (x)\right) + c (x) f (x).
\end{equation}
\end{remarku}

Then, as above, we have the result: If $f, g \in \tilde{H}_0^1$ then we may partially integrate, in $(A^*f, g)_0$, the terms containing the first and the second order derivatives as if the integrated terms are zero. That is, we have 
\begin{equation}\tag{7'}\label{eq:XIV.1.7'}
\begin{aligned}
&(A ^{*} f, g) _{0} =  - \int_ {\mathbb{R}^{m}} a ^{ij} f _ {x _{i}} g _ {x _{j}} \dd x - \int_ {\mathbb{R}^{m}} a _ {x _{i}} ^{ij} f g _ {x _{j}} \dd x - \int_ {\mathbb{R}^{m}} b ^{i} f g _ {x _{i}} \dd x + \int_ {\mathbb{R}^{m}} c f g \dd x.
\end{aligned}
\end{equation}

\begin{corollaryu} There exist positive constants $x$, $y$ and $\delta$ such that, for all sufficiently small positive constants $\alpha$, 
\begin{align}
\alpha \delta \| f \| _{1} ^{2}& \leqq (f - \alpha A f, f) _{0} \leqq (1 + \alpha \gamma) \| f \| _{1} ^{2} \quad \text{ when } \quad f \in \hat {H} _{0} ^{1}, \label{eq:XIV.1.9} \\
\alpha \delta \| f \| _{1} ^{2}& \leqq (f - \alpha A ^{*} f, f) _{0} \leqq (1 + \alpha \gamma) \| f \| _{1} ^{2} \quad \text{ when } \quad f \in \hat {H} _{0} ^{1}, \notag \\
\left| (f - \alpha A f, g) _{0} \right|& \leqq (1 + \alpha \gamma) \left\| f \right\| _{1} \left\| g \right\| _{1} \quad \text{ when } \quad f, g \in \hat {H} _{0} ^{1}, \label{eq:XIV.1.10} \\
\left| (f - \alpha A ^{*} f, g) _{0} \right|& \leqq (1 + \alpha \gamma) \left\| f \| _{1} \right. \| g \| _{1} \quad \text{ when } \quad f, g \in \hat {H} _{0} ^{1}, \notag \\
\left| (A f, g) _{0} - (f, A g) _{0} \right|& \leqq \varkappa \| f \| _{1} \| g \| _{0} \quad \text{ when } \quad f, g \in \hat {H} _{0} ^{1}. \label{eq:XIV.1.11}
\end{align}
\end{corollaryu}

\begin{proof}
 \eqref{eq:XIV.1.9} and \eqref{eq:XIV.1.10} may be proved by \eqref{eq:XIV.1.3}, \eqref{eq:XIV.1.4}, \eqref{eq:XIV.1.7} and \eqref{eq:XIV.1.7} remembering the inequality 
\begin{equation}\label{eq:XIV.1.12}
2 \alpha | a b | \leqq \alpha (\nu | a | ^{2} + \nu^{-1} | b | ^{2})
\end{equation}
which is valid whenever $\alpha$ and $\nu$ are $>0$. In fact, we can use an estimate such as 
\[
\left| \int_ {\mathbb{R}^{m}} a _ {x _{i}} ^{ij} f _ {x _{i}} g \dd x \right| \leqq \sum\limits_{i=1} ^{m} m \eta (v \| f _ {x _{i}} \| _{0} ^{2} + v ^{-1} \| g \| _{0} ^{2}). \]

We also obtain \eqref{eq:XIV.1.11} from 
\[
(A f, g) _{0} - (f, A g) _{0} = - \int_ {\mathbb{R}^{m}} \left(2 a _ {x _{i}} ^{ij} f _ {x _{j}} g + a _ {x _{i} x _{j}} ^{ij} f g - 2 b ^{i} f _ {x _{i}} g - b _ {x _{i}} ^{i} f g\right) \dd x.
\]
\end{proof}

\begin{lemma}[concerning the existence of solutions of $u - \alpha A u = f$]\label{lem:XIV.1.Lemma.2} Let a positive number $\alpha_0$ be so chosen that the above Corollary is valid for $0 < \alpha \leqq \alpha_0$. Then, for any function $f(x) \in \hat{H}_0^1$, the equation 
\begin{equation}\label{eq:XIV.1.13}
u - \alpha A u = f (0 < \alpha \leqq \alpha_{0})
\end{equation}
admits a uniquely determined solution $u \in H_{0}^{1} \cap C^{\infty}(\mathbb{R}^{m})$. 
\end{lemma}

\begin{proof}
 Let us define a bilinear functional $\hat{B}(u,v) = (u - \alpha A^{*}u,v)_{0}$ defined for functions $u, v \in \hat{H}_{0}^{1}$. From the above Corollary, we have 
\begin{equation}\label{eq:XIV.1.14}
\left| \hat {B} (u, v) \right| \leqq (1 + \alpha \gamma) \left\| u \right\| _{1} \left\| v \right\| _{1}, \alpha \delta_{r} \| u \| _{1} ^{2} \leqq \hat {B} (u, u). \end{equation}

Hence we may extend $\dot{B}(u,v)$, by continuity, to a bilinear functional $B(u,v)$ defined for $u,v\in H_{0}^{1}$ and satisfying 
\begin{equation}\tag{14'}\label{eq:XIV.1.14'}
\left| B (u, v) \right| \leqq (1 + \alpha \gamma) \left\| u \| _{1} \right\| v \| _{1}, \alpha \delta \| u \| _{1} ^{2} \leqq B (u, u).
\end{equation}

The linear functional $F(u) = (u, f)_0$ defined on $H_0^1$, is a bounded linear functional by $|(u, f)_0| \leqq \|u\|_0 \|f\|_0 \leqq \|u\|_1 \|f\|_0$. Hence, by F. \textsc{Riesz}' representation theorem in the Hilbert space $H_0^1$, there exists a uniquely determined $v = v(f) \in H_0^1$ such that $(u, f)_0 = (u, v(f))_1$. Thus, by the Milgram-Lax theorem applied to the Hilbert space $H_0^1$, we have 
\begin{equation}\label{eq:XIV.1.15}
(u, f) _{0} = (u, v (f)) _{1} = B (u, S v (f)) \text{ for all } u \in H _{0} ^{1}, \text{ where }
\end{equation}

$S$ is a bounded linear operator on $H_{0}^{1}$ onto $H_{0}^{1}$. 


Let $u$ run over $C_0^\infty(\mathbb{R}^{m})$, and let $v_n \in \hat{H}_0^1$ be such that $\lim\limits_{n \to \infty} \|v_n - Sv(f)\|_1 = 0$. Then 
\[
\begin{aligned}
&B (u, S v (f))= \lim\limits _{n \rightarrow \infty} B (u, v _{n}) = \lim\limits _{n \rightarrow \infty} \hat {B} (u, v _{n}) = \lim\limits _{n \rightarrow \infty} (u - \alpha A ^{*} u, v _{n}) _{0} = (u - \alpha A ^{*} u, S v (f)) _{0},
\end{aligned}
\]
because the norm $\|\|_1$ is larger than the norm $\|\|_0$. Hence 
\begin{equation}\tag{15'}\label{eq:XIV.1.15'}
(u, f) _{0} = (u - \alpha A ^{*} u, S v (f)) _{0},
\end{equation}
that is, $Sv(f) \in H_0^1$ is a distribution solution of the equation \eqref{eq:XIV.1.13}. Hence, by the strict ellipticity of $(I - \alpha A)$ and by the fact $f \in C_0^\infty(\mathbb{R}^{m})$, we see, from the Corollary of Friedrichs' theorem in Chapter~\ref{chap:VI}, \ref{sec:VI.9}, that we may consider $u = Sv(f) \in H_0^1$ to be a $C^\infty(\mathbb{R}^{m})$ solution of \eqref{eq:XIV.1.13}. Hence $u = Sv(f) \in H_0^1 \cap C^\infty(\mathbb{R}^{m})$. 


The uniqueness of such a solution $u$ of \eqref{eq:XIV.1.13} may be proved as follows. Let a function $u \in H_0^1 \cap C^\infty(\mathbb{R}^{m})$ satisfy $u - \alpha A u = 0$. Thus $Au \in H_0^1 \cap C^\infty(\mathbb{R}^{m}) \subseteq H_0^0$ and so the expression $(u - \alpha A u, u)_0$ is defined and $= 0$. Let $u_n \in \hat{H}_0^1$ be such that $\lim\limits_{n \to \infty} \|u - u_n\|_1 = 0$. We obtain, by partial integration as in \eqref{eq:XIV.1.9}, 
\[
0 = (u - \alpha A u, u) _{0} = \lim\limits _{n \rightarrow \infty} (u - \alpha A u, u _{n}) _{0} \geqq \alpha \delta \| u \| _{1} ^{2}, \text{ that is, } u = 0.
\]
\end{proof}

\begin{corollary}\label{cor:XIV.1.Corollary.1} Positive constants $\hat{\alpha}_{0}$ and $\eta_{0}$ exist such that, for any $f\in\hat{H}_{0}^{1}$, the equation 
\begin{equation}\label{eq:XIV.1.16}
\alpha u - A u = f \left(0 < \hat {\alpha} _{0} + \lambda_{0} + \eta_{0} \leqq \alpha\right)
\end{equation}
admits a uniquely determined solution $u = u_{f} \in H_{0}^{1} \cap C^{\infty}(\mathbb{R}^{m})$, and we have the estimate 
\begin{equation}\label{eq:XIV.1.17}
\left\| u _{f} \right\| _{0} \leqq (\alpha - \lambda_{0} - \eta_{0}) ^{-1} \left\| f \right\| _{0}.
\end{equation}
\end{corollary}

\begin{proof}
 By Schwarz' inequality we have 
\begin{equation}\label{eq:XIV.1.18}
\left\| (\alpha I - A) u \right\| _{0} \cdot \left\| u \right\| _{0} \geqq \left| \left((\alpha I - A) u, u\right) _{0} \right| \quad \text{ for } \quad u \in \hat {H} _{0} ^{1}. \end{equation}

By partial integration, we have 
\[
\begin{aligned}
&\left((\alpha I - A) u, u\right) _{0} = \alpha \| u \| _{0} ^{2} + \int_ {\mathbb{R}^{m}} a ^{ij} u _ {x _{i}} u _ {x _{j}} \dd x + \int_ {\mathbb{R}^{m}} a _ {x _{j}} ^{ij} u _ {x _{j}} u \dd x - \int_ {\mathbb{R}^{m}} b ^{i} u _ {x _{i}} u \dd x - \int_ {\mathbb{R}^{m}} c u u \dd x.
\end{aligned}
\]

Hence, we have, by \eqref{eq:XIV.1.3}, \eqref{eq:XIV.1.4} and \eqref{eq:XIV.1.12}, 
\[
\begin{aligned}
((\alpha I - A) u, u) _{0}& \geqq \alpha \| u \| _{0} ^{2} + \lambda_{0} (\| u \| _{1} ^{2} - \| u \| _{0} ^{2})
\\
&\quad \quad - m ^{2} \eta [ v (\| u \| _{1} ^{2} - \| u \| _{0} ^{2}) + v ^{-1} \| u \| _{0} ^{2} + m ^{-2} \| u \| _{0} ^{2} ]
\\
&= (\alpha - \lambda_{0} - m ^{2} \eta (v ^{-1} - v + m ^{-2})) \| u \| _{0} ^{2} + (\lambda_{0} - m ^{2} \eta v) \| u \| _{1} ^{2}.
\end{aligned}
\]

Thus by \eqref{eq:XIV.1.18} we have, for $\eta_0 = m^2\eta (\nu^{-1} - \nu +m^{-2})$ 
\begin{equation}\tag{17'}\label{eq:XIV.1.17'}
\|(\alpha I-A)u\|_{0}\geqq(\alpha-\lambda_{0}-\eta_{0})\|u\|_{0}\quad\text{whenever }u\in\hat{H}_{0}^{1}.
\end{equation}
by taking $\nu>0$ so small that $(\lambda_{0}-m^{2}\eta\nu)$ and $\eta_{0}$ are >$0$. We then take $\hat{\alpha}_{0}$ so large that for $\alpha\geqq\hat{\alpha}_{0}+\lambda_{0}+\eta_{0}$ we can apply Lemma~\ref{lem:XIV.1.Lemma.2} in solving \eqref{eq:XIV.1.16}. 

Since the solution $u = u_f \in H_0^1 \cap C^\infty(\mathbb{R}^{m})$ of \eqref{eq:XIV.1.16} is approximated by $\|\|_1$-norm by a sequence of functions $\in \hat{H}_0^1$, we obtain the estimate \eqref{eq:XIV.1.17} from \eqref{eq:XIV.1.18} and \eqref{eq:XIV.1.17'}. 
\end{proof}

\begin{corollary}\label{cor:XIV.1.Corollary.2} Consider $A$ as an operator on $D(A) = (\alpha I - A)^{-1} \hat{H}_0^1 \subseteq H_0^0$ into $H_0^0$. Then the smallest closed extension $\hat{A}$ in $H_0^0$ of $A$ admits, for $\alpha > \hat{\alpha}_0 + \lambda_0 + \eta_0$, the resolvent $(\alpha I - \hat{A})^{-1}$ defined on $H_0^0$ into $H_0^0$ such that 
\begin{equation}\label{eq:XIV.1.19}
\left\| (\alpha I - \hat {A}) ^{-1} \right\| _{0} \leqq (\alpha - \lambda_{0} - \eta_{0}) ^{-1}.
\end{equation}
\end{corollary}

\begin{proof}
 Clear from Corollary~\ref{cor:XIV.1.Corollary.1} remembering the fact that $D(A)$ and $\hat{H}_0^1$ both are $\|\|_0$-dense in $H_0^0$. 
\end{proof}

Thus, by Corollary~\ref{cor:IX.7.Corollary.2} in Chapter~\ref{chap:IX}, \ref{sec:IX.7}, $\hat{A}$ is the infinitesimal generator of a semi-group $T_{t}$ of class $(C_0)$ in the $B$-space $H_0^0$ such that $\| T_t \|_0 \leqq e^{(\lambda_0+ \eta_0)t}$ for $t \geqq 0$. 

Actually we can prove 

\begin{theorem}\label{thm:XIV.1.Theorem.1} Let Complex- $\hat{H}_0^1$ be the space of all complex-valued functions $f \in C_0^\infty(\mathbb{R}^{m})$ with $\|f\|_1 = \left( \int_{\mathbb{R}^{m}} |f|^2 \dd x + \sum\limits_{j=1}^{m} \int_{\mathbb{R}^{m}} |f_{x_{j}}|^2 \dd x \right)^{1/2} < \infty$. 
\end{theorem}

Let $\tilde{H}_0^1$ and $\tilde{H}_0^0$ be the completions of Complex- $\hat{H}_0^1$ with respect to the norm $\|f\|_1$ and $\|f\|_0$, respectively. We know that $\tilde{H}_0^1 = (\text{ complex })$ Sobolev space $W^1(\mathbb{R}^{m})$ (see Chapter~\ref{chap:I}, \ref{sec:I.10}). It is also clear that $\tilde{H}_0^0 = (\text{ complex })$ Hilbert space $L^2(\mathbb{R}^{m})$. We consider $A$ as an operator defined on $D(A) = (\alpha I - A)^{-1}\tilde{H}_0^1 \subseteq L^2(\mathbb{R}^{m})$ into $L^2(\mathbb{R}^{m})$. Then the smallest closed extension $\tilde{A}$ in $L^2(\mathbb{R}^{m})$ of $A$ is the infinitesimal generator of a holomorphic semi-group $T_t$ of class $(C_0)$ in $L^2(\mathbb{R}^{m})$ such that $\|T_t\|_0 \leqq e^{(\lambda_0+ \eta_0)t}$ for $t \geqq 0$. 

\begin{proof}
 By the preceding Corollary~\ref{cor:XIV.1.Corollary.2} and the reality of the coefficients in the differential operator $A$, we see that the range $R(\alpha I - A) = (\alpha I - A) \cdot D(A)$ is $\|\|_0$-dense in $L^2(\mathbb{R}^{m})$ for $\alpha > \hat{\alpha}_0 + \lambda_0 + \eta_0$. Moreover, we have, for $(u + iv) \in L^2(\mathbb{R}^{m})$ such that $(u + iv) \in D(A)$, 
\begin{equation}\left\| (\alpha I - A) (u + i v) \right\| _{0} ^{2} = \left\| (\alpha I - A) u \right\| _{0} ^{2} + \left\| (\alpha I - A) v \right\| _{0} ^{2} \geqq (\alpha - \lambda_{0} - \eta_{0}) ^{2} (\| u \| _{0} ^{2} + \| v \| _{0} ^{2}). \notag\end{equation}

Hence the inverse $(\alpha I - \tilde{A})^{-1}$ is a bounded linear operator on $L^2 (\mathbb{R}^{m})$ into $L^2 (\mathbb{R}^{m})$ such that $\|(\alpha I - \tilde{A})^{-1}\|_0 \leqq (\alpha - \lambda_0 - \eta_0)^{-1}$ for $\alpha > \hat{\alpha}_0 + \lambda_0 + \eta_0$. 

Therefore, by the Theorem in Chapter~\ref{chap:IX}, \ref{sec:IX.10}, we have only to show that 
\begin{equation}\label{eq:XIV.1.20}
\varlimsup\limits_{| \tau | \uparrow \infty} | \tau | \| ((\alpha + i \tau) I - \tilde {A}) ^{-1} \| _{0} < \infty .
\end{equation}

We have, for $w \in D(A)$, 
\begin{equation}\label{eq:XIV.1.21}
\left\| ((\alpha + i \tau) I - A) w \right\| _{0} \left\| w \right\| _{0} \geqq \left| \left(\left((\alpha + i \tau) I - A\right) w, w\right) _{0} \right|.
\end{equation}

By partial integration, we obtain, as in the proof of \eqref{eq:XIV.1.17}, 
\begin{gather}\quad \left| \text{ Real part } \left(\left((\alpha + i \tau) I - A\right) w, w\right) _{0} \right| = \bigl\| \alpha | w | _{0} ^{2} + \text{ Real part } \left(\int_ {\mathbb{R}^{m}} a ^{ij} w _ {x _{i}} \overline {{w}} _ {x _{j}} \dd x \right.
\\\quad \left. + a _ {x _{i}} ^{ij} w _ {x _{j}} \overline {{w}} \dd x - \int_ {\mathbb{R}^{m}} b ^{i} w _ {x _{i}} \overline {{w}} \dd x - \int_ {\mathbb{R}^{m}} c w \overline {{w}} \dd x\right) \bigr| \geqq (\alpha - \lambda_{0} - \eta_{0}) \| w \| _{0} ^{2} \notag
\\\quad + (\lambda_{0} - m ^{2} \eta \nu) \| w \| _{1} ^{2}. \notag
\end{gather}

Similarly we have 
\begin{gather}\left| \text{ Imaginary part } \left(\left((\alpha + i \tau) I - A\right) w, w\right) _{0} \right| \notag \\\geqq | \tau | \| w \| _{0} ^{2} - m ^{2} \eta (\| w \| _{1} ^{2} + m ^{-2} \| w \| _{0} ^{2}) | \geqq (| \tau | - \eta) \| w \| _{0} ^{2} - m ^{2} \eta \| w \| _{1} ^{2}. \notag
\end{gather}

Suppose there exists a $w_0 \in D(A)$, $\|w_0\|_0 \neq 0$, such that 
$\left|\operatorname{Im}\left(\left(\left(\alpha I+i\tau\right)I-A\right)w_{0},w_{0}\right)_{0}\right|<2^{-1}\left(\left|\tau\right|-\eta\right)\left\|w_{0}\right\|_{0}^{2}$ for sufficiently large $\tau$ (or for sufficiently large $-\tau$). Then, for such large $\tau$ (or $-\tau$), we must have 
\[
m ^{2} \eta \| w _{0} \| _{1} ^{2} > 2 ^{-1} (| \tau | - \eta) \| w _{0} \| _{0} ^{2}
\]
and so 
$\left|\left(\text{ Real part } \left(\left(\alpha+i\tau\right)I-A\right)w_{0},w_{0}\right)_{0}\right|\geqq\left(\lambda_{0}-m^{2}\eta\nu\right)\frac{\left(\left|\tau\right|-\eta\right)}{2m^{2}\eta}\left\|w_{0}\right\|_{0}^{2}.$ Hence we have proved Theorem~\ref{thm:XIV.1.Theorem.1} by \eqref{eq:XIV.1.21}. 
\end{proof}

\begin{theorem}\label{thm:XIV.1.Theorem.2} For any $f \in L^2(\mathbb{R}^{m})$, $u(t, x) = (T_t f)(x)$ is infinitely differentiable in $t > 0$ and in $x \in \mathbb{R}^{m}$, and $u(t, x)$ satisfies the diffusion equation \eqref{eq:XIV.1.1} as well as the initial condition $\lim\limits_{t \downarrow 0} \|u(t, x) - f(x)\|_0 = 0$. 
\end{theorem}

\begin{proof}
 We denote by $T_{t}^{(k)}$ the $k$-th strong derivative in $L^2(\mathbb{R}^{m})$ of $T_{t}$ with respect to $t$. $T_{t}$ being a holomorphic semi-group of class $(C_0)$ in $L^2(\mathbb{R}^{m})$, we have $T_{t}^{(k)}f = \tilde{A}^{k}T_{t}f\in L^{2}(\mathbb{R}^{m})$ if $t > 0$ ( $k = 0,1,\ldots$). Since $\tilde{A}$ is the smallest closed extension in $L^2(\mathbb{R}^{m})$ of $A$, we see that $A^{k}T_{t}f\in L^{2}(\mathbb{R}^{m})$ for fixed $t > 0$ ( $k = 0,1,2,\ldots$) if we apply the differential operator $A^{k}$ in the distributional sense. Thus, by the Corollary of Friedrichs' theorem in Chapter~\ref{chap:VI}, \ref{sec:VI.9}, $u(t,x)$ is, for fixed $t > 0$, equal to a $C^\infty(\mathbb{R}^{m})$ function after a correction on a set of measure zero. 


Because of the estimate $\left\|T_{t}\right\|_{0}\leqq e^{(\lambda_{0}+\eta_{0})t}$, we easily see that, if we apply, in the distributional sense, the elliptic differential operator 
\[
\left(\frac {\partial^{2}}{\partial t ^{2}} + A\right)
\]
any number of times to $u(t, x)$, then the result is locally square integrable in the product space $\{t; 0 < t < \infty\} \times \mathbb{R}^{m}$. Thus again, by the Corollary of Friedrichs' theorem in Chapter~\ref{chap:VI}, \ref{sec:VI.9}, we see that $u(t, x)$ is equal to a function which is $C^\infty$ in $(t, x)$, $t > 0$ and $x \in \mathbb{R}^{m}$, after a correction on a set of measure zero of the product space. Thus we may consider that $u(t, x)$ is a genuine solution of \eqref{eq:XIV.1.1} satisfying the initial condition $\lim\limits_{t \downarrow 0} \| u(t, x) - f(x) \|_0 = 0$.
\end{proof}

\begin{remarku} The above obtained solution $u(t, x)$ satisfies the ``forward and backward unique continuation property'': 
\end{remarku}

if, for a fixed $t_0 > 0$, $u(t_0, x) \equiv 0$ on an open set $G \subseteq \mathbb{R}^{m}$; 


then $u(t, x) = 0$ for every $t > 0$ and every $x \in G$. 

\begin{proof}
 Since $T_{t}$ is a holomorphic semi-group of class $(C_0)$, we have 
\[
\lim\limits _{n \rightarrow \infty} \left\| T _ {t _{0} + h} f - \sum\limits_{k=0} ^{n} (k!) ^{-1} h ^{k} A ^{k} T _ {t _{0}} f \right\| _{0} = 0
\]
for a sufficiently small $h$. Hence, as in the proof of the completeness of the space $L^2(\mathbb{R}^{m})$, there exists a sequence $\{n'\}$ of natural numbers such that 
\[
u \left(t _{0} + h, x\right) = \lim\limits _ {n ^{\prime} \rightarrow \infty} \sum\limits_{k=0} ^ {n ^{\prime}} (k!) ^{-1} h ^{k} A ^{k} u \left(t _{0}, x\right) \text{ for a.e. } x \in \mathbb{R}^{m}. \]

By hypothesis given in ($22$), we have $A^{k}u(t_{0}, x) = 0$ in $G$ and so we must have $u(t_{0} + h, x) = 0$ in $G$ for sufficiently small $h$. Repeating the process we see that conclusion ($22$) is true. 
\end{proof}

\yossub{References}

The result of this section is adapted from K. \textsc{Yosida} \cite{yosida364}. As for the ``forward and backward unique continuation property'' ($22$), we have the more precise result that $u(t, x) = 0$ for every $t > 0$ and every $x \in \mathbb{R}^{m}$. For, S. \textsc{Mizohata} \cite{mizohata210} has proved a ``space-like unique continuation property'' of solutions $u(t, x)$ of a diffusion equation to the effect that $u(t, x) = 0$ for all $t > 0$ and all $x \in \mathbb{R}^{m}$ if $u(t, x) = 0$ for all $t > 0$ and all $x \in G$. Concerning the holomorphic character of the semi-group $T_{t}$ obtained above, see also R. S. \textsc{Phillips} \cite{phillips256}. It is to be noted here that the unique continuation property of the solutions of a heat equation $\partial u/\partial t = \Delta u$ was first proposed and solved by H. \textsc{Yamabe}-S. ITô \cite{yamabe342}. 

There is a fairly complete discussion of parabolic equations from the view point of the theory of dissipative operators. See R. S. \textsc{Phillips} \cite{phillips257}. 


\yossection{2}{Integration of Diffusion Equations in a Compact Riemannian Space}
\label{sec:XIV.2}

Let $R$ be a connected, orientable $m$-dimensional $C^\infty$ Riemannian space with the metric 
\begin{equation}\label{eq:XIV.2.1}
\dd s ^{2} = g _{ij} (x) \dd x ^{i} \dd x ^{j}.
\end{equation}

Let $A$ be a second order linear partial differential operator in $R$ with real-valued $C^\infty$ coefficients: 
\begin{equation}\label{eq:XIV.2.2}
A = a ^{ij} (x) \frac {\partial^{2}}{\partial x ^{i} \partial x ^{j}} + b ^{i} (x) \frac {\partial}{\partial x ^{i}}.
\end{equation}

We assume that $a^{ij}$ is a symmetric contravariant tensor and that $b^{i}(x)$ satisfies the transformation rule 
\begin{equation}\label{eq:XIV.2.3}
\bar {b} ^{i} = b ^{k} \frac {\partial \bar {x} ^{i}}{\partial x ^{k}} + a ^{kj} \frac {\partial^{2} \bar {x} ^{i}}{\partial x ^{k} \partial x ^{j}}
\end{equation}
by the coordinate change $(x^{1}, x^{2}, \ldots, x^{m}) \to (\bar{x}^{1}, \bar{x}^{2}, \ldots, \bar{x}^{m})$ so that the value $(A f)(x)$ is determined independently of the choice of the local coordinates. We further assume that $A$ is strictly elliptic in the sense that there exist positive constants $\lambda_{0}$ and $\mu_{0}$ such that 
\begin{equation}\mu_{0} \sum\limits_{j=1} ^{m} \xi_{j} ^{2} \geqq a ^{ij} (x) \xi_{i} \xi_{j} \geqq \lambda_{0} \sum\limits_{j=1} ^{m} \xi_{j} ^{2} \text{ for every real vector } \left(\xi_{1}, \dots , \xi_{m}\right) \text{ and every } x \in \mathbb{R}^{m}. \label{eq:XIV.2.4}\end{equation}

We consider the Cauchy problem in the large on $R$ for the diffusion equation: to find solution $u(t, x)$ such that 
\begin{equation}\label{eq:XIV.2.5}
\begin{aligned}
&\partial u / \partial t  = A u, t > 0, u (0, x) = f (x) \text{ where } f (x) \text{ is a given } \text{ function on } R.
\end{aligned}
\end{equation}

We shall prove 

\begin{theoremu} If $R$ is compact so that $R$ is without boundary, then the equation \eqref{eq:XIV.2.5} admits, for any initial function $f \in C^{\infty}(R)$, a uniquely determined solution $u(t, x)$ which is $C^{\infty}$ in $(t, x)$, $t > 0$ and $x \in R$. This solution can be represented in the form 
\begin{equation}\label{eq:XIV.2.6}
u (t, x) = \int_{R} P (t, x, \dd y) f (y)
\end{equation}
where $P(t, x, E)$ is the transition probability of a Markov process on $R$. \end{theoremu}

\begin{proof}
Let us denote by $C(R)$ the $B$-space of real-valued continuous functions $f(x)$ on $R$ normed by $\|f\| = \sup |f(x)|$. We first prove 
\begin{align}
\text{ for any } f& \in C ^{\infty} (R) \text{ and any } n > 0, \text{ we have } \notag \\
\max\limits _{x} h (x)& \geqq f (x) \geqq \min\limits _{x} h (x) \text{ where } h (x) = f (x) - n ^{-1} (A f) (x). \label{eq:XIV.2.7}
\end{align}


Suppose that $f(x)$ attains its maximum at $x = x_0$. We choose a local coordinate system at $x_0$ such that $a^{ij}(x_0) = \delta_{ij} (= 1 \text{ or } 0 \text{ according as } i = j \text{ or } i \neq j)$. Such a choice is possible owing to condition \eqref{eq:XIV.2.4}. Then 
\[
\begin{aligned}
&h (x _{0})  = f (x _{0}) - n ^{-1} (A f) (x _{0}) = f (x _{0}) - n ^{-1} b ^{i} (x _{0}) \frac {\partial f}{\partial x _{0} ^{i}} - n ^{-1} \sum\limits_{i=1} ^{m} \frac {\partial^{2} f}{\partial (x _{0} ^{i}) ^{2}} \geqq f (x _{0}),
\end{aligned}
\]
since we have, at the maximum point $x_0$, 
\[
\frac {\partial f}{\partial x _{0} ^{i}} = 0 \quad \text{ and } \quad \frac {\partial^{2} f}{\partial (x _{0} ^{i}) ^{2}} \leqq 0.
\]

Thus $\max h(x) \geqq f(x)$. Similarly we have $f(x) \geqq \min h(x)$. 

We shall consider $A$ as an operator on $D(A) = C^{\infty}(R) \subseteq C(R)$ into $C(R)$. Then, by \eqref{eq:XIV.2.7}, we see that the inverse $(I - n^{-1}A)^{-1}$ exists for $n > 0$ and $\|(I - n^{-1}A)^{-1}g\| \leqq \|g\|$ for $g$ in the range $R(I - n^{-1}A) =$ $(I - n^{-1}A) \cdot D(A)$. This range is, for sufficiently large $n$, strongly dense in $C(R)$. For, as in the preceding section, we have the result: For any $g \in C^{\infty}(R)$ and for sufficiently large $n > 0$, the equation $u - n^{-1}Au = g$ admits a uniquely determined solution $u \in C^{\infty}(R)$. Because, by the compactness of the Riemann space $R$, Lemma~\hyperref[lem:XIV.1.Lemma.1]{$1$} concerning the partial integration in the preceding section may be adapted to our Riemannian space $R$ without boundary. Moreover, $C^{\infty}(R)$ is strongly dense in $C(R)$ as may be seen by the regularization of functions in $C(R)$ (see Proposition~\ref{prop:I.1.Proposition.8} in Chapter~\ref{chap:I}, \ref{sec:I.1}). 

Hence the smallest closed extension $\overline{A}$ in $C(R)$ of the operator $A$ satisfies the conditions: 

for sufficiently large $n > 0$, the resolvent $(I - n^{-1}\overline{A})^{-1}$ 
exists as a bounded linear operator on $C(R)$ into $C(R)$ 
\begin{align}
\text{ such that } \left\| (I - n ^{-1} \overline {{A}}) ^{-1} \right\|& \leqq 1, \label{eq:XIV.2.8} \\
((I - n ^{-1} \bar {A}) ^{-1} h) (x)& \geqq 0 \text{ on } R \text{ whenever } h (x) \geqq 0 \text{ on } R, \label{eq:XIV.2.9} \\
(I - n ^{-1} \overline {{{A}}}) ^{-1} \cdot 1& = 1. \label{eq:XIV.2.10}
\end{align}

The positivity of the operator $(I - n^{-1}\overline{A})^{-1}$ given in \eqref{eq:XIV.2.9} is clear from \eqref{eq:XIV.2.7}. The equation \eqref{eq:XIV.2.10} follows from $A\cdot 1 = 0$. 

Therefore $\overline{A}$ is the infinitesimal generator of a contraction semigroup $T_{t}$ in $C(R)$ of class $(C_0)$. As in the preceding section, we see, by the strict ellipticity of $A$, that, for any $f \in C^{\infty}(R)$, the function $u(t,x) = (T_{t}f)(x)$ is a $C^{\infty}$ function in $(t,x)$ for $t > 0$ and $x \in R$ so that $u(t,x)$ is a genuine solution of \eqref{eq:XIV.2.5}. 

Since the dual space of the space $C(R)$ is the space of Baire measures in $R$, we easily prove the latter part of the Theorem remembering \eqref{eq:XIV.2.9} and \eqref{eq:XIV.2.10}. 
\end{proof}

\yossection{3}{Integration of Wave Equations in a Euclidean Space \texorpdfstring{$\mathbb{R}^{m}$}{Rm}}
\label{sec:XIV.3}

Consider a wave equation 
\begin{equation}\label{eq:XIV.3.1}
\partial^{2} u / \partial t ^{2} = A u, - \infty < t < \infty ,
\end{equation}
where the differential operator 
\begin{equation}\label{eq:XIV.3.2}
A = a ^{ij} (x) \frac {\partial^{2}}{\partial x _{i} \partial x _{j}} + b ^{i} (x) \frac {\partial}{\partial x ^{i}} + c (x) \quad \left(a ^{ij} (x) = a ^{ji} (x)\right)
\end{equation}
is strictly elliptic in an $m$-dimensional Euclidean space $\mathbb{R}^{m}$. We assume that the real-valued $C^{\infty}$ coefficients $a, b$, and $c$ satisfy conditions \eqref{eq:XIV.3.3} and \eqref{eq:XIV.1.4} in Chapter~\ref{chap:XIV}, \ref{sec:XIV.1}. As done there, we shall denote by $\hat{H}_{0}^{1}$ the space of all real-valued $C_{0}^{\infty}(\mathbb{R}^{m})$ functions $f(x)$ normed by 
\[
\left\| f \right\| _{1} = \left(\int_ {\mathbb{R}^{m}} f ^{2} \dd x + \sum\limits_{j=1} ^{m} \int_ {\mathbb{R}^{m}} f _ {x _{j}} ^{2} \dd x\right) ^{1/2}
\]
and let $H_0^1$ (and $H_0^0$) be the completion of $\hat{H}_0^1$ with respect to the norm $\|f\|_1$ (and with respect to the norm $\|f\|_0 = \left( \int_{\mathbb{R}^{m}} f^2 \dd x \right)^{1/2}$). 

\begin{lemmau} For any pair $\{f, g\}$ of elements $\in \hat{H}_0^{1'}$ the equation 
\begin{equation}\label{eq:XIV.3.3}
\left(\begin{pmatrix}I & 0 \\ 0 & I\end{pmatrix} - n ^{-1} \begin{pmatrix}0 & I \\ A & 0\end{pmatrix}\right) \begin{pmatrix}u \\ v\end{pmatrix} = \begin{pmatrix}f \\ g\end{pmatrix}
\end{equation}
admits, if the integer $n$ be such that $|n^{-1}|$ is sufficiently small, a uniquely determined solution $\{u, v\}$, $u$ and $v \in H_0^1 \cap C^\infty(\mathbb{R}^{m})$, satisfying the estimate 
\begin{equation}\label{eq:XIV.3.4}
\begin{gathered}
\left((u - \alpha_{0} A u, u) _{0} + \alpha_{0} (v, v) _{0}\right) ^{1/2} \leqq (1 - \beta | n | ^{-1}) ^{-1} \left((f - \alpha_{0} A f, f) _{0} + \alpha_{0} (g, g) _{0}\right) ^{1/2},
\end{gathered}
\end{equation}
with positive constants $\alpha_0$ and $\beta$ which are independent of $n$ and $\{f, g\}$. \end{lemmau}

\begin{proof}
Let $u_1 \in H_0^1 \cap C^\infty(\mathbb{R}^{m})$ and $v_1 \in H_0^1 \cap C^\infty(\mathbb{R}^{m})$ be respectively the solutions of 
\begin{equation}\label{eq:XIV.3.5}
u _{1} - n ^{-2} A u _{1} = f \quad \text{ and } \quad v _{1} - n ^{-2} A v _{1} = g.
\end{equation}


The existence of such solutions for sufficiently small $|n^{-1}|$ was proved in Lemma~\ref{lem:XIV.1.Lemma.2} in Chapter~\ref{chap:XIV}, \ref{sec:XIV.1}. Then 
\begin{equation}\label{eq:XIV.3.6}
u = u _{1} + n ^{-1} v _{1}, v = n ^{-1} A u _{1} + v _{1}
\end{equation}
satisfies \eqref{eq:XIV.3.3}, i.e., we have $u - n^{-1}v = f$, $v - n^{-1}Au = g$. 

We next prove \eqref{eq:XIV.3.4}. We remark that 
$A u = n (v - g) \in H_0^1 \cap C^\infty(\mathbb{R}^{m}) \subseteq H_0^0$ and hence, by $f, g \in C_0^\infty(\mathbb{R}^{m})$, 
\[
A v = n (A u - A f) \in H _{0} ^{1} \cap C ^{\infty} (\mathbb{R}^{m}) \subseteq H _{0} ^{0}.
\]

We have thus, by \eqref{eq:XIV.3.3}, 
and 
\begin{gather}(f - \alpha_{0} A f, f) _{0}  = (u - n ^{-1} v - \alpha_{0} A (u - n ^{-1} v), u - n ^{-1} v) _{0} \notag
\\= (u - \alpha_{0} A u, u) _{0} - 2 n ^{-1} (u, v) _{0} + \alpha_{0} n ^{-1} (A u, v) _{0} + \alpha_{0} n ^{-1} (A v, u) _{0} + n ^{-2} (v - \alpha_{0} A v, v) _{0} \notag
\\\alpha_{0} (g, g) _{0}  = \alpha_{0} (v - n ^{-1} A u, v - n ^{-1} A u) _{0} \notag
\\= \alpha_{0} (v, v) _{0} - \alpha_{0} n ^{-1} (v, A u) _{0} - \alpha_{0} n ^{-1} (A u, v) _{0} + \alpha_{0} n ^{-2} (A u, A u) _{0}. \notag
\end{gather}

By a limiting process, we can prove that \eqref{eq:XIV.3.9}, \eqref{eq:XIV.3.10} and \eqref{eq:XIV.1.11} in Chapter~\ref{chap:XIV}, \ref{sec:XIV.1} are valid for $f = u$ and $g = v$. Thus, by \eqref{eq:XIV.1.12} in Chapter~\ref{chap:XIV}, \ref{sec:XIV.1}, there exists a positive constant $\beta$ satisfying 
\[
\begin{aligned} ((f - \alpha_{0} A f, f) _{0} + \alpha_{0} (g, g) _{0}) ^{1/2} & \geqq ((u - \alpha_{0} A u, u) _{0} + \alpha_{0} (v, v) _{0} \\ & - \alpha_{0} | n ^{-1} | | (A u, v) _{0} - (A v, u) _{0} | - 2 | n ^{-1} | | (u, v) _{0} |) ^{1/2} \\ & \geqq (1 - \beta | n ^{-1} |) ((u - \alpha_{0} A u, u) _{0} + \alpha_{0} (v, v) _{0}) ^{1/2} \end{aligned}
\]
for sufficiently large $|\pmb{n}|$. 

The above estimate for the solutions $\{u, v\}$ belonging to $H_0^1 \cap C^\infty(\mathbb{R}^{m})$ shows that such solutions are uniquely determined by $\{f, g\}$. 
\end{proof}

\begin{corollaryu} The product space $H_{0}^{1} \cap H_{0}^{0}$ of vectors 
\begin{equation}\label{eq:XIV.3.7}
\begin{pmatrix}u \\ v\end{pmatrix} = \{u, v \} ^{\prime}, \text{ where } u \in H _{0} ^{1} \text{ and } v \in H _{0} ^{0},
\end{equation}
is a $B$-space by the norm 
\begin{equation}\label{eq:XIV.3.8}
\left\| \begin{pmatrix}u \\ v\end{pmatrix} \right\| = \| \{u, v \} ^{\prime} \| = (B (u, u) + \alpha_{0} (v, v) _{0}) ^{1/2},
\end{equation}
where $B(f, g)$ is the extension by continuity with respect to the norm $\|\|_1$ of the bilinear functional 
\[
\hat {B} (f, g) = (f - \alpha_{0} A f, g) _{0} \text{ defined for } f, g \in \hat {H} _{0} ^{1}.
\]
\end{corollaryu}

We know that $B(u, u)^{1/2}$ is equivalent to the norm $\|u\|_1$ (see Chapter~\ref{chap:XIV}, \ref{sec:XIV.1}): 
\begin{equation}\label{eq:XIV.3.9}
\alpha_{0} \delta \| u \| _{1} ^{2} \leqq B (u, u) \leqq (1 + \alpha_{0} \gamma) \| u \| _{1} ^{2}.
\end{equation}

Let the domain $D(\mathfrak{A})$ of the operator 
\begin{equation}\label{eq:XIV.3.10}
\mathfrak {A} = \begin{pmatrix}0 & I \\ A & 0\end{pmatrix}
\end{equation}
be the vectors $\{u, v\}' \in H_0^1 \times H_0^0$ such that $u, v \in H_0^0$ are given by \eqref{eq:XIV.3.6}. Then the Lemma shows that the range of the operator 
\[
\operatorname{Im} - n ^{-1} \mathfrak {A}, \text{ where } \operatorname{Im} = \begin{pmatrix}I & 0 \\ 0 & I\end{pmatrix},
\]
contains all the vectors $\{f, g\}'$ such that $f, g \in H_0^1$. Thus the smallest closed extension $\overline{\mathfrak{A}}$ in $H_0^1 \times H_0^0$ of the operator $\mathfrak{A}$ is such that the operator $(\mathfrak{I} - n^{-1} \overline{\mathfrak{A}})$ with integral parameter $n$ admits, for sufficiently large $|n|$, an everywhere in $H_0^1 \times H_0^0$ defined inverse $(\mathfrak{I} - n^{-1} \overline{\mathfrak{A}})^{-1}$ satisfying 
\begin{equation}\label{eq:XIV.3.11}
\left\| (\operatorname{Im} - n ^{-1} \overline {{{\mathfrak {A}}}}) ^{-1} \right\| \leqq (1 - \beta | n ^{-1} |) ^{-1}.
\end{equation}

We are now prepared to prove 

\begin{theoremu} For any pair $\{f(x), g(x)\}$ of $C_0^\infty(\mathbb{R}^{m})$ functions, the equation \eqref{eq:XIV.3.1} admits a $C^\infty$ solution $u(t, x)$ satisfying the initial condition 
\begin{equation}\label{eq:XIV.3.12}
u (0, x) = f (x), u _{t} (0, x) = g (x)
\end{equation}
and the estimate 
\begin{equation}\label{eq:XIV.3.13}
(B (u, u) + \alpha_{0} (u _{t}, u _{t}) _{0}) ^{1/2} \leqq \exp (\beta | t |) (B (f, f) + \alpha_{0} (g, g) _{0}) ^{1/2}.
\end{equation}
\end{theoremu}

\begin{remarku} Formula \eqref{eq:XIV.3.9} shows that $B(u, u)$ is comparable to the potential energy of the wave (= the solution of \eqref{eq:XIV.3.1}) $u(t, x)$, and $(u_t, u_t)_0$ is comparable to the kinetic energy of the wave $u(t, x)$. Thus \eqref{eq:XIV.3.13} means that the total energy of the wave $u(t, x)$ does not increase more rapidly than $\exp(\beta |t|)$ when the time $t$ tends to $\pm\infty$. This is a kind of energy inequality which governs wave equations in general. 
\end{remarku}

\begin{proof}[Proof of Theorem.]
Estimate \eqref{eq:XIV.3.11} shows that $\overline{\mathfrak{A}}$ is the infinitesimal generator of a group $T_{t}$ of class $(C_0)$ in $H_0^1 \times H_0^0$ such that
\begin{equation}\label{eq:XIV.3.14}
\left\| T _{t} \right\| \leqq \exp (\beta | t |), - \infty < t < \infty .
\end{equation}

By hypothesis we have, for $k = 0, 1, 2, \ldots$, 
\[
\overline {{\mathfrak {A}}} ^{k} \begin{pmatrix}f \\ g\end{pmatrix} = \mathfrak {A} ^{k} \begin{pmatrix}f \\ g\end{pmatrix} \in C _{0} ^{\infty} (\mathbb{R}^{m}) \times C _{0} ^{\infty} (\mathbb{R}^{m}) \subseteq H _{0} ^{1} \times H _{0} ^{0}.
\]

Hence, if we put 
\[
\begin{pmatrix}u (t, x) \\ v (t, x)\end{pmatrix} = T _{t} \begin{pmatrix}f (x) \\ g (x)\end{pmatrix},
\]
then, by the commutativity of $\overline{A}$ with $T_{t}$, we have 
\[
\frac {\partial^{k}}{\partial t ^{k}} \begin{pmatrix}u (t, x) \\ v (t, x)\end{pmatrix} = \frac {\partial^{k} T _{t}}{\partial t ^{k}} \begin{pmatrix}f (x) \\ g (x)\end{pmatrix} = \mathfrak {A} ^{k} \begin{pmatrix}u (t, x) \\ v (t, x)\end{pmatrix} \in H _{0} ^{1} \times H _{0} ^{0}
\]
for $k = 0, 1, 2, \ldots$. Here we denote by $\partial^k T_t / \partial t^k$ the $k$ th strong derivative in $H_0^1 \times H_0^0$. Therefore, by $H_0^1 \subseteq H_0^0 = L^2(\mathbb{R}^{m})$ and the strict ellipticity of the operator $A$, we see, as in the proof of Theorem~\ref{thm:XIV.1.Theorem.2} in Chapter~\ref{chap:XIV}, \ref{sec:XIV.1}, that $u(t, x)$ is $C^\infty$ in $(t, x)$ for $-\infty < t < \infty$, $x \in \mathbb{R}^{m}$, and satisfies equation \eqref{eq:XIV.3.1} with \eqref{eq:XIV.3.12} and estimate \eqref{eq:XIV.3.13}. 
\end{proof}

\begin{remarku} The result of the present section is adapted from K. \textsc{Yosida} \cite{yosida365}. Cf. J. L. \textsc{Lions} \cite{liggett178}. P. D. \textsc{Lax} has kindly communicated to the present author that the method of integration given in this section is very similar to that announced by him in Abstract $180$, Bull. Amer. Math. Soc. $58$, $192$ ($1952$). It is to be noted here that our method can be modified to the integration of wave equations in an open domain of a Riemannian space. There is another approach to the integration of wave equations based upon the theory of dissipative semi-groups. See R. S. \textsc{Phillips} \cite{phillips258} and \cite{phillips259}. The method is closely connected with the theory of symmetric positive systems by K. \textsc{Friedrichs} \cite{friedrichs70}. Cf. also P. \textsc{Lax}-R. S. \textsc{Phillips} \cite{lewy173}. 
\end{remarku}


\yossection{4}{Integration of Temporally Inhomogeneous Equations of Evolution in a \texorpdfstring{$B$}{B}-space}
\label{sec:XIV.4}

We shall be concerned with the integration of the equation 
\begin{equation}\label{eq:XIV.4.1}
\dd x (t) / \dd t = A (t) x (t), a \leqq t \leqq b.
\end{equation}

Here the unknown $x(t)$ is an element of a $B$-space $X$, depending on a real parameter $t$, while $A(t)$ is a given, in general unbounded, linear operator with domain $D(A(t))$ and range $R(A(t))$, both in $X$, depending also on $t$. 

T. \textsc{Kato} \cite{kato136}, \cite{kato137} was the first to make a successful attack on the problem of integration of \eqref{eq:XIV.4.1}. He assumed the following four conditions: 

\begin{enumerate}[(i)]
\item The domain $D(A(t))$ is independent of $t$ and is strongly dense in $X$ such that, for $\alpha > 0$, the resolvent $(I - \alpha A(t))^{-1}$ exists as a bounded linear operator $\in L(X, X)$ with the norm $\leqq 1$.
\item The operator $B(t, s) = (I - A(t))(I - A(s))^{-1}$ is uniformly bounded in norm for $t \geqq s$.
\item $B(t, s)$ is, at least for some $s$, of bounded variation in $t$ in norm, that is, for every partition $s = t_{0} < t_{1} < \cdots < t_{n} = t$ of $[a, b]$, 
\[
\sum\limits_{j=0} ^{n-1} \left\| B (t _{j+1}, s) - B (t _{j}, s) \right\| \leqq N (a, b) < \infty .
\]
\item $B(t, s)$ is, at least for some $s$, weakly differentiable in $t$ and $\partial B(t, s)/\partial t$ is strongly continuous in $t$.
\end{enumerate}
Under these conditions, \textsc{Kato} proved that the limit 
\[
U (t, s) x _{0} = \underset {\max \left| t _{j+1} - t _{j} \right|\rightarrow 0} {\slim} \prod\limits_{j=n-1} ^{0} \exp \left(\left(t _{j+1} - t _{j}\right) A (t _{j})\right) x _{0}
\]
exists for every $x_0 \in X$ and gives the unique solution of \eqref{eq:XIV.4.1} with the initial condition $x(s) = x_0$ at least when $x_0 \in D(A(s))$. 

Kato's method is thus an abstraction of the classical polygon method of Cauchy for ordinary differential equation $\dd x(t) / \dd t = a(t)x(t)$. Although his method is very simple and natural in its idea, the proof is somewhat lengthy because it is concerned with a general partition of the interval $[s,t]$. \textsc{Kato} \cite{kato136} has shown that the proof can be simplified when the space $X$ is reflexive. For the reflexive $B$-space case, see also K. \textsc{Yosida} \cite{yosida371}. 

In this section, we shall be concerned with an equi-partition of a fixed interval, say, [$0$, $1$], independently of $s$ and $t$, to the effect that \textsc{Kato}'s original method be modified so as to yield a fairly simplified presentation. 

We shall assume the following four conditions which are essentially the same as \textsc{Kato}'s conditions (i) through (iv) above: 
\begin{equation}
\label{eq:XIV.4.2}
\begin{minipage}{0.92\linewidth}
\everymath{}\let\limits\nolimits
$D (A (t))$ is independent of $t$ and it is dense in $X.$
\end{minipage}
\end{equation}

\begin{equation}\label{eq:XIV.4.3}
\parbox{0.92\linewidth}{\everymath{}\let\limits\nolimits For every $\lambda \geqq 0$ and $t, 0 \leqq t \leqq 1$, the resolvent $(\lambda I - A(t))^{-1}$ exists as an operator $\in L(X, X)$ such that $\| (\lambda I - A(t))^{-1} \| \leqq \lambda^{-1}$ for $\lambda > 0$.}
\end{equation}

\begin{equation}\label{eq:XIV.4.4}
A (t) A (s) ^{-1} \in L (X, X) \quad \text{ for } \quad 0 \leqq s, t \leqq 1.
\end{equation}

\begin{equation}\label{eq:XIV.4.5}
\parbox{0.92\linewidth}{\everymath{}\let\limits\nolimits For every $x \in X$, $(t - s)^{-1} C(t, s) x = (t - s)^{-1}(A(t) A(s)^{-1} - I)x$ is bounded and uniformly strongly continuous in $t$ and $s$, $t \neq s$, and \slim $k \cdot C\left(t, t - \frac{1}{k}\right)x = C(t)x$ exists uniformly in $t$ so that $C(t) \in L(X, X)$ is strongly continuous in $t$.}
\end{equation}


\begin{remarku} Condition \eqref{eq:XIV.4.4} is stated for the sake of convenience. It is implied by (2) and the closed graph theorem, since $A(s)$ is a closed linear operator by ($3$). Moreover, conditions (2) and ($3$) imply that $A(s)$ is the infinitesimal generator of a contraction semi-group $\{\exp(tA(s)); t \geqq 0\}$ of class $(C_0)$. Hence we have (see Chapter~\ref{chap:IX}): 
\begin{align}
    \exp (t A (s)) x = & \slim _{n \rightarrow \infty} \exp \left(t A (s) (I - n ^{-1} A (s)) ^{-1}\right) x \text{ uniformly in } t, 0 \leqq t \leqq 1; \label{eq:XIV.4.6} \\
    \exp (t _{1} A (s)) \cdot \exp (t _{2} A (s)) = & \exp ((t _{1} + t _{2}) A (s)); \label{eq:XIV.4.7} \\
    \frac {\dd \exp (t A (s)) y}{\dd t} = & A (s) \exp (t A (s)) y = \exp (t A (s)) \cdot A (s) y, y\in D (A (s)), \notag \\  \text{where } \dd / \dd t & \text{ means differentiation in the strong topology of } X; \label{eq:XIV.4.8} \\
    \slim _ {t \to t _{0}} \exp (t A (s)) x = & \exp (t _{0} A (s)) x. \label{eq:XIV.4.9}
\end{align}
\end{remarku}

We are now able to state our result. 

\begin{theorem}\label{thm:XIV.4.Theorem.1} For any positive integer $k$ and $0 \leqq s \leqq t \leqq 1$, define the operator $U_k(t, s) \in L(X, X)$ through 
\begin{equation}\label{eq:XIV.4.1a}
\left.\begin{aligned}
    U _{k} (t, s)& = \exp \left((t - s) A \left(\frac {i - 1}{k}\right)\right) \quad \text{ for } \quad \frac {i - 1}{k} \leqq s \leqq t \leqq \frac {i}{k}, (1 \leqq i \leqq k), \\
    U _{k} (t, r)& = U _{k} (t, s) U _{k} (s, r) \quad \text{ for } \quad 0 \leqq r \leqq s \leqq t \leqq 1. 
\end{aligned}\right\}
\end{equation}

Then, for every $x \in X$ and $0 \leqq s \leqq t \leqq 1$, 
\begin{equation}\label{eq:XIV.4.11}
\slim _{k \to \infty} U _{k} (t, s) x = U (t, s) x \text{ exists uniformly in } t \text{ and } s.
\end{equation}

Moreover, if $y \in D(A(0))$ then the Cauchy problem 
\begin{equation}\tag{1'}\label{eq:XIV.4.1'}
\frac {\dd x (t)}{\dd t} = A (t) x (t), x (0) = y \text{ and } x (t) \in D (A (t)), 0 \leqq t \leqq 1
\end{equation}
is solved by $x(t) = U(t, 0)y$ which satisfies the estimate $\| x(t) \| \leqq \| y \|$.
\end{theorem}

\begin{proof} By ($3$), \eqref{eq:XIV.4.6} and \eqref{eq:XIV.4.1a}, we obtain 
\begin{equation}\label{eq:XIV.4.12}
\| U _{k} (t, s) x \| \leqq \| x \| (k = 1, 2, \dots ,; 0 \leqq s \leqq t \leqq 1; x \in X).
\end{equation}

We also need, for the operator 
\begin{equation}\label{eq:XIV.4.13}
W _{k} (t, s) = A (t) U _{k} (t, s) A (s) ^{-1},
\end{equation}
the following estimate which is the key for the whole proof: 
\begin{equation}\label{eq:XIV.4.14}
\begin{yoslines}
&\| W _{k} (t, s) x \| \leqq (1 + k ^{-1} N) ^{2} \cdot \exp (N (t - s)) \cdot \| x \|, N = \sup\limits _{0 \leqq s, t \leqq 1, s \neq t} \| (t - s) ^{-1} C (t, s) \|.
\end{yoslines}
\end{equation}

To prove this, we shall rewrite $W_{k}(t,s)$, remembering \eqref{eq:XIV.4.1a} and the commutativity $A(s)^{-1} \exp((t - s)A(s)) = \exp((t - s)A(s))A(s)^{-1}$: 
\[
\begin{aligned}
&\quad W _{k} (t, s) = A (t) A \left(\frac {[ k t ]}{k}\right) ^{-1} U _{k} \left(t, \frac {[ k t ]}{k}\right) A \left(\frac {[ k t ]}{k}\right) A \left(\frac {[ k t ] - 1}{k}\right) ^{-1} U _{k} \left(\frac {[ k t ]}{k}, \frac {[ k t ] - 1}{k}\right) \dots
\\
&\quad \dots A \left(\frac {[ k s ] + 2}{k}\right) A \left(\frac {[ k s ] + 1}{k}\right) ^{-1} U _{k} \left(\frac {[ k s ] + 2}{k}, \frac {[ k s ] + 1}{k}\right) A \left(\frac {[ k s ] + 1}{k}\right) A \left(\frac {[ k s ]}{k}\right) ^{-1} \times
\\
&\quad \times U _{k} \left(\frac {[ k s ] + 1}{k}, s\right) A \left(\frac {[ k s ]}{k}\right) A (s) ^{-1}.
\end{aligned}
\]

Expanding the right side and again remembering \eqref{eq:XIV.4.1a}, we obtain 
\[
\begin{aligned}
W _{k} (t, s)& = \left(I + C \left(t, \frac {[ k t ]}{k}\right)\right) \left\{U _{k} (t, s) + \sum\limits_{ku=[ks]+1} ^{[kt]} U _{k} (t, u) C \left(u, u - \frac {1}{k}\right) U _{k} (u, s) + \right.
\\
&\quad \left. + \sum\limits_{kv=[ks]+1} ^{[kt]} U _{k} (t, v) C \left(v, v - \frac {1}{k}\right) \sum\limits_{ku=[ks]+1} ^{[kv]} U _{k} (v, u) C \left(u, u - \frac {1}{k}\right) U _{k} (u, s) + \dots \right\} \times
\\
&\quad \times \left(I + C \left(\frac {[ k s ]}{k}, s\right)\right),
\end{aligned}
\]
that is, 
\begin{equation}\label{eq:XIV.4.15}
\begin{aligned} 
    W _{k} (t, s) & = \left(I + C \left(t, \frac {[ k t ]}{k}\right)\right) \left\{U _{k} (t, s) + W _{k} ^{(1)} (t, s) + W _{k} ^{(2)} (t, s) + \dots \right\} \\
&\quad \times \left(I + C \left(\frac {[ k s ]}{k}, s\right)\right), \\ 
    W _{k} ^{(1)} (t, s) & = \sum\limits_{ku=[ks]+1} ^{[kt]} U _{k} (t, u) C \left(u, u - \frac {1}{k}\right) U _{k} (u, s), \\ 
    W _{k} ^{(m+1)} (t, s) & = \sum\limits_{ku=[ks]+1} ^{[kt]} U _{k} (t, u) C \left(u, u - \frac {1}{k}\right) W _{k} ^{(m)} (u, s), (m = 1, 2, \dots , [ k t ] - 1). 
\end{aligned}
\end{equation}

We have $\left\| C\left(u, u - \frac{1}{k}\right)x\right\| \leqq \frac{1}{k} N \| x\|$ by \eqref{eq:XIV.4.14}. Hence from \eqref{eq:XIV.4.12} we conclude that 
\begin{equation}\label{eq:XIV.4.16}
\begin{aligned}
&\left\| W _{k} ^{(1)} (t, s) x \right\| \leqq (t - s) N \| x \|, \dots \left\| W _{k} ^{(m)} (t, s) x \right\| \leqq \frac {(t - s) ^{m}}{m !} N ^{m} \| x \|.
\end{aligned}
\end{equation}

Combined with the definition of $N$, we have proved \eqref{eq:XIV.4.14} by \eqref{eq:XIV.4.12}.

Incidentally, we have proved that 
\begin{equation}\label{eq:XIV.4.17}
U _{k} (t, s) y \in D (A (t)) \quad \text{ when } \quad y \in D (A (s)).
\end{equation}

Hence 
\[
U _{k} (t, s) y = \exp \left(\left(t - \frac {[ k t ]}{k}\right) A \left(\frac {[ k t ]}{k}\right)\right) U _{k} \left(\frac {[ k t ]}{k}, s\right) y
\]
is differentiable at $t \neq \frac{i}{k} (i = 0, 1, \ldots, k)$ and 
\begin{equation}\label{eq:XIV.4.18}
\frac {\dd U _{k} (t , s) y}{\dd t} = A \left(\frac {[ k t ]}{k}\right) U _{k} (t, s) y, y \in D (A (0)).
\end{equation}

Similarly, $U_{k}(t,s)y$ is differentiable at $s\neq \frac{i}{k} (i = 0,1,\dots ,k)$ and 
\begin{equation}\label{eq:XIV.4.19}
\frac {\dd U _{k} (t , s) y}{\dd s} = - U _{k} (t, s) A \left(\frac {[ k s ]}{k}\right) y, y \in D (A (0)).
\end{equation}

These derivatives are bounded and are strongly continuous in $t$ and $s$ except at $t = \frac{i}{k}$ and $s = \frac{i}{k}$. For the proof, rewrite the right side of \eqref{eq:XIV.4.18} as $= C\left(\frac{[kt]}{k}, t\right) W_k(t, s) A(s) A(0)^{-1}x, y = A(0)^{-1}x$, and make use of ($5$), \eqref{eq:XIV.4.9}, \eqref{eq:XIV.4.12}, \eqref{eq:XIV.4.14} and \eqref{eq:XIV.4.15}; similarly for \eqref{eq:XIV.4.19}. 

Since $U_{n}(t,s)U_{k}(s,0)A(0)^{-1}x$ is strongly continuous in $s$ by \eqref{eq:XIV.4.9}, we have, remembering the fact $U_{k}(s,0)A(0)^{-1}x\in D(A(s))$ together with \eqref{eq:XIV.4.18} and \eqref{eq:XIV.4.19}, 
\[
\begin{aligned}
&\left(U _{k} (t, 0) - U _{n} (t, 0)\right) A (0) ^{-1} x = [ U _{n} (t, s) U _{k} (s, 0) A (0) ^{-1} x ] _{s=0} ^{s=t}
\\
&= \int_{0} ^{t} \frac {d}{\dd s} \left\{U _{n} (t, s) U _{k} (s, 0) A (0) ^{-1} x \right\} \dd s = \int_{0} ^{t} U _{n} (t, s) \left\{A \left(\frac {[ k s ]}{k}\right) - A \left(\frac {[ n s ]}{n}\right) \right\} A \left(\frac {[ k s ]}{k}\right) ^{-1}
\\
&\quad A \left(\frac {[ k s ]}{k}\right) U _{k} (s, 0) A (0) ^{-1} x \dd s
\\
&= \int_{0} ^{t} - U _{n} (t, s) C \left(\frac {[ n s ]}{n}, \frac {[ k s ]}{k}\right) A \left(\frac {[ k s ]}{k}\right) A (s) ^{-1} W _{k} (s, 0) x \dd s.
\end{aligned}
\]

Hence, \eqref{eq:XIV.4.12} and \eqref{eq:XIV.4.14} yield 
\begin{gather*}{\| U _{k} (t, 0) A (0) ^{-1} x - U _{n} (t, 0) A (0) ^{-1} x \|} \leqq  \int_{0} ^{t} \left\| C \left(\frac {[ n s ]}{n}, \frac {[ k s ]}{k}\right) \left(I + C \left(\frac {[ k s ]}{k}, s\right)\right) W _{k} (s, 0) x \right\| d s \\\leqq  N \int_{0} ^{t} \left| \frac {[ k s ]}{k} - \frac {[ n s ]}{n} \right| \cdot \left\| \left(I + C \left(\frac {[ k s ]}{k}, s\right)\right) \right. \cdot \left. \cdot W _{k} (s, 0) x \right\| d s \\\leqq  N \int_{0} ^{t} \left| \frac {[ k s ]}{k} - \frac {[ n s ]}{n} \right| \cdot \left(1 + N \left| s - \frac {[ k s ]}{k} \right|\right) g \cdot \left(1 + \frac {1}{k} N\right) \exp (N s) \cdot \| x \| d s.
\end{gather*}

This proves that $\slim\limits_{k \to \infty} U_k(t, 0) A(0)^{-1} x$ exists uniformly in $t, 0 \leqq t \leqq 1$, 
and so, by (2) and \eqref{eq:XIV.4.12}, we prove that, for every $x \in X$, $\slim\limits_{k \to \infty} U_k(t, 0) x = U(t, 0) x$ exists uniformly in $t, 0 \leqq t \leqq 1$. 

Similar argument shows that, for every $x \in X$,
\begin{equation}\label{eq:XIV.4.20}
\parbox{0.92\linewidth}{\everymath{}\let\limits\nolimits \slim $U_{k}(t,s)x = U(t,s)x$ exists uniformly in $t$ and $s,0 \leqq s \leqq t \leqq 1$.}
\end{equation} 

Therefore, by \eqref{eq:XIV.4.9}, we see that $U(t, s)$ is uniformly strongly continuous in $t$ and $s$. 

We also prove easily, remembering \eqref{eq:XIV.4.15}, \eqref{eq:XIV.4.16} and ($20$), that as $k \to \infty$ $W_{k}(t, 0)x$ strongly converges boundedly in $t$ such that 
\begin{equation}\label{eq:XIV.4.21}
\begin{yoslines}& \slim _{k \rightarrow \infty} W _{k} (t, 0) x = W (t, 0) x = U (t, 0) x + W ^{(1)} (t, 0) x + W ^{(2)} (t, 0) x \dots ,\\& \text{ where } W ^{(1)} (t, 0) x = \int_{0} ^{t} U (t, s) C (s) U (s, 0) x \dd s,\\& \dots\\& W ^{(m+1)} (t, 0) x = \int_{0} ^{t} U (t, s) C (s) W ^{(m)} (s, 0) x \dd s \quad (m = 1, 2, \dots).\end{yoslines}
\end{equation}

Therefore, if $y \in D(A(0))$ then the limits 
$\slim\limits_{k \to \infty} U_k(t,0) y = U(t,0) y$ and $\slim\limits_{k \to \infty} A(t) U_k(t,0) y = W(t,0) A(0) y$ both exist boundedly and uniformly in $t$ and $W(t,0) A(0) y$ is uniformly strongly continuous in $t$. $A(t)$ being a closed linear operator by ($3$), we have proved that, for every $y \in D(A(0))$, 
\begin{equation}\label{eq:XIV.4.22}
\begin{yoslines}& U (t, 0) y \in D (A (t)) \quad \text{ and }\\& A (t) U (t, 0) y = \slim\limits_{k \to \infty} A (t) A \left(\frac {[ k t ]}{k}\right) ^{-1} A \left(\frac {[ k t ]}{k}\right) U _{k} (t, 0) y = W (t, 0) A (0) y,\\& \text{ where the \slim \, exists boundedly and uniformly in } t, 0 \leqq t \leqq 1.\end{yoslines}
\end{equation}
Hence, by letting $k \to \infty$ in 
\[
U _{k} (t, 0) y - y = \int_{0} ^{t} \left(\frac {d}{\dd s} U _{k} (s, 0) y\right) \dd s = \int_{0} ^{t} A \left(\frac {[ k s ]}{k}\right) U _{k} (s, 0) y \dd s
\]

We obtain 
\[
U (t, 0) y - y = \int_{0} ^{t} A (s) U (s, 0) y \dd s.
\]


As proved in \eqref{eq:XIV.4.22}, the integrand is strongly continuous in $s$ so that we have proved that $x(t) = U(t,0)y$ solves Cauchy problem \eqref{eq:XIV.4.1'}.
\end{proof}

\begin{remarku} Theorem~\ref{thm:XIV.4.Theorem.1} and its proof are adapted from K. \textsc{Yosida} \cite{yosida373}. It is to be noted here that an inequality similar to our \eqref{eq:XIV.4.14} was obtained in T. \textsc{Kato} \cite{kato136} by solving a Volterra type integral equation. The appropriateness of the equi-partition of the interval [$0$,$1$] as revealed in \eqref{eq:XIV.4.1a} is suggested by a paper by J. \textsc{Kisyński} \cite{kisynski150}. 
\end{remarku}

The approximation $x_{k}(t) = U_{k}(t,0)y$ to the true solution $x(t) = U(t,0)y$ of \eqref{eq:XIV.4.1'} is a kind of difference approximation to the differential equation given in \eqref{eq:XIV.4.1'}. However, the true difference approximation would be to use the scheme 
\begin{equation}\label{eq:XIV.4.23}
\frac {x \left(t _{j}\right) - x \left(t _{j-1}\right)}{t _{j} - t _{j-1}} = A \left(t _{j}\right) x \left(t _{j}\right),
\end{equation}
choosing the backward difference. This scheme gives 
\[
x \left(t _{j}\right) = \left(I - \left(t _{j} - t _{j-1}\right) A \left(t _{j}\right)\right) ^{-1} x \left(t _{j-1}\right).
\]

Therefore, suggested by \eqref{eq:XIV.4.1a}, we are lead to the approximation 
\begin{equation}\tag{10'}\label{eq:XIV.4.10'}
\begin{yoslines}
&V _{k} (t, s) = \left(I - \left(t - \frac {[ k t ]}{k}\right) A \left(\frac {[ k t ]}{k}\right)\right) ^{-1} \left(I - \frac {1}{k} A \left(\frac {[ k t ] - 1}{k}\right)\right) ^{-1} \dots \\
&\dots \left(I - \frac {1}{k} A \left(\frac {[ k s ] + 1}{k}\right)\right) ^{-1} \left(I - \left(\frac {[ k s ] + 1}{k} - s\right) A \left(\frac {[ k s ]}{k}\right)\right) ^{-1}, (0 \leqq s \leqq t \leqq 1).
\end{yoslines}
\end{equation}

From the point of view of the numerical analysis, this approximation $V_{k}(t,s)$ would be much more practical than the approximation $U_{k}(t,s)$, since in the construction of the approximation $U_{k}(t,s)$ we have to appeal to the construction of $\exp(tA(s))$. 

We shall prove (Cf. the last section in T. \textsc{Kato} \cite{kato143}). 

\begin{theorem}\label{thm:XIV.4.Theorem.2} Under the same conditions (2) through ($5$), the Cauchy problem given in \eqref{eq:XIV.4.1'} is, for every initial data $y \in D(A(0))$, solved by 
\[
x (t) = \slim\limits_{k \rightarrow \infty} V _{k} (t, 0) y.
\]
\end{theorem}

\begin{proof}
 From ($3$) it follows that 
\begin{equation}\tag{12'}\label{eq:XIV.4.12'}
\| V _{k} (t, s) x \| \leqq \| x \| \quad (k = 1, 2, \dots ; 0 \leqq s \leqq t \leqq 1; x \in X). \end{equation}

We next see, from 
\[
V _{k} (t, s) = V _{k} \left(t, \frac {[ k s ] + 1}{k}\right) \left(I - \left(\frac {[ k s ] + 1}{k} - s\right) A \left(\frac {[ k s ]}{k}\right)\right) ^{-1},
\]
that $V_{k}(t, s)$ is differentiable at $s \neq \frac{i}{k} (i = 0, 1, \ldots, k)$ and 
\begin{equation}\tag{19'}\label{eq:XIV.4.19'}
\frac {\dd V _{k} (t , s) y}{\dd s} = - V _{k} (t, s) \left(I - \left(\frac {[ k s ] + 1}{k} - s\right) A \left(\frac {[ k s ]}{k}\right)\right) ^{-1} A \left(\frac {[ k s ]}{k}\right) y.
\end{equation}

We see, by $y \in D(A(0)) = D(A(s))$, ($3$) and \eqref{eq:XIV.4.12'}, that this derivative is bounded and strongly continuous in $s$ except at $s = \frac{i}{k}$. Since 
$V_{k}(t,s)U_{k}(s,0)A(0)^{-1}x$ is strongly continuous in $s$, we obtain, remembering the fact $U_{k}(s,0)A(0)^{-1}x\in D(A(0))$, 
\[
\begin{aligned}
& \left(U _{k} (t, 0) - V _{k} (t, 0)\right) A (0) ^{-1} x = \left[ V _{k} (t, s) U _{k} (s, 0) A (0) ^{-1} x \right] _{s=0} ^{s=t}
\\
= & \int_{0} ^{t} \frac {d}{\dd s} \left(V _{k} (t, s) U _{k} (s, 0) A (0) ^{-1} x\right) \dd s
\\
= & \int_{0} ^{t} V _{k} (t, s) \left(A \left(\frac {[ k s ]}{k}\right)\right. \left.- A \left(\frac {[ k s ]}{k}\right) \left(I - \left(\frac {[ k s ] + 1}{k} - s\right) A \left(\frac {[ k s ]}{k}\right)\right) ^{-1}\right) U _{k} (s, 0) A (0) ^{-1} x \dd s
\\
= & \int_{0} ^{t} V _{k} (t, s) \left(I - \left(I - \left(\frac {[ k s ] + 1}{k} - s\right) A \left(\frac {[ k s ]}{k}\right)\right) ^{-1}\right)\left(I + C \left(\frac {[ k s ]}{k}, s\right)\right) W _{k} (s, 0) x \dd s.
\end{aligned}
\]

Thus we have 
\begin{equation}\label{eq:XIV.4.24}
\begin{aligned}
& \left\| \left(U _{k} (t, 0) - V _{k} (t, 0)\right) A (0) ^{-1} x \right\|
\\
&\leqq \int_{0} ^{t} \left\| \left(I - \left(I - \left(\frac {[ k s ] + 1}{k} - s\right) A \left(\frac {[ k s ]}{k}\right)\right) ^{-1}\right) \left(I + C \left(\frac {[ k s ]}{k}, s\right)\right) W _{k} (s, 0) x \right\| \dd s.
\end{aligned}
\end{equation}

From ($5$) and \eqref{eq:XIV.4.21}, it follows that 
\[
\slim _{k \rightarrow \infty} \left(I + C \left(\frac {[ k s ]}{k}, s\right)\right) W _{k} (s, 0) x = W (s, 0) x \text{ uniformly in } s, 0 \leqq s \leqq 1.
\]

On the other hand, we have, for $z \in D(A(s))$, 
\begin{equation}\label{eq:XIV.4.25}
\begin{aligned} & - \left(I - \left(I - \left(\frac {[ k s ]}{k} - s\right) A \left(\frac {[ k s ]}{k}\right)\right) ^{-1}\right) z \\ & = \left(\frac {[ k s ]}{k} - s\right) \left(I - \left(\frac {[ k s ]}{k} - s\right) A \left(\frac {[ k s ]}{k}\right)\right) ^{-1} A \left(\frac {[ k s ]}{k}\right) A (s) ^{-1} A (s) z. \end{aligned}
\end{equation}

By ($3$) and ($5$), we see that this tends strongly to $0$ as $k \to \infty$. The integrand in \eqref{eq:XIV.4.24} is bounded in $k$ and $s$. Thus, since $D(A(s)) = D(A(0))$ is dense in $X$, we see from \eqref{eq:XIV.4.25} that 
\[
\slim _{k \to \infty} (U _{k} (t, 0) - V _{k} (t, 0)) A (0) ^{-1} x = 0.
\]

Therefore, by virtue of Theorem~\ref{thm:XIV.4.Theorem.1}, we have proved Theorem~\ref{thm:XIV.4.Theorem.2}. 
\end{proof}

\begin{remarku} The above proof is communicated from H. \textsc{Fujita} who is suggested by the last section of T. \textsc{Kato} \cite{kato143}. 
\end{remarku}

Other approaches. K. \textsc{Yosida} \cite{yosida366} and \cite{yosida371} devised the idea of approximating Cauchy problem \eqref{eq:XIV.4.1'} by 
\begin{equation}\tag{1''}\label{eq:XIV.4.1''}
\frac {\dd x _{k} (t)}{\dd t} = A (t) \left(I - \frac {1}{k} A (t)\right) ^{-1} x _{k} (t), x _{k} (0) = y \in D (A (0)), 0 \leqq t \leqq 1.
\end{equation}

Applying the Cauchy polygon method to equation \eqref{eq:XIV.4.1''}, J. \textsc{Kisyński} \cite{kisynski150} proved the strong convergence of $\{x_{k}(t)\}$ to the solution of \eqref{eq:XIV.4.1'}. Another method of integration of \eqref{eq:XIV.4.1} has been devised by J. L. \textsc{Lions} \cite{liggett179}. He assumes the operator $A(t)$ to be an elliptic differential operator with smooth coefficients depending on $t$, and seeks distributional solutions by transforming equation \eqref{eq:XIV.4.1} into an integrated form in concrete function spaces such as the Sobolev space $W^{k,p}(\Omega)$ and its variants. We also refer the reader to a paper by O. A. \textsc{Ladyzhenskaya}-I. M. \textsc{Visik} \cite{ladyzhenskaya170} which is motivated by a similar idea as Lions'.\yosstarfootnote{See also Supplementary Notes, p.~\pageref{notes:XIV}.} 


\yossection{5}{The Method of Tanabe and Sobolevski}
\label{sec:XIV.5}

Let $X$ be a complex $B$-space, and consider the equation of evolution in $X$ with a given \emph{inhomogeneous term} $f(t)$: 
\begin{equation}\label{eq:XIV.5.1}
\dd x (t) / \dd t = A x (t) + f (t), a \leqq t \leqq b.
\end{equation}

Then the solution $x(t) \in X$ with the initial condition $x(a) = x_0 \in X$ is given formally by the so-called \emph{Duhamel Principle} from the solution $\exp((t - a)A)x$ of the \emph{homogeneous equation} $\dd x/\dd t = Ax$: 
\begin{equation}\label{eq:XIV.5.2}
x (t) = \exp ((t - a) A) x _{0} + \int_{a} ^{t} \exp ((t - s) A) \cdot f (s) \dd s.
\end{equation}

This suggests that a temporally inhomogeneous equation in $X$: 
\begin{equation}\label{eq:XIV.5.3}
\dd x (t) / \dd t = A (t) x (t), a \leqq t \leqq b,
\end{equation}
may be solved formally as follows. We rewrite equation \eqref{eq:XIV.5.3} in the form 
\begin{equation}\label{eq:XIV.5.4}
\dd x (t) / \dd t = A (a) x (t) + (A (t) - A (a)) x (t).
\end{equation}

By virtue of the formalism \eqref{eq:XIV.5.2}, the solution $x(t)$ of \eqref{eq:XIV.5.4} with the initial condition $x(a) = x_{0}$ will be given as the solution of an abstract integral equation 
\begin{equation}\label{eq:XIV.5.5}
\begin{aligned}
&x (t) = \exp ((t - a) A (a)) x _{0} + \int_{a} ^{t} \exp ((t - s) A (s)) (A (s) - A (a)) x (s) \dd s.
\end{aligned}
\end{equation}

Solving \eqref{eq:XIV.5.5} formally by successive approximation, we obtain approximate solutions: 
\[
\begin{aligned}
&x _{1} (t) = \exp ((t - a) A (a)) x _{0}, \dots x _{n+1} (t) = \exp ((t - a) A (a)) x _{0}
\\
&+ \int_{a} ^{t} \exp ((t - s) A (s)) (A (s) - A (a)) x _{n} (s) \dd s.
\end{aligned}
\]

Hence the solution $x(t)$ of \eqref{eq:XIV.5.5} would be given formally by 
\begin{equation}\label{eq:XIV.5.6}
x (t) = \exp ((t - a) A (a)) x _{0} + \int_{a} ^{t} \exp ((t - s) A (s)) R (s, a) x _{0} \dd s,
\end{equation}
where 
\begin{equation}\label{eq:XIV.5.7}
\begin{aligned}
&R (t, s) = \sum\limits_{m=1} ^{\infty} R _{m} (t, s), R _{1} (t, s) = (A (t) - A (s)) \exp ((t - s) A (s)), s < t, = 0, s \geqq t,
\\
&R _{m} (t, s) = \int_{s} ^{t} R _{1} (t, \sigma) R _{m-1} (\sigma , s) \dd \sigma \quad (m = 2, 3, \dots).
\end{aligned}
\end{equation}

Justification for the above formal procedure of integration has been given by H. \textsc{Tanabe} \cite{tanabe315} using the theory of holomorphic semi-groups as given in Chapter~\ref{chap:IX}, \ref{sec:IX.10}. We shall follow Tanabe's approach, and assume the following conditions: 

\begin{equation}\label{eq:XIV.5.8}
\parbox{0.92\linewidth}{\everymath{}\let\limits\nolimits For each $t \in [a, b]$, $A(t)$ is a closed linear operator with domain dense in $X$ and range $X$ such that the resolvent set $\varrho(A(t))$ of $A(t)$ contains a fixed angular domain $\Theta$ of the complex $\lambda$-plane consisting of the origin $0$ plus the set $\{\lambda; -\theta < \arg \lambda < \theta \text{ with } \theta > \pi/2\}$. The resolvent $(\lambda I - A(t))^{-1}$ is strongly continuous in $t$ uniformly in $\lambda$ on any compact set $\subseteq \Theta$.}
\end{equation}

\begin{equation}\label{eq:XIV.5.9}
\parbox{0.92\linewidth}{\everymath{}\let\limits\nolimits There exists positive constants $M$ and $N$ such that, for $\lambda \in \Theta$ and $t \in [a, b]$, we have $\|(\lambda I - A(t))^{-1}\| \leqq N(|\lambda| - M)^{-1}$ whenever $|\lambda| > M$ with $N = 1$ for real $\lambda$.}
\end{equation}

\begin{equation}\label{eq:XIV.5.10}
\parbox{0.92\linewidth}{\everymath{}\let\limits\nolimits The domain $D(A(t))$ of $A(t)$ is independent of $t$ so that, by the closed graph theorem in Chapter~\ref{chap:II}, \ref{sec:II.6}, the operator $A(t)A(s)^{-1}$ is in $L(X,X)$. It is assumed that there exists a positive constant $K$ such that $\|A(t)A(s)^{-1} - A(r)A(s)^{-1}\| \leqq K|t - r|$ for $s, t$ and $r \in [a,b]$.}
\end{equation}



Under these conditions, we can prove 

\begin{theoremu} For any $x_0 \in X$ and $s$ with $a \leqq s < b$, the equation 
\begin{equation}\tag{3'}\label{eq:XIV.5.3'}
\dd x (t) / \dd t = A (t) x (t), x (s) = x _{0}, s < t \leqq b
\end{equation}
admits a uniquely determined solution $x(t) \in \mathbf{X}$. This solution is given by 
\begin{align}
x (t)& = U (t, s) x (s) = U (t, s) x _{0}, \text{ where } \label{eq:XIV.5.11} \\
&\left.\begin{yoslines}& U (t, s) = \exp ((t - s) A (s)) + W (t, s),\\& W (t, s) = \int_{s} ^{t} \exp ((t - \sigma) A (\sigma)) R (\sigma , s) \dd \sigma \text{ with } R (t, s) \quad \text{ given  by \eqref{eq:XIV.5.7}. }\end{yoslines}\right\} \label{eq:XIV.5.12}
\end{align}
\end{theoremu}

For the proof, we prepare three lemmas. 

\begin{lemma}\label{lem:XIV.5.Lemma.1} $R(t,s)$ satisfies, with a constant $C$, 
\begin{equation}\label{eq:XIV.5.13}
\left\| R (t, s) \right\| \leqq K C \cdot \exp (K C (t - s)),
\end{equation}
and $R(t, s)$ is strongly continuous in $a \leqq s < t \leqq b$. 
\end{lemma}

\begin{proof}
 By \eqref{eq:XIV.5.8} and \eqref{eq:XIV.5.9}, we see that each $A(s)$ generates a holomorphic semi-group which is given by (see Chapter~\ref{chap:IX}, \ref{sec:IX.10}) 
\begin{equation}\label{eq:XIV.5.14}
\exp (t A (s)) = (2 \pi i) ^{-1} \int_ {C ^{\prime}} e ^{\lambda t} (\lambda I - A (s)) ^{-1} \dd \lambda , \text{ where } C ^{\prime} \text{ is }
\end{equation}
a smooth contour running from $\infty e^{-i\theta}$ to $\infty e^{i\theta}$ in $\Theta$. 

Hence, by $A(s)(\lambda I - A(s)) = \lambda (\lambda I - A(s))^{-1} - I$, we have, for $(b - a) > t > 0$, 
\begin{equation}\label{eq:XIV.5.15}
\begin{aligned}
&\left\| \exp (t A (s)) \right\| \leqq C \quad \text{ and } \quad \left\| A (s) \exp (t A (s)) \right\| \leqq C t ^{-1}, \\
&\qquad \text{where the positive constant } C \text{ is independent of } t > 0 \text{ and } s\in [a,b].
\end{aligned}
\end{equation}


We have, by \eqref{eq:XIV.5.7}, 
\[
R _{1} (t, s) = \left(A (t) - A (s)\right) A (s) ^{-1} A (s) \exp ((t - s) A (s)), t > s,
\]
and so, by \eqref{eq:XIV.5.10} and \eqref{eq:XIV.5.15}, 
\begin{equation}\label{eq:XIV.5.16}
\left\| R _{1} (t, s) \right\| \leqq K C.
\end{equation}

It is also clear, from \eqref{eq:XIV.5.8} and \eqref{eq:XIV.5.14}, that $R_{1}(t,s)$ is strongly continuous in $a \leqq s < t \leqq b$. Next, by induction, we obtain 
\[
\begin{aligned}
&\left\| R _{m} (t, s) \right\|  \leqq \int_{s} ^{t} \left\| R _{1} (t, \sigma) \right\| \cdot \left\| R _{m-1} (\sigma , s) \right\| \dd \sigma \leqq \int_{s} ^{t} (K C) ^{m} (\sigma - s) ^{m-2} (\underline {{m - 2}}) ^{-1} \dd \sigma \\
&= (K C) ^{m} (t - s) ^{m-1} (| m - 1) ^{-1},
\end{aligned}
\]
and hence \eqref{eq:XIV.5.13} is obtained. In the same way, we see that $R(t, s)$ is strongly continuous in $a \leqq s < t \leqq b$. 
\end{proof}

\begin{lemma}\label{lem:XIV.5.Lemma.2} For $s < \tau < t$, we have 
\begin{equation}\label{eq:XIV.5.17}
\left\| R (t, s) - R (\tau , s) \right\| \leqq C _{1} \left(\frac {t - \tau}{t - s} + (t - \tau) \log \frac {t - s}{t - \tau}\right),
\end{equation}
where $C_1$ is a positive constant which is independent of $s, \tau$ and $t$. 
\end{lemma}

\begin{proof}
 We have, by \eqref{eq:XIV.5.7}, 
\[
\begin{aligned}
R _{1} (t, s) - R _{1} (\tau , s)& = \left(A (t) - A (\tau)\right) \exp \left((t - s) A (s)\right)
\\
&\quad \quad + \left(A (\tau) - A (s)\right) \left[ \exp \left((t - s) A (s)\right) - \exp \left((\tau - s) A (s)\right) \right]. \end{aligned}
\]

By \eqref{eq:XIV.5.10} and \eqref{eq:XIV.5.15}, the norm of the first term on the right is dominated by $KC(t - \tau)(t - s)^{-1}$. The second term on the right is 
\[
\begin{aligned}
&= (A (\tau) - A (s)) \int_{\tau - s} ^{t-s} \frac {d}{\dd \sigma} \exp (\sigma A (s)) = (A (\tau) - A (s)) A (s) ^{-1} \times \int_{\tau - s} ^{t-s} A (s) ^{2} \exp (\sigma A (s)) \dd \sigma ,
\end{aligned}
\]
and we have, by \eqref{eq:XIV.5.15}, 
\[
\begin{aligned}
&\left\| \int_{\tau - s} ^{t-s} A (s) ^{2} \exp (\sigma A (s)) \dd \sigma \right\| \leqq \int_{\tau - s} ^{t-s} \| (A (s) \exp (2 ^{-1} \sigma A (s))) ^{2} \| \dd \sigma \leqq \int_{\tau - s} ^{t-s} (2 C / \sigma) ^{2} \dd \sigma
\\
&\quad = 4 C ^{2} \left[ \frac {- 1}{\sigma} \right] _{\tau - s} ^{t-s} = 4 C ^{2} (t - \tau) (t - s) ^{-1} (\tau - s) ^{-1}.
\end{aligned}
\]

Therefore 
\begin{equation}\label{eq:XIV.5.18}
\left\| R _{1} (t, s) - R _{1} (\tau , s) \right\| \leqq K C (1 + 4 C) \frac {t - \tau}{t - s}.
\end{equation}

On the other hand, by \eqref{eq:XIV.5.7}, 
\[
\begin{aligned}
&\sum\limits_{m=2} ^{\infty} R _{m} (t, s) - \sum\limits_{m=2} ^{\infty} R _{m} (\tau , s) = \int_{s} ^{t} R _{1} (t, \sigma) R (\sigma , s) \dd \sigma
\\
&- \int_{s} ^{\tau} R _{1} (\tau , \sigma) R (\sigma , s) \dd \sigma = \int_{\tau} ^{t} R _{1} (t, \sigma) R (\sigma , s) \dd \sigma + \int_{s} ^{\tau} (R _{1} (t, \sigma) - R _{1} (\tau , \sigma)) R (\sigma , s) \dd \sigma .
\end{aligned}
\]

The norm of the first term on the right is dominated by 
\[
\int_{\tau} ^{t} \left\| R _{1} (t, \sigma) \right\| \left\| R (\sigma , s) \right\| \dd \sigma \leqq K ^{2} C ^{2} \exp (K C (b - a)) (t - \tau).
\]

The norm of the second term on the right is, by \eqref{eq:XIV.5.13} and \eqref{eq:XIV.5.18}, dominated by 
\[
\begin{aligned}
&\int_{s} ^{\tau} \left\| R _{1} (t, \sigma) - R _{1} (\tau , \sigma) \right\| \left\| R (\sigma , s) \right\| \dd \sigma
\\
&\leqq K ^{2} C ^{2} (1 + 4 C) \exp (K C (b - a)) \int_{s} ^{\tau} (t - \tau) (t - \sigma) ^{-1} \dd \sigma = K _{1} (t - \tau) \log \frac {t - s}{t - \tau}.
\end{aligned}
\]

Therefore we obtain \eqref{eq:XIV.5.17}. 
\end{proof}

\begin{lemma}\label{lem:XIV.5.Lemma.3} For $s < t$, we have 
\begin{equation}\label{eq:XIV.5.19}
\begin{aligned} & {\left\| A (t) \left\{\exp ((t - s) A (t)) - \exp ((t - s) A (s)) \right\} \right\| \leqq C _{2}, \text{ where }} \\ & {\quad C _{2} \text{ is a positive constant independent of } s \text{ and } t.} \end{aligned}
\end{equation}
\end{lemma}

\begin{proof}
 We obtain, from \eqref{eq:XIV.5.14}, 
\[
\begin{gathered}
A (t) \left\{\exp ((t - s) A (t)) - \exp ((t - s) A (s)) \right\}
\\
= (2 \pi i) ^{-1} \int_{C} \dot {e} ^{\lambda (t - s)} A (t) (\lambda I - A (t)) ^{-1} (A (t) - A (s)) (\lambda I - A (s)) ^{-1} \dd \lambda . \end{gathered}
\]

On the other hand, we have $A(t)(\lambda I - A(t))^{-1} = \lambda (\lambda I - A(t))^{-1} - I$, and so, by \eqref{eq:XIV.5.9}, 
\begin{equation}\label{eq:XIV.5.20}
\left\| A (t) (\lambda I - A (t)) ^{-1} \right\| \leqq \frac {| \lambda |}{| \lambda | - M} + 1 \text{ for } \lambda \in \Theta \text{ and } t \in [ a, b ].
\end{equation}

Hence we obtain \eqref{eq:XIV.5.19} by \eqref{eq:XIV.5.10} and 
\[
\begin{aligned}
&{\left\| (A (t) - A (s)) (\lambda I - A (s)) ^{-1} \right\|} {\qquad \leqq \left\| (A (t) - A (s)) A (s) ^{-1} \right\| \left\| A (s) (\lambda I - A (s)) ^{-1} \right\|.}
\end{aligned}
\]
\end{proof}

\begin{proof}[Proof of the Theorem.]
We rewrite $W(t, s)$ given in \eqref{eq:XIV.5.12} as follows:
\[
\begin{aligned}
W (t, s)& = \int_{s} ^{t} \exp ((t - \tau) A (t)) R (t, s) \dd \tau
\\
\qquad +& \int_{s} ^{t} \left\{\exp ((t - \tau) A (\tau)) - \exp ((t - \tau) A (t)) \right\} R (\tau , s) \dd \tau
\\
\qquad +& \int_{s} ^{t} \exp ((t - \tau) A (t)) (R (\tau , s) - R (t, s)) \dd \tau .
\end{aligned}
\]

By approximating the integrals by Riemann sums and making use of the closure property of the operator $A(t)$, we see that we can apply $A(t)$ to each term of the right side of the above equality. For, by \eqref{eq:XIV.5.19}, we can apply $A(t)$ to the second term of the right side; and also to the third term of the right side by \eqref{eq:XIV.5.15} and \eqref{eq:XIV.5.17}; we also have, by $A(t) \exp((t - \tau)A(t)) = -\dd\exp((t - \tau)A(t))/\dd\tau$, 
\[
A (t) \int_{s} ^{t} \exp ((t - \tau) A (t)) R (t, s) \dd \tau = \left\{\exp ((t - s) A (t)) - I \right\} R (t, s).
\]

Hence we obtain 
\begin{equation}\label{eq:XIV.5.21}
\begin{aligned}
A (t) U (t, s)& = A (t) \exp ((t - s) A (s)) + \{\exp ((t - s) A (t)) - I \} R (t, s)
\\
\quad +& \int_{s} ^{t} A (t) \left\{\exp ((t - \tau) A (\tau)) - \exp ((t - \tau) A (t)) \right\} R (\tau , s) \dd \tau
\\
\quad +& \int_{s} ^{t} A (t) \exp ((t - \tau) A (t)) (R (\tau , s) - R (t, s)) \dd \tau .
\end{aligned}
\end{equation}

The above proof shows that $A(t)U(t,s)$ is strongly continuous in $a \leqq s < t \leqq b$ and 
\begin{equation}\label{eq:XIV.5.22}
\left\| A (t) W (t, s) \right\| \leqq C _{3} \text{ and } \left\| A (t) U (t, s) \right\| \leqq C _{3} (t - s) ^{-1},
\end{equation}
where $C_{3}$ is a positive constant independent of $s$ and $t$. 

We next define, for $s < (t - h) < t$, 
\begin{equation}\label{eq:XIV.5.23}
U _{h} (t, s) = \exp ((t - s) A (s)) + \int_{s} ^{t-h} \exp ((t - \tau) A (\tau)) R (\tau , s) \dd \tau .
\end{equation}

Since a holomorphic semi-group $\exp(tA(u))$ is differentiable in $t > 0$, we have 
\[
\begin{aligned}
\frac {\partial}{\partial t} U _{h} (t, s)& = A (s) \exp ((t - s) A (s)) + \exp (h A (t - h)) R (t - h, s)
\\
\qquad +& \int_{s} ^{t-h} A (\tau) \exp ((t - \tau) A (\tau)) R (\tau , s) \dd \tau .
\end{aligned}
\]

Hence we have, by \eqref{eq:XIV.5.7}, 
\begin{equation}\label{eq:XIV.5.24}
\begin{aligned}
&\frac {\partial}{\partial t} U _{h} (t, s) - A (t) U _{h} (t, s)  = \exp (h A (t - h)) R (t - h, s) - R _{1} (t, s) - \int_{s} ^{t-h} R _{1} (t, \tau) R (\tau , s) \dd \tau .
\end{aligned}
\end{equation}

By \eqref{eq:XIV.5.8}, \eqref{eq:XIV.5.13} and \eqref{eq:XIV.5.14}, $\exp(hA(t-h))R(t-h,s)$ tends strongly to $R(t,s)$ as $h\downarrow 0$. Thus we have 
\begin{equation}\label{eq:XIV.5.25}
\begin{gathered}
\slim _{h \downarrow 0} \left(\frac {\partial}{\partial t} U _{h} (t, s) - A (t) U _{h} (t, s)\right) x _{0}
\\
= \left(R (t, s) - R _{1} (t, s) - \int_{s} ^{t} R _{1} (t, \sigma) R (\sigma , s) \dd \sigma\right) x _{0}, x _{0} \in X.
\end{gathered}
\end{equation}

The right side must be $0$ as may easily be proved by \eqref{eq:XIV.5.7}. Since we have 
\[
\slim _{h \downarrow 0} A (t) U _{h} (t, s) x _{0} = A (t) U (t, s) x _{0}
\]
by the reasoning used in proving \eqref{eq:XIV.5.21}, we obtain from \eqref{eq:XIV.5.25} 
\begin{equation}\label{eq:XIV.5.26}
\slim _{h \downarrow 0} \frac {\partial}{\partial t} U _{h} (t, s) x _{0} = A (t) U (t, s) x _{0} \text{ for } t > s \text{ and } x _{0} \in X.
\end{equation}

The right side of \eqref{eq:XIV.5.26} being strongly continuous in $t > s$, we see, by integrating \eqref{eq:XIV.5.26} and remembering $\slim\limits_{h\downarrow 0}U_h(t,s)x_0 = U(t,s)x_0$, that 
\begin{equation}\label{eq:XIV.5.27}
\frac {\partial}{\partial t} U (t, s) x _{0} = A (t) U (t, s) x _{0} \text{ for } t > s \text{ and } x _{0} \in X.
\end{equation}

Therefore, $x(t) = U(t, s)$ $x_{0}$ is the desired solution of $\eqref{eq:XIV.5.3}$. The uniqueness of the solution may be proved as in the preceding section. 
\end{proof}

\yossub{Comments and References}

The above Theorem and proof are adapted from H. \textsc{Tanabe} \cite{tanabe315}. To illustrate his idea we have made Tanabe's conditions somewhat stronger. Thus, e.g., the condition \eqref{eq:XIV.5.9} may be replaced by a weaker one: 
\[
\left\| A (t) A (s) ^{-1} - A (r) A (s) ^{-1} \right\| \leqq K _{1} | t - r | ^{\varrho} \text{ with } 0 < \varrho < 1.
\]

For the details, see the above cited paper by H. \textsc{Tanabe} which is a refinement of H. \textsc{Tanabe} \cite{tanabe316} and \cite{tanabe317}. It is to be noted that the Russian school independently has developed a similar method. See, e.g., P. E. \textsc{Sobolevski} \cite{sobolevski308} and the reference cited in this paper. Cf. E. T. \textsc{Poulsen} \cite{poulsen264}. 

Komatsu's work. H. \textsc{Komatsu} \cite{komatsu153} made an important remark regarding Tanabe's result given above. Let $\Delta$ be a convex complex neighbourhood of the real segment $[a, b]$ considered as embedded in the complex plane. Suppose that $A(t)$ is defined for $t \in \Delta$ and satisfies \eqref{eq:XIV.5.8} and \eqref{eq:XIV.5.9} in which the phrase `` $t \in [a, b]$ '' is replaced by `` $t \in \Delta$ ''. Assume, moreover, that there exists a bounded linear operator $A_0$ which maps $X$ onto $D$, the domain of $A(t)$ which is assumed to be independent of $t \in \Delta$, in a one-to-one manner and such that $B(t) = A(t)A_0$ is strongly holomorphic in $t \in \Delta$. Under these assumptions, \textsc{Komatsu} proved that the operator $U(t, s)$ constructed as above is strongly holomorphic in $t \in \Delta$ if 
\[
\left| \arg (t - s) \right| < \theta_{0} \text{ with a certain } \theta_{0} \text{ satisfying } 0 < \theta_{0} < \pi / 2.
\]

The result may be applied to the ``forward and backward unique continuation'' of solutions of temporally inhomogeneous diffusion equations as in Chapter~\ref{chap:XIV}, \ref{sec:XIV.1}. In this connection, we refer to H. \textsc{Komatsu} \cite{komatsu154}, \cite{komatsu155} and T. \textsc{Kotaké}-M. \textsc{Narasimhan} \cite{konishi162}. 

Kato's work. To get rid of the assumption that the domain $D(A(t))$ is independent of $t$, T. \textsc{Kato} \cite{kato139} proved that, in the above Theorem, we may replace condition \eqref{eq:XIV.5.10} by the following: 

For a certain positive integer $k$, the domain $D\left((-A(t))^{1/k}\right)$ is independent of $t$. (Here $(-A(t))^{1/k}$ is the fractional power as defined in Chapter~\ref{chap:IX}, \ref{sec:IX.11}.) Further, there exist constants $K_{2} > 0$ and $\gamma$ with $1 - k^{-1} < \gamma \leqq 1$ such that, for $s, t \in [a, b]$, $\left\| (-A(t))^{1/k}(-A(s))^{1/k} - I \right\| \leqq K_{2}|t - s|^{\gamma}$. 

\yosparagraphtag{10'}\label{eq:XIV.5.10'} 

Tanabe's and the Kato-Tanabe recent work. With the same motivation as \textsc{Kato}, H. \textsc{Tanabe} \cite{kato134} devised a method to replace condition \eqref{eq:XIV.5.10} by 
$A(t)^{-1}$ is once strongly differentiable in $a \leqq t \leqq b$ 
and such that, for positive constants $K_{3}$ and $\alpha$, 
\begin{equation}\tag{10''}\label{eq:XIV.5.10''}
\left\| \frac {\dd A (t) ^{-1}}{\dd t} - \frac {\dd A (s) ^{-1}}{\dd s} \right\| \leqq K _{3} | t - s | ^{\alpha}.
\end{equation}

Further, there exist positive constants $N$ and $\varrho$ with $0 < \varrho \leqq 1$ such that 
\[
\left\| \frac {\partial}{\partial t} (\lambda I - A (t)) ^{-1} \right\| \leqq N | \lambda | ^{q-1}, q < 1.
\]

For details, see \textsc{Kato}-\textsc{Tanabe} \cite{kato141}. Their point is to start with the first approximation $\exp((t-a)A(t))x_0$ instead of $\exp((t-a)A(a))x_0$. 

Nelson's work on Feynman integrals. The semi-group method of integration of Schrödinger equations gives an interpretation of Feynman integrals. See E. \textsc{Nelson} \cite{nelson227}. 

The Agmon-Nirenberg work. Agmon-Nirenberg \cite{agmon3} discussed the behaviour as $t \uparrow \infty$ of solutions of the equation $\frac{1}{i} \frac{\dd u}{\dd t} - Au = 0$ in some $B$-spaces. 


\yossection{6}{Non-linear Evolution Equations 1 (The Kōmura-Kato Approach)}
\label{sec:XIV.6}

Let $X$ be a real or complex Banach space, endowed with a semi-scalar product $[x, y]$ such that (see Chapter~\ref{chap:IX}, \ref{sec:IX.8}) 
\begin{equation}\label{eq:XIV.6.1}
\begin{gathered}
\left[ \alpha_{1} x _{1} + \alpha_{2} x _{2}, y \right] = \alpha_{1} \left[ x _{1}, y \right] + \alpha_{2} \left[ x _{2}, y \right], | [ x, y ] | \leqq \| x \| \cdot \| y \| \quad \text{ and }
\\
{}[ x, x ] = \| x \| ^{2},
\end{gathered}
\end{equation}
and consider a family $\{T_t; t \geqq 0\}$ of non-linear mappings of $X$ into $X$ satisfying conditions: 

$\int T_{t}T_{s} = T_{t+s}, T_{0} =$ the identity mapping $I$ (semi-group property), 
$\left\|T_{t}x-T_{t}y\right\|\leqq\left\|x-y\right\|\text{ (contraction property) }$, 
and strong continuity in $t$ of $T_{t}$. As in the case of linear mappings, we define the infinitesimal generator $A$ of $\{T_{t}; t \geqq 0\}$ by 
\[
A \cdot x = \slim _{h \downarrow 0} h ^{-1} \left(T _{t} x - x\right).
\]

Then $A$ must be dissipative in the following sense: 
\begin{equation}\label{eq:XIV.6.2}
\operatorname{Re} [ A x - A y, x - y ] \leqq 0.
\end{equation}

The proof is easy, since we have 
\begin{gather}\operatorname{Re} \left[ h ^{-1} (T _{h} x - x) - h ^{-1} (T _{h} y - y), x - y \right] = {} h ^{-1} \operatorname{Re} \left[ T _{h} x - T _{h} y, x - y \right] - h ^{-1} [ x - y, x - y ] \notag \\\leqq {} h ^{-1} \| T _{h} x - T _{h} y \| \cdot \| x - y \| - h ^{-1} \| x - y \| ^{2} \leqq {} h ^{-1} \| x - y \| ^{2} - h ^{-1} \| x - y \| ^{2} = 0. \notag
\end{gather}

A celebrated example due to Y. \textsc{Kōmura} \cite{komura158} is this: Let $X = \mathbb{R}^{1}$ with $[x, y] = x \cdot y$ and $\|x\| = |x|$, and 
\[
T _{t} x = \begin{cases}\max (x - t, 0) & \text{ for } \quad x > 0, \\ x & \text{ for } \quad x \leqq 0,\end{cases} \quad A x = \begin{cases}- 1 & \text{ for } \quad x > 0, \\ 0 & \text{ for } \quad x \leqq 0.\end{cases}
\]

By virtue of \eqref{eq:XIV.6.2} we can prove that, for all $\lambda > 0$, the mapping $(I - \lambda A)$ of $D(A)$ into $X$ has an inverse $J_{\lambda} = (I - \lambda A)^{-1}$, because 
\[
x _{1} - \lambda A x _{1} = z = x _{2} - \lambda A x _{2}
\]
implies 
\[
\begin{aligned}
0 &= [ (x _{1} - \lambda A x _{1}) - (x _{2} - \lambda A x _{2}), x _{1} - x _{2} ]
\\
\quad &= [ x _{1} - x _{2}, x _{1} - x _{2} ] - \lambda [ A x _{1} - A x _{2}, x _{1} - x _{2} ] \text{ so that }
\\
0 &= \| x _{1} - x _{2} \| ^{2} - \lambda \operatorname{Re} [ A x _{1} - A x _{2}, x _{1} - x _{2} ] \geqq \| x _{1} - x _{2} \| ^{2}, \text{ i.e., } x _{1} = x _{2}.
\end{aligned}
\]

Therefore, the theory of linear contraction semi-groups (Chapter~\ref{chap:IX}, \ref{sec:IX.8}) suggests us to approximate the non-linear evolution equation 
$\frac{\dd u(t)}{\dd t} = A u(t)$ for $t \geqq 0$ with $u(0) = x_0 \in D(A)$ by equations $(\lambda > 0)$ $\frac{\dd u^{(\lambda)}(t)}{\dd t} = A_{\lambda} u^{(\lambda)}(t)$ for $t \geqq 0$ with $u^{(\lambda)}(0) = x_0$, where $A_{\lambda} = \lambda^{-1}(J_{\lambda} - 1)$, under the assumption that $D(J_{\lambda}) = R(I - \lambda A) = X$, expecting that we shall have $u(t) = \slim\limits_{\lambda \downarrow 0} u^{(\lambda)}(t)$. 

In the above example of \textsc{Kōmura}, we have $D(J_{\lambda}) = (-\infty, 0] \cup (\lambda, \infty)$ which does not coincide with $X = \mathbb{R}^{1}$. To obtain $D(J_{\lambda}) = X = \mathbb{R}^{1}$, it will be necessary to extend $A$ to a multi-valued mapping $\hat{A}$ by 
\[
\hat {A} x = \begin{cases}- 1 & \text{ for } x > 0, \\ - [ 0, 1 ] & \text{ for } x = 0, \\ 0 & \text{ for } x < 0,\end{cases}
\]
for which the dissipative property is preserved in the following sense: $\operatorname{Re}\langle y_1 - y_2, x_1 - x_2 \rangle \leqq 0$ for any $\{x_i, y_i\}$ with $y_i \in \hat{A} x_i (i = 1, 2)$. \yosparagraphtag{2'}\label{eq:XIV.6.2'} 

In this way, we are led to the following setup for our further discussions. Let $X$ be a Banach space whose dual space is denoted by $X'$. An element of $X \times X$ will be written in the form $\{x, y\}$ where $x$ and $y$ are both $\in X$. For a subset $A$ of $X \times X$, the following notations will conveniently be used: 

$1$) $D(A) = \{x; \{x, y\} \in A \text{ for some } y\}$, 
$2$) $R(A) = \{y; \{x, y\} \in A \text{ for some } x\}$, 
$3$) $A^{-1} = \{\{y, x\}; \{x, y\} \in A\}$, 
$4$) $\lambda A = \{\{x, \lambda y\}; \{x, y\} \in A$ and for real $\lambda\}$, 
$5$) $A + B = \{\{x, y + z\}; \{x, y\} \in A \text{ and } \{x, z\} \in B\}$, 
$6$) $A x = \{y; \{x, y\} \in A$ for a fixed $x \in D(A)\}$, 
$7$) $\|Ax\| = \inf \{|y\|; y \in Ax\}$, 
$8$) $J_{\lambda} = (I - \lambda A)^{-1} = \{\{x - \lambda y, x\} ; \{x, y\} \in A\}$ for real $\lambda$ and $I = \{\{x, x\} ; x \in X\}$, 
$9$) $A_{\lambda} = \{\{x - \lambda y, y\}; \{x, y\} \in A\}$; 

we have 
\begin{equation}\label{eq:XIV.6.3}
\lambda A _{\lambda} = J _{\lambda} - I.
\end{equation}

If $A x$ consists of a single element for all $x \in D(A)$, then $A$ is the graph of a uniquely determined function ( $= \text{ single-valued mapping }$) from $D(A)$ with values in $X$; in this case, we may identify the set $A$ and the above function. 

To state the notion of dissipative sets in $X \times X$, we give 

\begin{definition}\label{def:XIV.6.Definition.1} By the duality map of $X$ into $X'$ we mean the multi-valued mapping $F$ from $X$ into $X'$ defined by 
\begin{equation}\label{eq:XIV.6.4}
F (x) = \{f \in X ^{\prime}; \langle x, f \rangle = \| x \| ^{2} = \| f \| ^{2} \}.
\end{equation}
\end{definition}

That $F(x)$ is not void is clear from the \textsc{Hahn}-\textsc{Banach} theorem. If $X$ is a Hilbert space, then F. \textsc{Riesz}' representation theorem asserts that $F(x)$ consists of $x$ alone and $\langle y, F(x) \rangle = (y, x)$, the scalar product of $y$ and $x$. 

\begin{definition}\label{def:XIV.6.Definition.2} $A$ set $A$ in $X \times X$ is called a dissipative set if, for arbitrary two points $\{x_{1}, y_{1}\}$ and $\{x_{2}, y_{2}\}$ of $A$, there exists an $f \in F(x_{1} - x_{2})$ such that $\operatorname{Re}\langle y_{1} - y_{2}, f \rangle \leqq 0$. 
\end{definition}

\begin{remarku} $A$ is a dissipative set iff $-A$ is a monotonic set in the sense of G. \textsc{Minty} \cite{minty203}. 
\end{remarku}

We also give 

\begin{definition}\label{def:XIV.6.Definition.3} Let $D$ be a subset of $X$, and $T$ be a function from $D$ into $X$. $T$ is called a Lipschitzian mapping (function) with Lipschitz constant $k > 0$ if $\|Tx_{1} - Tx_{2}\| \leqq k\|x_{1} - x_{2}\|$ for all $x_{1}, x_{2} \in D$. If we can take $k = 1$, then $T$ is called a contraction mapping (function). We have 
\end{definition}

\begin{lemma}[T. \textsc{Kato} {\cite{kato144}}]\label{lem:XIV.6.Lemma.1} Let $x, y \in X$. Then $\|x - \lambda y\| \geqq \|x\|$ for all $\lambda > 0$ iff there exists an $f \in F(x)$ such that $\operatorname{Re}\langle y, f \rangle \leqq 0$. 
\end{lemma}

\begin{proof}
 The assertion is trivial if $x = 0$. So we shall assume $x \neq 0$ in the following. If $\operatorname{Re}\langle y, f \rangle \leqq 0$ for some $f \in F(x)$, then $\|x\|^2 = \langle x, f \rangle = \operatorname{Re}\langle x, f \rangle \leqq \operatorname{Re}\langle x - \lambda y, f \rangle \leqq \|x - \lambda y\| \cdot \|f\|$. Since $\|x\| = \|f\|$, we obtain $\|x\| \leqq \|x - \lambda y\|$. 

Suppose, conversely, that $\|x\| \leqq \|x - \lambda y\|$ for all $\lambda > 0$. For each $\lambda > 0$, let $f_{\lambda} \in F(x - \lambda y)$ and put $g_{\lambda} = f_{\lambda} / \|f_{\lambda}\|$ so that $\|g_{\lambda}\| = 1$. Then $\|x\| \leqq \|x - \lambda y\| = \langle x - \lambda y, g_{\lambda} \rangle = \operatorname{Re}\langle x, g_{\lambda} \rangle - \lambda \operatorname{Re}\langle y, g_{\lambda} \rangle \leqq \|x\| - \lambda \operatorname{Re}\langle y, g; \rangle$ by $\|g_{\lambda}\| = 1$. Hence 
\[
\liminf\limits _{\lambda \downarrow 0} \operatorname{Re} \left\langle x, g _{\lambda} \right\rangle = \| x \| \quad \text{ and } \quad - \lambda \operatorname{Re} \left\langle y, g _{\lambda} \right\rangle \geqq 0.
\]

Since the closed unit sphere of the dual space $X'$ is weakly* compact (see Theorem~\hyperref[thm:AV.1.Theorem.1]{$1$} in the Appendix to Chapter~\ref{chap:V}, \ref{sec:AV.1}) the sequence $\{g_{1/n}\}$ has a weak* accumulation point $g \in X'$ with $\|g\| \leqq 1$. Thus we see that $g$ must satisfy $\operatorname{Re}\langle x, g \rangle \geqq \|x\|$ and $\operatorname{Re}\langle y, g \rangle \leqq 0$, that is, we must have $\|g\| = 1$ and $\langle x, g \rangle = \|x\|$. Therefore $f = \|x\| \cdot g$ satisfies $f \in F(x)$ and $\operatorname{Re}\langle y, f \rangle \leqq 0$. 
\end{proof}

\begin{corollaryu} $A$ is a dissipative set iff 
\begin{equation}\label{eq:XIV.6.5}
\begin{aligned}
&{\| (x _{1} - \lambda y _{1}) - (x _{2} - \lambda y _{2}) \| \geqq \| x _{1} - x _{2} \| \text{ whenever } \lambda > 0} {\quad \text{ and } \quad \{x _{i}, y _{i} \} \in A (i = 1, 2).}
\end{aligned}
\end{equation}
\end{corollaryu}

We have the following 

\begin{proposition}\label{prop:XIV.6.Proposition.1} If $A$ is a dissipative set and $\lambda > 0$, then $J_{\lambda}$ and $A_{\lambda}$ are both single-valued mappings, and 
\begin{align}
\left\| J _{\lambda} x _{1} - J _{\lambda} x _{2} \right\|& \leqq \left\| x _{1} - x _{2} \right\| \quad \text{ for } \quad x _{1}, x _{2} \in D (J _{\lambda}), \label{eq:XIV.6.6} \\
\left\| A _{\lambda} x _{1} - A _{\lambda} x _{2} \right\|& \leqq \frac {2}{\lambda} \left\| x _{1} - x _{2} \right\| \quad \text{ for } \quad x _{1}, x _{2} \in D (A _{\lambda}) = D (J _{\lambda}). \label{eq:XIV.6.7}
\end{align}
moreover, $A_{\lambda}$ is dissipative and 
\begin{align}
A J _{\lambda} x& = A (J _{\lambda} x) \ni A _{\lambda} x \quad \text{ for } \quad x \in D (J _{\lambda}), \label{eq:XIV.6.8} \\
\left\| A _{\lambda} x \right\|& \leqq \left\| A x \right\| \quad \text{ for all } \quad x \in D (A) \cap D (J _{\lambda}). \label{eq:XIV.6.9}
\end{align}
\end{proposition}

\begin{proof}
 That $J_{\lambda}$ and $A_{\lambda}$ are both single-valued is clear from \eqref{eq:XIV.6.5}. We also have \eqref{eq:XIV.6.6} by \eqref{eq:XIV.6.5}, and thus \eqref{eq:XIV.6.7} is proved by \eqref{eq:XIV.6.3} and \eqref{eq:XIV.6.6}. Next let $f \in F(x_1 - x_2)$. Then we obtain, by \eqref{eq:XIV.6.3}, \eqref{eq:XIV.6.4}, and \eqref{eq:XIV.6.6}, 
\[
\begin{gathered}
\operatorname{Re} \left\langle A _{\lambda} x _{1} - A _{\lambda} x _{2}, f \right\rangle = \lambda^{-1} \operatorname{Re} \left\langle (J _{\lambda} x _{1} - x _{1}) - (J _{\lambda} x _{2} - x _{2}), f \right\rangle
\\
= \lambda^{-1} \operatorname{Re} \left\langle J _{\lambda} x _{1} - J _{\lambda} x _{2}, f \right\rangle - \lambda^{-1} \left\langle x _{1} - x _{2}, f \right\rangle \leqq \lambda^{-1} \| J _{\lambda} x _{1} - J _{\lambda} x _{2} \| \cdot \| f \|
\\
- \lambda^{-1} \left\langle x _{1} - x _{2}, f \right\rangle \leqq \lambda^{-1} \| x _{1} - x _{2} \| ^{2} - \lambda^{-1} \| x _{1} - x _{2} \| ^{2} = 0. \end{gathered}
\]

This proves that $A_{\lambda}$ is dissipative. \eqref{eq:XIV.6.8} is clear. \eqref{eq:XIV.6.9} is proved as follows. By \eqref{eq:XIV.6.3} and \eqref{eq:XIV.6.6}, we obtain, for any $y \in Ax$, 
\[
\lambda \| A _{\lambda} x \| = \| J _{\lambda} x - x \| = \| J _{\lambda} x - J _{\lambda} (x - \lambda y) \| \leqq \| x - (x - \lambda y) \| = \lambda \| y \|.
\]
\end{proof}

\begin{lemma}[Y. \textsc{Kōmura} {\cite{komura158}}]\label{lem:XIV.6.Lemma.2} Let $A$ be a dissipative set $\subseteq X \times X$, and assume that $D(J_{\lambda}) = X$ for a $\lambda > 0$. Then $D(J_{\mu}) = X$ for every $\mu > 0$ with $0 \leqq |(\mu - \lambda)/\mu| < 1$. 
\end{lemma}

\begin{proof}
 Take any point $x \in X$, and consider a single-valued mapping $T$ defined by 
\[
X \ni z \rightarrow T z = J _{\lambda} \left(\frac {\lambda}{\mu} x + \frac {\mu - \lambda}{\mu} z\right). \]

Since $J_{\lambda}$ is a contraction mapping by \eqref{eq:XIV.6.6}, we obtain 
\[
\left\| T z - T w \right\| \leqq \left\| \left(\frac {\lambda}{\mu} x + \frac {\mu - \lambda}{\mu} z\right) - \left(\frac {\lambda}{\mu} x + \frac {\mu - \lambda}{\mu} w\right) \right\| = \left| \frac {\mu - \lambda}{\mu} \right| \cdot \| z - w \|.
\]
that is, $T$ is a Lipschitzian mapping with Lipschitz constant 
\[
\alpha = | (\mu - \lambda) / \mu | < 1.
\]

Thus, for $n > m$ and for any point $z \in X$, we have 
\[
\begin{aligned} & {\| T ^{n} z - T ^{m} z \| \leqq \alpha^{m} \| T ^{n-m} z - z \| \leqq \alpha^{m} (\| T z - z \| + \| T ^{2} z - T z \| + \dots)} \\ & {\quad \leqq \alpha^{m} (1 + \alpha + \alpha^{2} + \dots) \cdot \| T z - z \| \leqq \alpha^{m} (1 - \alpha) ^{-1} \| T z - z \|.} \end{aligned}
\]

Hence, by the completeness of the space $X$, $\slim\limits_{n\to 0}T^n z = y$ exists in $X$. $T$ being continuous as a Lipschitzian mapping, we obtain $T\cdot y = y$ by $y = \slim\limits_{n\to 0}T^{n+1}z = \slim\limits_{n\to 0}T(T^n z)$. Hence 
\[ y = J _{\lambda} \left(\frac {\lambda}{\mu} x + \frac {\mu - \lambda}{\mu} y\right) = J _{\lambda} \left(y - \lambda \left(\frac {1}{\mu} y - \frac {1}{\mu} x\right)\right),
\]
that is, $\left(\frac {1}{\mu} y - \frac {1}{\mu} x\right) \in A y$, and so $y - \mu z = x$ with $z \in Ay$. This proves that $J_{\mu}x = y$. Since $x$ was arbitrary, we must have $D(J_{\mu}) = X$. 
\end{proof}

\begin{corollaryu} Repeating the above argument, we can prove that $D(J_{\mu}) = X$ for all $\mu > 0$. 
\end{corollaryu}

We are now able to give 

\begin{definition}\label{def:XIV.6.Definition.4} $A$ dissipative set $A \subseteq X \times X$ is called a hyperdissipative set if $D(J_{\lambda}) = X$ for some $\lambda > 0$ and hence for all $\lambda > 0$. 
\end{definition}

\begin{proposition}\label{prop:XIV.6.Proposition.2} $A$ hyperdissipative set $A \subseteq X \times X$ must be maximally dissipative in the sense that there does not exist a dissipative set $B \subseteq X \times X$ which contains $A$ as a proper subset. 
\end{proposition}

\begin{proof}
 Let us assume that a dissipative set $B \subseteq X \times X$ contains $A$ as a subset. Let $y \in Bx$. Then, since $A$ is a hyperdissipative set, there exists a point $\{x_{1}, y_{1}\} \in A$ such that $x - y = x_{1} - y_{1}$. Since $A \subseteq B$, we must have $\{x_{1}, y_{1}\} \in B$ and hence we obtain $x = x_{1}$ and $y = y_{1}$ by \eqref{eq:XIV.6.5} applied to the dissipative set $B$. 
\end{proof}

We have thus prepared tools for the exposition of \textsc{Kōmura}'s approach to the integration in Hilbert space of non-linear evolution equations. Let $H$ be a real or complex Hilbert space with the scalar product $(x, z)$, and $A \subseteq H \times H$ be a hyperdissipative set. We shall be concerned with strong solutions of the initial value problem: 
\begin{equation}\label{eq:XIV.6.10}
\begin{yoslines}
&\frac {\dd u (t)}{\dd t} \in A u (t) \quad \text{ for almost every } t \text{ on the interval } [ 0, \infty), u (0) = x _{0} \in D (A),
\end{yoslines}
\end{equation}
where $u(t)$ defined on $[0, \infty)$ with values in $H$ is called a strong solution of \eqref{eq:XIV.6.10} if $u(t)$ is strongly absolutely continuous in $t$, $u(t)$ is strongly differentiable almost everywhere on $[0, \infty)$ such that the strong differential quotient $\dd u(t)/\dd t$ is Bochner integrable on any compact interval of $t$ and $u$ satisfies \eqref{eq:XIV.6.10}. We shall approximate this problem \eqref{eq:XIV.6.10} by 
\begin{equation}\tag{10'}\label{eq:XIV.6.10'}
\begin{yoslines}& \frac {\dd u ^{(\lambda)} (t)}{\dd t} - A _{\lambda} u ^{(\lambda)} (t) = 0 \quad \text{ for all } t \quad \text{ on } [ 0, \infty) \quad (\lambda > 0),\\& u ^{(\lambda)} (0) = x _{\lambda} = x _{0} - \lambda y _{0} \quad \text{ with a fixed } \quad y _{0} \in A x _{0}.\end{yoslines}
\end{equation}

Since $A_{\lambda}$ with $\lambda > 0$ is a Lipschitzian mapping, we see that equation \eqref{eq:XIV.6.10'} admits a uniquely determined solution $u^{(\lambda)}(t)$ which can be obtained by E. \textsc{Picard}'$S$ successive approximation: 
\begin{equation}\tag{10''}\label{eq:XIV.6.10''}
u _{n+1} ^{(\lambda)} (t) = x _{\lambda} + \int_{0} ^{t} A _{\lambda} u _{n} ^{(\lambda)} (s) \dd s \quad (n = 0, 1 \dots ; u _{0} ^{(\lambda)} (t) = x _{\lambda}).
\end{equation}

\begin{lemma}[Y. \textsc{Kōmura} {\cite{komura158}}]\label{lem:XIV.6.Lemma.3} For the solution $u^{(\lambda)}(t)$ of $\eqref{eq:XIV.6.10'}$, we have the following estimates: 
\begin{equation}\label{eq:XIV.6.11}
\left\| \frac {d}{\dd t} u ^{(\lambda)} (t) \right\| \leqq \left\| \frac {d}{\dd s} u ^{(\lambda)} (s) \right\| \quad \text{ for } 0 \leqq s < t,
\end{equation}
and 
\begin{equation}\label{eq:XIV.6.12}
\left\| u ^{(\lambda)} (t) - u ^{(\mu)} (t) \right\| ^{2} \leqq \left((\lambda - \mu) ^{2} + 4 t (\lambda + \mu)\right) \left\| y _{0} \right\| ^{2}.
\end{equation}
\end{lemma}

\begin{proof}
 Putting $u^{(\lambda)}(t + h) = v^{(\lambda)}(t)$ with $h > 0$, we obtain 
\[
\begin{aligned} \frac {d}{\dd t} \| v ^{(\lambda)} (t) - u ^{(\lambda)} (t) \| ^{2} & = 2 \operatorname{Re} \left(\frac {d}{\dd t} v ^{(\lambda)} (t) - \frac {d}{\dd t} u ^{(\lambda)} (t), v ^{(\lambda)} (t) - u ^{(\lambda)} (t)\right) \\ & = 2 \operatorname{Re} \left(A _{\lambda} v ^{(\lambda)} (t) - A _{\lambda} u ^{(\lambda)} (t), v ^{(\lambda)} (t) - u ^{(\lambda)} (t)\right) \leqq 0 \end{aligned}
\]
by the dissipative property of $A_{\lambda}$. Thus $\|u^{(\lambda)}(t + h) - u^{(\lambda)}(t)\|$ is monoton decreasing in $t$ and so we obtain \eqref{eq:XIV.6.11}. Therefore we have 
\begin{equation}\label{eq:XIV.6.13}
\left\| \frac {d}{\dd t} u ^{(\lambda)} (t) \right\| = \| A _{\lambda} u ^{(\lambda)} (t) \| \leqq \| A _{\lambda} u ^{(\lambda)} (0) \| = \| A _{\lambda} (x _{0} - \lambda y _{0}) \| = \| y _{0} \|. \end{equation}

By a similar argument as above, we obtain, by $\lambda A_{\lambda} = J_{\lambda} - I$, 
\begin{gather}\left\| u ^{(\lambda)} (t) - u ^{(\mu)} (t) \right\| ^{2} - \left\| u ^{(\lambda)} (0) - u ^{(\mu)} (0) \right\| ^{2} = \left\| u ^{(\lambda)} (t) - u ^{(\mu)} (t) \right\| ^{2} - \left\| (\lambda - \mu) y _{0} \right\| ^{2} \notag
\\= \int_{0} ^{t} \frac {d}{\dd s} \left\| u ^{(\lambda)} (s) - u ^{(\mu)} (s) \right\| ^{2} \dd s = 2 \int_{0} ^{t} \operatorname{Re} \left(A _{\lambda} u ^{(\lambda)} (s) - A _{\mu} u ^{(\mu)} (s), \right. \left. u ^{(\lambda)} (s) - u ^{(\mu)} (s)\right) \dd s \notag
\\= 2 \int_{0} ^{t} \operatorname{Re} \left(A _{\lambda} u ^{(\lambda)} (s) - A _{\mu} u ^{(\mu)} (s), \right. \left. J _{\lambda} u ^{(\lambda)} (s) - J _{\mu} u ^{(\mu)} (s)\right) \dd s \notag
\\\quad - 2 \int_{0} ^{t} \operatorname{Re} \left(A _{\lambda} u ^{(\lambda)} (s) - A _{\mu} u ^{(\mu)} (s), \right. \left. \lambda A _{\lambda} u ^{(\lambda)} (s) - \mu A _{\mu} u ^{(\mu)} (s)\right) \dd s. \notag
\end{gather}

The first term on the extreme right is non-negative, because $A$ is a dissipative set in $H \times H$ and $A_{\lambda}u^{(\lambda)}(s) \in A J_{\lambda}u^{(\lambda)}(s)$ by \eqref{eq:XIV.6.8}. The second term on the extreme right is, by \eqref{eq:XIV.6.13} smaller than $4t(\lambda + \mu)\|y_{0}\|^{2}$ so that 
\[
\left\| u ^{(\lambda)} (t) - u ^{(\mu)} (t) \right\| ^{2} - (\mu - \lambda) ^{2} \left\| y _{0} \right\| ^{2} \leqq 4 t (\lambda + \mu) \left\| y _{0} \right\| ^{2}.
\]

This proves \eqref{eq:XIV.6.12}. 
\end{proof}

\begin{corollaryu} 
\begin{equation}
\label{eq:XIV.6.14}
\begin{minipage}{0.92\linewidth}
\everymath{}\let\limits\nolimits
$\slim_{\lambda \downarrow 0} u^{(\lambda)}(t) = u(t)$ exists uniformly on every compact set of $t.$
\end{minipage}
\end{equation}
\end{corollaryu}

\begin{equation}
\label{eq:XIV.6.15}
\begin{minipage}{0.92\linewidth}
\everymath{}\let\limits\nolimits
$\slim _{\lambda \downarrow 0} J _{\lambda} u ^{(\lambda)} (t) = u (t)$ uniformly on every compact set of $t.$
\end{minipage}
\end{equation}

\begin{proof}
 ($14$) is clear from \eqref{eq:XIV.6.12}. We also have \eqref{eq:XIV.6.15} from ($14$), since we have 
\[
\left\| J _{\lambda} u ^{(\lambda)} (t) - u ^{(\lambda)} (t) \right\| = \lambda \left\| A _{\lambda} u ^{(\lambda)} (t) \right\| \leqq \lambda \left\| y _{0} \right\| \text{ by } (1 3). \]
\end{proof}

\begin{lemma}[Y. \textsc{Kōmura} {\cite{komura158}}]\label{lem:XIV.6.Lemma.4} $u(t) \in D(A)$ for all $t \geqq 0$. 
\end{lemma}

\begin{proof}
 By \eqref{eq:XIV.6.8} we have, for fixed $t > 0$, 
\begin{equation}\label{eq:XIV.6.16}
\left\{J _{\lambda} u ^{(\lambda)} (t), A _{\lambda} u ^{(\lambda)} (t) \right\} \in A \quad \text{ for every } \quad \lambda > 0. \end{equation}

Thus, setting $A_{\lambda}u^{(\lambda)}(t) = w^{(\lambda)}(t)$, we obtain $\|w^{(\lambda)}(t)\| \leqq \|y_0\|$ by \eqref{eq:XIV.6.13}. Hence, by the local weak compactness of the Hilbert space $H$, there exists, for fixed $t > 0$, a sequence $\{\lambda_n\}$ of positive numbers satisfying $\lambda_n \downarrow 0$ and such that 
\begin{equation}\label{eq:XIV.6.17}
\wlim _{n \to \infty} w ^ {(\lambda_{n})} (t) = w (t) \quad H \quad \text{ with } \quad \| w (t) \| \leqq \| y _{0} \|.
\end{equation}

Since $A$ is a dissipative set, we have, for every point $\{x, y\} \in A$, the inequality $\operatorname{Re}(y - w^{(\lambda_n)}(t), x - J_{\lambda_n}u^{(\lambda_n)}(t)) \leqq 0$. Hence, by letting $n \uparrow \infty$, we obtain $\operatorname{Re}(y - w(t), x - u(t)) \leqq 0$ by \eqref{eq:XIV.6.15}. This proves that 
\begin{equation}\label{eq:XIV.6.18}
\{u (t), w (t) \} \in A,
\end{equation}
because $A$ is a maximally dissipative set $\subseteq H \times H$ as proved in Proposition~\ref{prop:XIV.6.Proposition.2}. 
\end{proof}

\begin{lemma}[Y. \textsc{Kōmura} {\cite{komura158}}]\label{lem:XIV.6.Lemma.5} $u(t)$ is strongly absolutely continuous in $t$, and it is strongly differentiable at almost every $t \geqq 0$. 
\end{lemma}

\begin{proof}
 By \eqref{eq:XIV.6.13}, we obtain, for $0 \leqq t_{1} < t_{2} < \cdots$ with $\sum\limits_{i} (t_{i+1} - t_{i}) < \infty$, 
\[
\left\| u ^{\lambda} \left(t _{i+1}\right) - u ^{(\lambda)} \left(t _{i}\right) \right\| \leqq \int_ {t _{i}} ^ {t _{i+1}} \left\| \frac {d}{\dd s} u ^{(\lambda)} (s) \right\| \dd s \leqq \left(t _{i+1} - t _{i}\right) \cdot \left\| y _{0} \right\|. \]

Therefore, by letting $\lambda \downarrow 0$, we obtain 
\begin{equation}\label{eq:XIV.6.19}
\sum\limits_{i} \left\| u \left(t _{i+1}\right) - u \left(t _{i}\right) \right\| \leqq \sum\limits_{i} \left(t _{i+1} - t _{i}\right) \cdot \left\| y _{0} \right\|.
\end{equation}

This proves the strong absolute continuity in $t$ of $u(t)$. 

Next let $0 < t_0 < \infty$. Then, by \eqref{eq:XIV.6.19}, the set $\{u(t); 0 \leqq t \leqq t_0\}$ is compact in $H$, hence it is separable. Therefore, without losing the generality, we may assume that $H$ is separable. Thus let $\{x_k\}$ be a strongly dense countable sequence of $H$. Each numerical function $v_k(t) = (u(t), x_k)$ is absolutely continuous by \eqref{eq:XIV.6.19}, and so there exists a set $N_k$ of measure zero such that $v_k(t)$ is differentiable on $[0, t_0] - N_k$. In view of \eqref{eq:XIV.6.19}, we see that $u(t)$ is weakly differentiable on $[0, t_0] - \bigcup\limits_{k=1}^{\infty} N_k$, and the weak derivative $u'(t)$ satisfies $\|u'(t)\| \leqq \|y_0\|$. $H$ being assumed separable, the weakly measurable function $u'(t)$ is strongly measurable by \textsc{Pettis}' theorem (Chapter~\ref{chap:V}, \ref{sec:V.4}) and so, by the boundedness condition $\|u'(t)\| \leqq \|y_0\|$, $u'(t)$ is Bochner integrable. Therefore, by \textsc{Bochner}'$S$ theorem (Chapter~\ref{chap:V}, \ref{sec:V.5}), 
\begin{equation}\label{eq:XIV.6.20}
u (t) - u (0) = \int_{0} ^{t} u ^{\prime} (s) \dd s
\end{equation}
is strongly differentiable for almost every $t \in [0, t_0]$ with $u'(t)$ as the strong derivative. 
\end{proof}
We are now able to prove Y. \textsc{Kōmura}'s 

\begin{theoremu} $u(t)$ is a strong solution of the initial value problem \eqref{eq:XIV.6.10}. 
\end{theoremu}

\begin{proof}
 For a fixed positive number $t_0$, we define the space $L^2([0, t_0], H)$ of $H$-valued strongly measurable function $\tilde{x} = x(t)$ such that $\int_{0}^{t_0} \|x(s)\|^2 \dd s < \infty$, and normed by $\|\tilde{x}\|^2 = \int_{0}^{t_0} \|x(s)\|^2 \dd s$. It is easy to show that $L^2([0, t_0], H)$ is a Hilbert space with the scalar product 
\begin{equation}\label{eq:XIV.6.21}
(\tilde {x}, \tilde {y}) = \int_{0} ^ {t _{0}} (x (s), y (s)) \dd s,
\end{equation}
where $(x(s), y(s))$ is the scalar product in $H$ of $x(s)$ and $y(s)$. 

Thus our $\tilde{u} = u(t)$ can be considered as an element of $L^2([0, t_0], H)$, because $u(t)$ is a strongly continuous $H$-valued function. Moreover, the hyperdissipative set $A$ can naturally be extended to a hyperdissipative set $\tilde{A} \subseteq L^2([0, t_0], H) \times L^2(0, t_0], H$ by defining 
\begin{equation}\label{eq:XIV.6.22}
\tilde {A} \tilde {x} = \{\tilde {y} \in L ^{2} ([ 0, t _{0} ], H); y (t) \in A x (t) \text{ for almost every } t \in [ 0, t _{0} ] \}
\end{equation}

We have, by \eqref{eq:XIV.6.16} and \eqref{eq:XIV.6.13}, 
\begin{equation}\tag{16'}\label{eq:XIV.6.16'}
\{\widetilde {J} _{\lambda} \widetilde {u} ^{(\lambda)}, \widetilde {A} _{\lambda} \widetilde {u} ^{(\lambda)} \} \in \tilde {A} \quad \text{ for every } \lambda > 0.
\end{equation}

Again by \eqref{eq:XIV.6.13}, $\{\widetilde{A}_{\lambda}\widetilde{u}^{(\lambda)};\lambda >0\}$ is a norm bounded subset of the Hilbert space $L^2 ([0,t_0],H)$. Hence there exists a sequence $\{\lambda_n\}$ of positive numbers satisfying $\lambda_{n}\downarrow 0$ and such that 
\begin{equation}\tag{17'}\label{eq:XIV.6.17'}
\wlim \widetilde {A} _ {\lambda_{n}} \widetilde {u} ^ {(\lambda_{n})} = \tilde {w} \in L ^{2} ([ 0, t _{0} ], H) .
\end{equation}

Moreover, we have, by \eqref{eq:XIV.6.15}, 
\begin{equation}\tag{15'}\label{eq:XIV.6.15'}
\slim _{n \rightarrow \infty} \widetilde {J} _ {\lambda_{n}} \widetilde {u} ^ {(\lambda_{n})} = \tilde {u} \in L ^{2} ([ 0, t _{0} ], H).
\end{equation}
$\tilde{A}$ is a hyperdissipative set by the fact that $A$ is a hyperdissipative set. Thus, by Proposition~\ref{prop:XIV.6.Proposition.2}, $\tilde{A}$ is a maximally dissipative set. Therefore, we have, as in the proof of Lemma~\ref{lem:XIV.6.Lemma.4}, 
\begin{equation}\tag{18'}\label{eq:XIV.6.18'}
w (t) \in A u (t) \quad \text{ for almost every } \quad t \in [ 0, t _{0} ] .
\end{equation}

On the other hand, we have 
\[
u ^ {(\lambda_{n})} (t) - u ^ {(\lambda_{n})} (0) = \int_{0} ^{t} \frac {d}{\dd s} u ^ {(\lambda_{n})} (s) \dd s \quad (\text{ for } 0 \leqq t \leqq t _{0}),
\]
and so, for every $x \in H$, 
\[
(u ^ {(\lambda_{n})} (t), x) - \langle x _ {\lambda_{n}}, x \rangle = \int_{0} ^{t} \left(\frac {d}{\dd s} u ^ {(\lambda_{n})} (s), x\right) \dd s \quad (\text{ for } \quad 0 \leqq t \leqq t _{0})
\]
by Colollary $2$ in Chapter~\ref{chap:V}, \ref{sec:V.5}. Hence, we obtain, by ($14$) and \eqref{eq:XIV.6.17'}, 
\begin{equation}\label{eq:XIV.6.23}
(u (t), x) - (x _{0}, x) = \int_{0} ^{t} (w (s), x) \dd s = (\int_{0} ^{t} w (s) \dd s, x).
\end{equation}

Since $x \in H$ and $t_0 > 0$ were arbitrary, we obtain, by \eqref{eq:XIV.6.23}, \eqref{eq:XIV.6.18'} and \eqref{eq:XIV.6.20}, that $\frac{\dd u(t)}{\dd t} = u'(t) = w(t) \in Au(t)$ for almost every $t \geqq 0$. 
\end{proof}

\begin{remarku} T. \textsc{Kato} \cite{kato144} extended the above Theorem to the case where $X$ is a Banach space whose dual space is uniformly convex. An almost complete version of the Hille-Yosida theorem for non-linear contraction semi-groups was devised and proved in Hilbert space by Y. \textsc{Kōmura} \cite{komura159}. Cf. M. G. \textsc{Crandall}-A. \textsc{Pazy} \cite{clarkson36}, T. \textsc{Kato} \cite{kato145} and J. R. \textsc{Dorroh} \cite{doob43} cited in § $5$ of Y. \textsc{Kōmura} \cite{komura159}. In this paper of \textsc{Kōmura}, the most crucial point is the proof that the infinitesimal generator is densely defined, and this proof was later fairly simplified by T. \textsc{Kato} \cite{kato145}. Recently, H. \textsc{Brezis} has published a comprehensive treatise on nonlinear semi-groups in Hilbert space. This book (H. \textsc{Brezis} \cite{brezis28}) contains an extensive bibliography which is not necessarily restricted to the Hilbert space. 
\end{remarku}


\yossection[7. Non-linear Evolution Equations 2]{7}{Non-linear Evolution Equations 2 (The Approach through the Crandall-Liggett Convergence Theorem)}
\label{sec:XIV.7}

Let $X$ be a real or complex Banach space, and $A$ the infinitesimal generator of an equi-continuous semi-group $\{T_t; t \geqq 0\}$ of linear operators $\in L(X, X)$ of class $(C_0)$. Then the Lemma in Chapter~\ref{chap:IX}, \ref{sec:IX.12} gives the proof of $E$. Hille's approximation of $T_t$: 

\begin{equation}
\label{eq:XIV.7.1}
\begin{minipage}{0.92\linewidth}
\everymath{}\let\limits\nolimits
For every $x_0 \in X,  T_t x_0 = \slim\limits_{n \to \infty} (I - n^{-1} t A)^{-n} \cdot x_0,$ and the convergence is uniform on every compact interval of $t.$
\end{minipage}
\end{equation}

The titled convergence theorem is suggested by \eqref{eq:XIV.7.1}. It reads as follows. 

\begin{theorem}[M. \textsc{Crandall}-T. \textsc{Liggett} {\cite{clarkson37}}]\label{thm:XIV.7.Theorem.1} Let $A$ be a hyperdissipative set $\subseteq X \times X$. Then, for all $x_0 \in D(A)$, 

\begin{equation}
\tag{1'}\label{eq:XIV.7.1'}
\begin{minipage}{0.92\linewidth}
\everymath{}\let\limits\nolimits
$\slim\limits_{n\to \infty}(J_{t/n})^n \cdot x_0$ exists uniformly on every compact interval of $t,$ where $J_{\lambda} = (I - \lambda A)^{-1}$ as in the preceding Section~\hyperref[sec:XIV.6]{6}.
\end{minipage}
\end{equation}
\end{theorem}

The following proof is adapted from S. \textsc{Rasmussen} \cite{rasmussen269} which seems to give a good modification of the original proof by \textsc{Crandall}-\textsc{Liggett}: 

\begin{proof}
 The first step. We shall prepare the following ($2$)-($4$): 

\begin{equation}\label{eq:XIV.7.2}
\|A_{\lambda}x\| \leqq \|Ax\| \text{ for all } x \in D(A), \text{ and for all } \lambda > 0,
\end{equation}

\begin{equation}\label{eq:XIV.7.3}
\left\| \prod\limits_{i=1} ^{n} J _ {\lambda_{i}} x - x \right\| = \| J _ {\lambda_{n}} J _ {\lambda_{n-1}} \dots J _ {\lambda_{1}} x - x \| \leqq \left(\sum\limits_{i=1} ^{n} \lambda_{i}\right) \| A x \|
\end{equation}
for all $x \in D(A)$ and for all $\lambda_i > 0$ ( $i = 1, 2, \ldots, \overline{n}$). 

\begin{equation}\label{eq:XIV.7.4}
J_{\mu}\left(\frac{\mu}{\lambda} x + \frac{\lambda - \mu}{\lambda} J_{\lambda}x\right) = J_{\lambda}x \text{ for all } x \in X \text{ and for all } \lambda, \mu > 0.
\end{equation}


($2$) was proved in the preceding Section~\hyperref[sec:XIV.6]{$6$}. The proof of ($3$) is obtained by the contraction property of $J_{\lambda}$, ($2$) and $\lambda A_{\lambda} = J_{\lambda} - I$ as follows. 
\[
\begin{aligned}
\left\| \prod\limits_{i=1} ^{n} J _ {\lambda_{i}} x - x \right\|& = \left\| \sum\limits_{i=1} ^{n} \left(\prod\limits_{j=i} ^{n} J _ {\lambda_{j}} x - \prod\limits_{j=i+1} ^{n} J _ {\lambda_{j}} x\right) \right\| \leqq \sum\limits_{i=1} ^{n} \left\| \prod\limits_{j=i} ^{n} J _ {\lambda_{j}} x - \prod\limits_{j=i+1} ^{n} J _ {\lambda_{j}} x \right\|
\\
&\leqq \sum\limits_{i=1} ^{n} \| J _ {\lambda_{i}} x - x \| \leqq \sum\limits_{i=1} ^{n} \| \lambda_{i} A _ {\lambda_{i}} x \| \leqq \left(\sum\limits_{i=1} ^{n} \lambda_{i}\right) \| A x \|.
\end{aligned}
\]

We next prove ($4$). Since $A$ is a hyperdissipative set, there exists, for any $x \in X$ and $\lambda > 0$, a point $\{x_1, y_1\} \in A$ such that $x = x_1 - \lambda y_1$. Thus $J_\lambda x = x_1$ and so 
\[
\frac {\mu}{\lambda} x + \frac {\lambda - \mu}{\lambda} J _{\lambda} x = \frac {\mu}{\lambda} (x _{1} - \lambda y _{1}) + \frac {\lambda - \mu}{\lambda} x _{1} = x _{1} - \mu y _{1},
\]
proving that $J_{\mu}(x_1 - \mu y_1) = x_1 = J_{\lambda}x$. 

The second step. Let $x \in D(A)$ and $\lambda > 0$. Let furthermore $\{\mu_n; n \geqq 1\}$ be a sequence such that $0 < \mu_n \leqq \lambda$ for all $n$. We then define 
\begin{align}
A _{n,m}& = \left\| \prod\limits_{i=j} ^{n} J _ {\mu_{i}} x - J _{\lambda} ^{m} x \right\| \text{ for } n, m = 0, 1, 2, \dots \label{eq:XIV.7.5} \\
&\quad \left(\text{ We set } \prod\limits_{i=1} ^{0} J _ {\mu_{i}} = I = J _{\lambda} ^{0} \text{ for convenience. }\right) \notag
\end{align}

Moreover, let 
\begin{equation}\label{eq:XIV.7.6}
t _{n} = \sum\limits_{i=1} ^{n} \mu_{i} (n = 0, 1, 2, \dots ; t _{0} = 0)
\end{equation}
and 
\begin{equation}\label{eq:XIV.7.7}
\alpha_{n} = \mu_{n} / \lambda \quad \text{ and } \quad \beta_{n} = 1 - \alpha_{n} = (\lambda - \mu_{n}) / \lambda \quad \text{ for } \quad n = 1, 2, \dots
\end{equation}

Then we have 
\begin{equation}\label{eq:XIV.7.8}
A _{n,m} \leqq \left\{\left[ (m \lambda - t _{n}) ^{2} + m \lambda^{2} \right] ^{1/2} + \left[ (m \lambda - t _{n}) ^{2} + \lambda t _{n} \right] ^{1/2} \right\} \| A x \|.
\end{equation}

For the proof, we first show 
\begin{equation}\label{eq:XIV.7.9}
A _{n,m} \leqq \alpha_{n} A _{n-1,m-1} + \beta_{n} A _{n-1,m} \quad \text{ for } \quad n, m = 1, 2, \dots .
\end{equation}

In fact we have, by ($4$) and the contraction property of $J_{\mu_{n}}$, 
\[
\begin{aligned}
&A _{n,m} = \left\| \prod\limits_{i=1} ^{n} J _ {\mu_{i}} x - J _{\lambda} ^{m} x \right\| = \left\| \prod\limits_{i=1} ^{n} J _ {\mu_{i}} x - J _ {\mu_{n}} \left(\frac {\mu_{n}}{\lambda} J _{\lambda} ^{m-1} x + \frac {\lambda - \mu_{n}}{\lambda} J _{\lambda} ^{m} x\right) \right\|
\\
&\leqq \left\| \prod\limits_{i=1} ^{n-1} J _ {\mu_{i}} x - \left(\frac {\mu_{n}}{\lambda} J _{\lambda} ^{m-1} x + \frac {\lambda - \mu_{n}}{\lambda} J _{\lambda} ^{m} x\right) \right\| \leqq \alpha_{n} \left\| \prod\limits_{i=1} ^{n-1} J _ {\mu_{i}} x - J _{\lambda} ^{m-1} x \right\| + \beta_{n} \left\| \prod\limits_{i=1} ^{n-1} J _ {\mu_{i}} x - J _{\lambda} ^{m} x \right\|
\\
&= \alpha_{n} A _{n -1 \cdot m -1} + \beta_{n} A _{n-1,m}.
\end{aligned}
\]

We shall prove \eqref{eq:XIV.7.8} by \eqref{eq:XIV.7.9} and induction. We start with the proof that $A_{0,m}$ satisfies \eqref{eq:XIV.7.8} for all $m = 0, 1, \ldots$. In fact, we have 
\[
\begin{aligned} A _{0,m} & = \| x - J _{\lambda} ^{m} x \| \leqq m \lambda \| A x \| \quad \text{ (by \eqref{eq:XIV.7.3} and \eqref{eq:XIV.7.5}) } \\ & \leqq \left\{\left[ (m \lambda - t _{0}) ^{2} + m \lambda^{2} \right] ^{1/2} + \left[ \left(m \lambda - t _{0}\right) ^{2} + \lambda t _{0} \right] ^{1/2} \right\} \| A x \| (\text{since $t_{0}$ = 0}) \end{aligned}
\]

Let now $n$ and $m$ be arbitrary and assume that $A_{n,m}$ and $A_{n,m-1}$ both satisfy \eqref{eq:XIV.7.8}. We shall then show that $A_{n+1,m}$ satisfies \eqref{eq:XIV.7.8}. In fact, we have 
\[
\begin{aligned}
A_{n+1,m}&\leqq\alpha_{n+1}A_{n,m-1}+\beta_{n+1}A_{n,m}
  \quad\text{(by \eqref{eq:XIV.7.9})}\\
&\leqq\alpha_{n+1}\bigl\{[((m-1)\lambda-t_n)^2+(m-1)\lambda^2]^{1/2}\\
&\qquad+[((m-1)\lambda-t_n)^2+\lambda t_n]^{1/2}\bigr\}\|Ax\|\\
&\quad+\beta_{n+1}\bigl\{[(m\lambda-t_n)^2+m\lambda^2]^{1/2}
  +[(m\lambda-t_n)^2+\lambda t_n]^{1/2}\bigr\}\|Ax\|\\
&\leqq(\alpha_{n+1}+\beta_{n+1})^{1/2}
  \bigl\{\alpha_{n+1}[((m-1)\lambda-t_n)^2+(m-1)\lambda^2]\\
&\qquad+\beta_{n+1}[(m\lambda-t_n)^2+m\lambda^2]\bigr\}^{1/2}\|Ax\|\\
&\quad+(\alpha_{n+1}+\beta_{n+1})^{1/2}
  \bigl\{\alpha_{n+1}[((m-1)\lambda-t_n)^2+\lambda t_n]\\
&\qquad+\beta_{n+1}[(m\lambda-t_n)^2+\lambda t_n]\bigr\}^{1/2}\|Ax\|
  \quad\text{(by Schwarz' inequality)}\\
&=\bigl\{m^2\lambda^2-2\alpha_{n+1}m\lambda^2+\alpha_{n+1}\lambda^2
  -2m\lambda t_n+2\alpha_{n+1}\lambda t_n+t_n^2\\
&\qquad+m\lambda^2-\alpha_{n+1}\lambda^2\bigr\}^{1/2}\|Ax\|\\
&\quad+\bigl\{m^2\lambda^2-2\alpha_{n+1}m\lambda^2+\alpha_{n+1}\lambda^2
  -2m\lambda t_n+2\alpha_{n+1}\lambda t_n+t_n^2+\lambda t_n\bigr\}^{1/2}\|Ax\|\\
&\hspace{7em}\text{(by $\beta_{n-1}=1-\alpha_{n-1}$)}\\
&=\bigl\{m^2\lambda^2-2\mu_{n+1}m\lambda+\mu_{n+1}\lambda
  -2m\lambda t_n+2\mu_{n+1}t_n+t_n^2+m\lambda^2-\mu_{n+1}\lambda\bigr\}^{1/2}\|Ax\|\\
&\quad+\bigl\{m^2\lambda^2-2\mu_{n+1}m\lambda+\mu_{n+1}\lambda
  -2m\lambda t_n+2\mu_{n+1}t_n+t_n^2+\lambda t_n\bigr\}^{1/2}\|Ax\|\\
&=\bigl\{(m\lambda-(\mu_{n+1}+t_n))^2-\mu_{n+1}^2+m\lambda^2\bigr\}^{1/2}\|Ax\|\\
&\quad+\bigl\{(m\lambda-(\mu_{n+1}+t_n))^2-\mu_{n+1}^2
  +\lambda(t_n+\mu_{n+1})\bigr\}^{1/2}\|Ax\|\\
&\leqq\bigl\{[(m\lambda-t_{n+1})^2+m\lambda^2]^{1/2}
  +[(m\lambda-t_{n+1})^2+\lambda t_{n+1}]^{1/2}\bigr\}\|Ax\|,
\end{aligned}
\]
and we have thus proved that $A_{n+1,m}$ satisfies \eqref{eq:XIV.7.8}. We can also prove that $A_{n,0}$ satisfies \eqref{eq:XIV.7.8} for all $n = 1, 2, \ldots$. In fact, we obtain by \eqref{eq:XIV.7.3} 
\[
\begin{aligned}
A _{n,0}& = \left\| \prod\limits_{i=1} ^{n} J _ {u _{i}} x - x \right\| \leqq \sum\limits_{i=1} ^{n} \mu_{i} \cdot \| A x \| = t _{n} \| A x \|
\\
\quad& \leqq \{[ (0 \cdot \lambda - t _{n}) ^{2} + 0 \cdot \lambda^{2} ] ^{1/2} + [ (0 \cdot \lambda - t _{n}) ^{2} + \lambda t _{n} ] ^{1/2} \} \| A x \|.
\end{aligned}
\]

Therefore, we have completed the induction and hence \eqref{eq:XIV.7.8} is proved. 

The third step. Let $t$ be a positive number, and consider a partition $\Delta$ of the closed interval $[0, t]$. 
\[
\Delta : 0 = t _{0} < t _{1} < \dots < t _{i-1} < t _{i} < t _{i+1} < \dots < t _{n} = t.
\]

For this partition, we define 
\[
\Delta_{i} = \left(t _{i} - t _{i-1}\right) \quad (i = 1, 2, \dots , n) \quad \text{ and } \quad | \Delta | = \max\limits _{1 \leqq i \leqq n} \left(t _{i} - t _{i-1}\right).
\]

Take another partition 
\[
\Delta^{\prime}: 0 = t _{0} ^{\prime} < t _{1} ^{\prime} < \dots < t _{j-1} ^{\prime} < t _{j} ^{\prime} < t _{j+1} ^{\prime} < \dots < t _{k} ^{\prime} = t
\]
and define 
\[
\Delta_{j} ^{\prime} = \left(t _{j} ^{\prime} - t _{j-1} ^{\prime}\right) \quad (j = 1, 2, \dots , k) \quad \text{ and } \quad | \Delta^{\prime} | = \max\limits _{1 \leqq j \leqq k} \left(t _{j} ^{\prime} - t _{j-1} ^{\prime}\right)
\]

Then, for $x_0 \in D(A)$, we obtain, by taking 
$\lambda = \max (|\Delta |, |\Delta '|)$ and $m =$ the largest integer $\leqq t / \lambda$, 
the following inequality implied from \eqref{eq:XIV.7.8}: 
\begin{equation}\label{eq:XIV.7.10}
\begin{yoslines}& \left\| \prod\limits_{i=1} ^{n} J _ {\Delta_{i}} x _{0} - \prod\limits_{j=1} ^{k} J _ {\Delta_ {j ^{\prime}}} x _{0} \right\| \leqq \left\| \prod\limits_{i=1} ^{n} J _ {\Delta_{i}} x _{0} - J _{\lambda} ^{m} x _{0} \right\| + \left\| J _{\lambda} ^{m} x _{0} - \prod\limits_{j=1} ^{k} J _ {\Delta_ {j ^{\prime}}} x _{0} \right\|\\& \leqq 2 \left\{\left[ (m \lambda - t) ^{2} + m \lambda^{2} \right] ^{1/2} + \left[ (m \lambda - t) ^{2} + \lambda t \right] ^{1/2} \right\} \| A x _{0} \|\\& \leqq 2 \left\{\left[ \lambda^{2} + t \lambda \right] ^{1/2} + \left[ \lambda^{2} + \lambda t \right] ^{1/2} \right\} \| A x \|.\end{yoslines}
\end{equation}

Therefore, by taking $\Delta_{i} = t / n$ ( $i = 1,2,\ldots,n$) and $\Delta j = t / k$ ( $j = 1,2,\ldots,k$) in \eqref{eq:XIV.7.10}, we obtain \eqref{eq:XIV.7.1'}. 
\end{proof}

\begin{remarku} We have incidentally proved that 
\begin{equation}
\tag{1''}\label{eq:XIV.7.1''}
\begin{minipage}{0.92\linewidth}
\everymath{}\let\limits\nolimits
$T _{t} x _{0} = \slim _{n \rightarrow \infty} (I - n ^{-1} t A) ^{-n} x _{0} = \slim (I - n ^{-1} A) ^{-[nt]} x _{0}$ uniformly on every compact interval of $t,$ where $[ n t ] =$ the largest integer $\leqq n t.$
\end{minipage}
\end{equation}
\end{remarku}

The scope of this Theorem will be seen from Theorem~\ref{thm:XIV.7.Theorem.2}, Theorem~\ref{thm:XIV.7.Theorem.3} and Theorem~\ref{thm:XIV.7.Theorem.4} given below. 

\begin{theorem}\label{thm:XIV.7.Theorem.2} Let $A$ be a hyperdissipative closed set of $X \times X$. Then, for every $x_{0} \in D(A)$, the following two conditions are mutually equivalent.
\begin{enumerate}[i)]
\item $u(t)$ is a strong solution of the initial value problem 
\begin{equation}\label{eq:XIV.7.11}
\frac {\dd u (t)}{\dd t} \in A u (t) \quad \text{ for almost every } \quad t \geqq 0 \quad \text{ and } \quad u (0) = x _{0},
\end{equation}
as formulated in the preceding Section~\hyperref[sec:XIV.6]{$6$}. 

\item 
\begin{equation}
\label{eq:XIV.7.12}
\begin{minipage}{0.92\linewidth}
\everymath{}\let\limits\nolimits
$u (t) = \slim _{n \rightarrow \infty} (I - n ^{-1} A) ^{-[nt]} x _{0}$ and $u (t)$ is strongly differentiable at almost every $t \geqq 0.$
\end{minipage}
\end{equation}
\end{enumerate}
\end{theorem}

We first give the \textsc{Brezis}-\textsc{Pazy} \cite{brezis29} proof of the assertion i) $\to$ ii), and the proof of the assertion ii) $\to$ i) due to \textsc{Crandall}-\textsc{Liggett} \cite{clarkson37} will be given later; after the proof of Theorem~\ref{thm:XIV.7.Theorem.3}. We shall begin with a key Lemma due to T. \textsc{Kato} \cite{kato144}. 

\begin{lemmau} Let $S$ be the set of those $s \geqq 0$ at which the strong solution $u(s)$ of \eqref{eq:XIV.7.11} is strongly differentiable. Then, we have, at almost every $s \in S$, 
\begin{equation}\label{eq:XIV.7.13}
2 ^{-1} \frac {d}{\dd s} \| u (s) \| ^{2} = \| u (s) \| \frac {\dd \| u (s) \|}{\dd s} = \operatorname{Re} \left\langle \frac {\dd u (s)}{\dd s}, f \right\rangle \text{ whenever } f \in F (u (s)).
\end{equation}
\end{lemmau}

\begin{proof}
 $u(s)$ being a strong solution of \eqref{eq:XIV.7.11}, the strong derivative $\dd u(s)/\dd s$ is Bochner integrable on every compact interval of $s$ so that 
\[
u (s) - u (0) = \int_{0} ^{s} \frac {\dd u (s)}{\dd s} \dd s
\]
by \textsc{Bochner}'s theorem in Chapter~\ref{chap:V}, \ref{sec:V.5}. Hence $\|u(s)\|$ is of bounded variation on every compact interval of $s$. Therefore, $\|u(s)\|$ is differentiable at almost every $s \geqq 0$. 

Since $\operatorname{Re}\langle u(t), f \rangle \leqq \|u(t)\| \cdot \|u(s)\|$ and $\operatorname{Re}\langle u(s), f \rangle = \|u(s)\|^2$ by $f \in F(u(s))$, we obtain 
\[
\operatorname{Re} \left\langle u (t) - u (s), f \right\rangle \leqq \| u (s) \| \left(\| u (t) \| - \| u (s) \|\right).
\]

Dividing the both side by $(t - s)$ and letting $t \rightarrow s$ from above and from below, we obtain at almost every $s \in S$, 
\[
\operatorname{Re} \left\langle \frac {\dd u (s)}{\dd s}, f \right\rangle \leqq \| u (s) \| \frac {\dd \| u (s) \|}{\dd s} \text{ and } \operatorname{Re} \left\langle \frac {\dd u (s)}{\dd s}, f \right\rangle \geqq \| u (s) \| \frac {\dd \| u (s) \|}{\dd s}.
\]

We have thus proved \eqref{eq:XIV.7.13}. 
\end{proof}

\begin{corollaryn}[H. \textsc{Brezis}-A. \textsc{Pazy} \cite{brezis28}] For the strong solution $u(t)$ of \eqref{eq:XIV.7.11}, we have
\begin{equation}\label{eq:XIV.7.14}
\left\| \frac {\dd u (t)}{\dd t} \right\| \leqq \| A x _{0} \| = \inf\limits _ {x \in A x _{0}} \| z \| \quad \text{ at almost every } t \geqq 0.
\end{equation}

\end{corollaryn}
\begin{proof}
 We put $\dd u(t) / \dd t = y(t)$ so that, by \eqref{eq:XIV.7.11}, $y(t) \in Au(t)$ at almost every $t \geqq 0$. For such $t$ and for every $z \in Ax_0$, there exists, by the dissipative property of $A$, an $f_0 \in F(u(t) - x_0)$ such that $\operatorname{Re}\langle y(t) - z, f_0 \rangle \leqq 0$. Hence, by the same argument which enabled us to prove \eqref{eq:XIV.7.13}, we obtain 
\[
2 ^{-1} \frac {d}{\dd t} \| u (t) - x _{0} \| ^{2} = \| u (t) - x _{0} \| \frac {d}{\dd t} \| u (t) - x _{0} \| = \operatorname{Re} \left\langle \frac {\dd u (t)}{\dd t} - 0, f _{0} \right\rangle
\]
$= \operatorname{Re}\langle y(t) - z,f_0\rangle +\operatorname{Re}\langle z,f_0\rangle \leqq \operatorname{Re}\langle z,f_0\rangle$ at almost every $t\geqq 0$ 

Thus, at almost every $t \geqq 0$, we have 
\[
\left\| u (t) - x _{0} \right\| \frac {d}{\dd t} \left\| u (t) - x _{0} \right\| \leqq | \langle z, f _{0} \rangle | \leqq \| z \| \cdot \| u (t) - x _{0} \|
\]
and so, since $z \in Ax_0$ was arbitrary, 
\begin{equation}\label{eq:XIV.7.15}
\left\| u (t) - x _{0} \right\| \leqq t \left\| A x _{0} \right\| \text{ at every } t \geqq 0.
\end{equation}

Next, let $h > 0$ and put $u(t + h) = v(t)$. Then 
$\frac{\dd v(t)}{\dd t} \in A v(t)$ for almost every $t \geqq 0$ and $v(0) = u(h)$. \yosparagraphtag{11'}\label{eq:XIV.7.11'} 

Thus, similarly as above, we obtain, for a certain $f \in F(v(t) - u(t))$, 
\[
2 ^{-1} \frac {d}{\dd t} \| v (t) - u (t) \| ^{2} = \operatorname{Re} \left\langle \frac {\dd v (t)}{\dd t} - \frac {\dd u (t)}{\dd t}, f \right\rangle \leqq 0 \text{ at almost every } t \geqq 0.
\]

Therefore, we have 
\[
\left\| u (t + h) - u (t) \right\| \leqq \left\| u (s + h) - u (s) \right\| \text{ whenever } t \geqq s \geqq 0,
\]
and so, combined with \eqref{eq:XIV.7.15}, we finally obtain 
\begin{equation}\tag{15'}\label{eq:XIV.7.15'}
\left\| u (t + h) - u (t) \right\| \leqq \left\| u (h) - u (0) \right\| \leqq h \left\| A x _{0} \right\|.
\end{equation}

This proves \eqref{eq:XIV.7.14}. 
\end{proof}

\begin{proof}[Proof of the assertion i) $\to$ ii) in Theorem~\ref{thm:XIV.7.Theorem.2}.]
The existence of the right hand limit of \eqref{eq:XIV.7.12} is guaranteed by Theorem~\ref{thm:XIV.7.Theorem.1}. Let $T$ be an arbitrary positive constant. Besides step functions $u_{n}(t) = (I - n^{-1}A)^{-[nt]}x_{0}$, we define on [$0$, $T$] another sequence of functions $v_{n}(t)$ by
\begin{equation}\label{eq:XIV.7.16}
\begin{aligned}
&v _{n} (t)  = u _{n} (j / n) + (t - j / n) \cdot n \cdot [ u _{n} ((j + 1) / n) - u _{n} (j / n) ] \\
&\quad (\text{ for } \quad j / n \leqq t \leqq (j + 1) / n; j = 0, 1, \dots , [ n T ] - 1) = u _{n} (t) \quad (\text{ for } \quad n ^{-1} [ n T ] \leqq t \leqq T).
\end{aligned}
\end{equation}

Clearly, $v_{n}(t)$ is strongly differentiable on $[0, T]$ except for finite number of points of the form $t = j / n$, and for $j / n < t < (j + 1) / n$ we have 
\begin{equation}\label{eq:XIV.7.17}
\begin{yoslines}
&\frac {\dd v _{n} (t)}{\dd t}  = n [ u _{n} ((j + 1) / n) - u _{n} (j / n) ] = n [ (I - n ^{-1} A) ^{-1} u _{n} (j / n) - u _{n} (j / n) ] \\
&= A _{1/n} u _{n} (j / n) \quad (\text{ by } \quad A _{\lambda} = \lambda^{-1} (J _{\lambda} - I)) = A _{1/n} u _{n} (t) = A _{1/n} v _{n} (j / n).
\end{yoslines}
\end{equation}

By \eqref{eq:XIV.6.7}, \eqref{eq:XIV.6.8} and \eqref{eq:XIV.6.9} of the preceding Section and by 
\begin{align}
A u _{n} (j / n)& = A J _{1/n} u _{n} ((j - 1 / n) \ni A _{1/n} u _{n} ((j - 1) / n), \text{ we obtain } \notag \\
\left\| A _{1/n} u _{n} (j / n) \right\|& \leqq \left\| A u _{n} (j / n) \right\| \leqq \left\| A _{1/n} u _{n} ((j - 1) / n) \right\| \leqq \left\| A u _{n} ((j - 1) / n) \right\|. \label{eq:XIV.7.18}
\end{align}

Thus, by \eqref{eq:XIV.7.17}, we have 
\begin{equation}\label{eq:XIV.7.19}
\left\| \frac {\dd v _{n} (t)}{\dd t} \right\| = \| A _{1/n} u _{n} (j / n) \| \leqq \| A u _{n} (0) \| = \| A x _{0} \| \text{ for a.e. } t \geqq 0.
\end{equation}

We have incidentally proved from the proof of \eqref{eq:XIV.7.17}-\eqref{eq:XIV.7.19} that, whenever $j / n \leqq t \leqq (j + 1) / n$, 
\[
\begin{aligned} & {\| v _{n} (t) - u _{n} (t) \| = (t - j / n) \cdot n \cdot \| u _{n} ((j + 1) / n) - u _{n} (j / n) \|} \\ & {\qquad \leqq n ^{-1} \| A _{1/n} u _{n} (j / n) \| \leqq n ^{-1} \| A x _{0} \|,} \end{aligned}
\]
and hence 
\begin{equation}\label{eq:XIV.7.20}
\left\| v _{n} (t) - u _{n} (t) \right\| \leqq n ^{-1} \left\| A x _{0} \right\| \quad \text{ for all } \quad t \geqq 0.
\end{equation}

We put $\dd u(t) / \dd t = y(t)$ and $n[u_n(t) - u_n(t - 1/n)] = y_n(t)$. Then, by \eqref{eq:XIV.7.11}, we have $y(t) \in Au(t)$ for almost every $t \geqq 0$. We have also $y_n(t) \in Au_n(t)$ when $t \neq j/n$ ( $j = 0, 1, \ldots$), because 
\[
y _{n} (t) = A _{1/n} u _{n} (t - 1 / n) \in A (I - n ^{-1} A) ^{-1} u _{n} (t - 1 / n) = A u _{n} (t)
\]
by \eqref{eq:XIV.6.8} of the preceding Section~\hyperref[sec:XIV.6]{$6$}. A being a dissipative set, we obtain, by \eqref{eq:XIV.6.5} of the preceding Section~\hyperref[sec:XIV.6]{$6$}, that for almost every $t \geqq 0$ the inequality 
\[
\left\| [ n \cdot u _{n} (t) - y _{n} (t) ] - [ n \cdot u (t) - y (t) ] \right\| \geqq \left\| n \cdot u _{n} (t) - n \cdot u (t) \right\|.
\]

Therefore we have, for almost every $t \geqq 0$, 
\begin{equation}\label{eq:XIV.7.21}
\begin{aligned}
&\left\|\frac{\dd u(t)}{\dd t}-n[u(t)-u(t-1/n)]\right\|\\
&=\bigl\|n[u_n(t)-u_n(t-1/n)]-n[u(t)-u(t-1/n)]+y(t)-y_n(t)\bigr\|\\
&\geqq\bigl\|[n u_n(t)-y_n(t)]-[n u(t)-y(t)]\bigr\|
  -\bigl\|n u_n(t-1/n)-n u(t-1/n)\bigr\|\\
&\geqq n\|u_n(t)-u(t)\|-n\|u_n(t-1/n)-u(t-1/n)\|.
\end{aligned}
\end{equation}

We here extend $u_{n}(t)$ and $u(t)$ for negative value of $t$ by 
\begin{equation}\label{eq:XIV.7.22}
\begin{yoslines}
&u _{n} (t) = x _{0} - n ^{-1} z _{0}, \text{ where } z _{0} \text{ is any fixed point } A x _{0}, u (t) = x _{0}.
\end{yoslines}
\end{equation}

Then, integrating the resulting inequality \eqref{eq:XIV.7.21} on $[0, \theta]$ where $n^{-1} \leqq \theta \leqq T$, we obtain 
\[
\begin{aligned}
-& \int_{-1/n} ^{0} n \| u _{n} (t) - u (t) \| \dd t + \int_{\theta -1 / n} ^{\theta} n \| u _{n} (t) - u (t) \| \dd t
\\
&\leqq \int_{0} ^{\theta} \left\| \frac {\dd u (t)}{\dd t} - n [ u (t) - u (t - 1 / n) ] \right\| \dd t,
\end{aligned}
\]
that is, 
\begin{equation*}
\begin{aligned}
&n \int_{\theta -1 / n} ^{\theta} \| u _{n} (t) - u (t) \| \dd t  \leqq \int_{0} ^{\theta} \left\| \frac {\dd u (t)}{\dd t} - n [ u (t) - u (t - 1 / n) ] \right\| \dd t + n \int_{-1/n} ^{0} \| u _{n} (t) - x _{0} \| \dd t \\
&\leqq \int_{0} ^{\theta} \left\| \frac {\dd u (t)}{\dd t} - n [ u (t) - u (t - 1 / n) ] \right\| \dd t + n ^{-1} \| z _{0} \| . \quad \text{ (by \eqref{eq:XIV.7.22}) }
\end{aligned}
\end{equation*}

Adding these inequalities for $\theta = 1 / n, 2 / n, \ldots, N / n$ with $N = [nT]$ yields 
\[
\begin{aligned}
&n \int_{0} ^{N/n} \| u _{n} (t) - u (t) \| \dd t  \leqq N \int_{0} ^{T} \left\| \frac {\dd u (t)}{\dd t} - n [ u (t) - u (t - 1 / n) ] \right\| \dd t + n ^{-1} N \| z _{0} \|
\end{aligned}
\]
and therefore 
\begin{equation}\label{eq:XIV.7.23}
\begin{aligned}
&\int_{0} ^{N/n} \| u _{n} (t) - u (t) \| \dd t  \leqq T \int_{0} ^{T} \left\| \frac {\dd u (t)}{\dd t} - n [ u (t) - u (t - 1 / n) ] \right\| \dd t + n ^{-1} T \cdot \| z _{0} \| .
\end{aligned}
\end{equation}

Since, by the assumption, $\slim\limits_{n\to \infty}n[u(t) - u(t - 1 / n)] = \dd u(t) / \dd t$ for almost every $t$ on $[0,T]$, we obtain $\lim\limits_{n\to \infty}\int_0^T \|u_n(t) - u(t)\|\dd t = 0$ by \eqref{eq:XIV.7.23} 
and by 
\[
\left\| \frac {\dd u (t)}{\dd t} - n [ u (t) - u (t - 1 / n) ] \right\| \leqq 2 \cdot \| z _{0} \| ,
\]
which is implied by \eqref{eq:XIV.7.14} and $n[u(t) - u(t - 1/n)] = \int_{t-1/n}^{t}\frac{\dd u(s)}{\dd s} \quad \dd s$. Therefore, combined with \eqref{eq:XIV.7.20}, we obtain 
\begin{equation}\label{eq:XIV.7.24}
\lim\limits _{n \rightarrow \infty} \int_{0} ^{T} \| v _{n} (t) - u (t) \| \dd t = 0.
\end{equation}

On the other hand, we have, by \eqref{eq:XIV.7.14}, \eqref{eq:XIV.7.19} and by the same argument which enabled us to prove \eqref{eq:XIV.7.13}, 
\[
2 ^{-1} \frac {d}{\dd t} \| v _{n} (t) - u (t) \| ^{2} \leqq \| v _{n} (t) - u (t) \| \cdot \left\| \frac {\dd v _{n} (t)}{\dd t} - \frac {\dd u (t)}{\dd t} \right\|
\]
$\leqq \|v_n(t) - u(t)\| \cdot 2 \cdot \|z_0\|$ for almost every $t \in [0, T]$. 

Therefore, by $v_{n}(0) = u(0) = x_{0}$, we obtain 
\[
\left\| v _{n} (t) - u (t) \right\| ^{2} \leqq 4 \int_{0} ^{t} \left\| v _{n} (s) - u (s) \right\| \dd s \cdot \left\| z _{0} \right\| \quad \text{ for } \quad 0 \leqq t \leqq T.
\]

Hence, by \eqref{eq:XIV.7.24}, we obtain $\slim\limits_{n\to \infty}v_n(t) = u(t)$ uniformly on $[0,T]$. Thus, by \eqref{eq:XIV.7.20}, we have proved that $\slim\limits_{n\to \infty}u_n(t) = u(t)$ uniformly on $[0,T]$. 
\end{proof}

\begin{theorem}\label{thm:XIV.7.Theorem.3} Let $A$ be a hyperdissipative set $\subseteq X \times X$, and let $x_0 \in D(A)$ and $\lambda > 0$. As in the preceding Section~\hyperref[sec:XIV.6]{$6$}, the initial value problem 
\begin{equation}\label{eq:XIV.7.25}
\frac {\dd u ^{(\lambda)} (t)}{\dd t} = A _{\lambda} u ^{(\lambda)} (t) \quad \text{ for } \quad t \geqq 0 \quad \text{ and } \quad u ^{(\lambda)} (0) = x _{0}
\end{equation}
has a uniquely determined solution $u^{(\lambda)}(t)$, because $A_{\lambda}$ is a Lipschitzian mapping. We have then 
\begin{equation}\label{eq:XIV.7.26}
\begin{aligned}
\slim_{\lambda \downarrow 0} u^{(\lambda)}(t) &\ \text{ exists uniformly on every compact interval of } t \text{ and }
\\
\slim_{\lambda \downarrow 0} u^{(\lambda)}(t) &= \slim_{n \rightarrow \infty} (I - n^{-1} t A)^{-n} x _{0}.
\end{aligned}
\end{equation}
\end{theorem}

\begin{remarku} If the strong solution $u(t)$ of \eqref{eq:XIV.7.11} exists, then, as in the linear case, $u(t)$ can be given by Hille's type approximation $\slim\limits_{n\to \infty}(I - n^{-1}tA)^{-n}\cdot x_0$ as well as by Yosida's type approximation $\slim\limits_{\lambda \downarrow \infty}u^{(\lambda)}(t)$. 
\end{remarku}

\begin{proof}[Proof of Theorem~\ref{thm:XIV.7.Theorem.3}.]
Since $A_{1} = (J_{1} - I)$ is dissipative, we obtain, as in the case of \eqref{eq:XIV.7.15'},
\begin{equation}\label{eq:XIV.7.27}
\left\| u ^{(1)} (t) - u ^{(1)} (0) \right\| \leqq t \cdot \left\| A _{1} u ^{(1)} (0) \right\| = t \cdot \left\| A _{1} x _{0} \right\| = t \cdot \left\| J _{1} x _{0} - x _{0} \right\|.
\end{equation}

On the other hand, it is easy to verify that the solution $u^{(1)}(t)$ of \eqref{eq:XIV.7.25} with $\lambda = 1$ is given by 
\begin{equation}\label{eq:XIV.7.28}
u ^{(1)} (t) = e ^{-t} x _{0} + \int_{0} ^{t} e ^{s-t} J _{1} u ^{(1)} (s) \dd s.
\end{equation}

Hence 
\[
u ^{(1)} (t) - J _{1} ^{n} x _{0} = e ^{-t} \left(x _{0} - J _{1} ^{n} x _{0}\right) + \int_{0} ^{t} e ^{s-t} \left[ J _{1} u ^{(1)} (s) - J _{1} ^{n} x _{0} \right] \dd s
\]
and so, by the fact that $J_{1}$ is a contraction mapping, we obtain 
\begin{equation}\label{eq:XIV.7.29}
\left\| u ^{(1)} (t) - J _{1} ^{n} x _{0} \right\| \leqq n e ^{-t} \left\| \left(J _{1} - I\right) x _{0} \right\| + \int_{0} ^{t} e ^{s-t} \left\| u ^{(1)} (s) - J _{1} ^{n-1} x _{0} \right\| \dd s.
\end{equation}

From this we shall derive 
\begin{equation}\label{eq:XIV.7.30}
\left\| u ^{(1)} (n) - J _{1} ^{n} x _{0} \right\| \leqq \sqrt {n} \left\| (J _{1} - I) x _{0} \right\|.
\end{equation}

We may assume that $(J_{1}-I)x_{0}=A_{1}x_{0}\neq0$. For, if otherwise, the solution $u^{(1)}(t)$ of \eqref{eq:XIV.7.25} with $\lambda=1$ is given by $u^{(1)}(t)\equiv x_{0}$ and hence \eqref{eq:XIV.7.30} would be trivial. The assumption $(J_{1}-I)x_{0}\neq0$, combined with \eqref{eq:XIV.7.29}, gives 
\begin{equation}\varphi_{n} (t) \leqq n e ^{-t} + \int_{0} ^{t} e ^{s-t} \varphi_{n-1} (s) \dd s, \text{ where } \varphi_{n} (t) = \left\| u ^{(1)} (t) - J _{1} ^{n} x _{0} \right\| \cdot \left\| \left(J _{1} - I\right) x _{0} \right\| ^{-1}. \label{eq:XIV.7.31}\end{equation}

We shall prove, by induction with respect to $n$, that 
\begin{equation}\label{eq:XIV.7.32}
\varphi_{n} (t) \leqq \{(n - t) ^{2} + t \} ^{1/2}.
\end{equation}

This gives \eqref{eq:XIV.7.30} by putting $t = n$. Inequality \eqref{eq:XIV.7.32} is clear for $n = 0$. If \eqref{eq:XIV.7.32} 
is true for $(n - 1)$, then $\varphi_n(t) \leqq n e^{-t} + \int_0^t e^{s-t}\{ (n - 1 - s\}^2 + s\}^{1/2} \dd s$. Denoting 
the right hand side by $\psi_{n}(t)e^{-t}$, we have to show that 
\[
\psi_{n} (t) \leqq e ^{t} \left\{(n - t) ^{2} + t \right\} ^{1/2}.
\]

Since $\psi_{n}(0) = n$, it is sufficient to prove that 
\[
\psi_{n} ^{\prime} (t) = e ^{t} \left\{(n - 1 - t) ^{2} + t \right\} ^{1/2} \leqq \frac {d}{\dd t} e ^{t} \left\{(n - t) ^{2} + t \right\} ^{1/2}.
\]

The last inequality is easily checked and so we have proved \eqref{eq:XIV.7.30}. 

We are now prepared to prove \eqref{eq:XIV.7.26}. Put $v^{(\lambda)}(t) = u^{(\lambda)}(\lambda t)$. Then, by $A_{\lambda} = \lambda^{-1}(J_{\lambda} - I)$, we have 
\begin{equation}\tag{25'}\label{eq:XIV.7.25'}
\frac {\dd v ^{(\lambda)} (t)}{\dd t} = \left(J _{\lambda} - I\right) v ^{(\lambda)} \text{ for all } t \geqq 0 \text{ with } v ^{(\lambda)} (0) = x _{0}.
\end{equation}

Taking $J_{\lambda}$ for $J_{1}$ and $v^{(\lambda)}$ for $u^{(\lambda)}$ in \eqref{eq:XIV.7.30} and making use of \eqref{eq:XIV.6.9} of the preceding Section~\hyperref[sec:XIV.6]{$6$}, we obtain 
\begin{equation}\label{eq:XIV.7.33}
\begin{aligned}
&\| v ^{(\lambda)} (n) - J _{\lambda} ^{n} x _{0} \|  = \| u ^{(\lambda)} (n \lambda) - J _{\lambda} ^{n} x _{0} \| \leqq \sqrt {n} \| (J _{\lambda} - I) x _{0} \| \leqq \sqrt {n} \lambda \| A _{\lambda} x _{0} \| \leqq \sqrt {n} \lambda \| A x _{0} \|.
\end{aligned}
\end{equation}

Let $\lambda_{k} \downarrow 0$ and let $n_{k} = [t / \lambda_{k}]$ so that $n_{k} \lambda_{k} \uparrow t$. Hence, by \eqref{eq:XIV.7.1'} and \eqref{eq:XIV.7.10}, we obtain $\slim\limits_{k \to \infty} J_{\lambda_k}^{n_k} x_0 = \slim\limits_{n \to \infty} (I - n^{-1} tA)^{-n} x_0$. Therefore, combined with \eqref{eq:XIV.7.33}, we have 
\[
\left\| u ^ {\left(\lambda_{k}\right)} \left(n _{k} \lambda_{k}\right) - J _ {\lambda_{k}} ^ {n _{k}} x _{0} \right\| \leqq \sqrt {n _{k}} \lambda_{k} \left\| A x _{0} \right\| \leqq n _{k} ^{-1/2} t \left\| A x _{0} \right\|.
\]

Since $\{\lambda_k\}$ was arbitrary, we have thus proved \eqref{eq:XIV.7.26}. 
\end{proof}

\begin{remarku} The above proof is adapted from H. \textsc{Brezis}-A. \textsc{Pazy} \cite{brezis29}. It is to be noted here that inequality \eqref{eq:XIV.7.30} is due to I. \textsc{Miyadera}-S. \textsc{Oharu}\cite{miyadera204}. 
\end{remarku}

\begin{proof}[Proof of the assertion ii) $\rightarrow$ i) in Theorem~\ref{thm:XIV.7.Theorem.2}.]
Let $u(t) = s\text{- }\lim\limits_{n\to \infty}(I - n^{-1}tA)^{-n}x_0$ be strongly differentiable at $t_0 > 0$ with strong derivative $y = \dd u(t) / \dd t|_{t=t_0}$ so that
\begin{equation}\label{eq:XIV.7.34}
u \left(t _{0} + h\right) = u \left(t _{0}\right) + h y + o (h), \text{ where } s \underset {h \downarrow 0} {\lim} o (h) / h = 0.
\end{equation}

Since \eqref{eq:XIV.7.19} and \eqref{eq:XIV.7.20} implies that $t \to u(t)$ is a Lipschitzian mapping with Lipschitz constant $\|Ax_0\|$, we obtain i) if we can prove 
\begin{equation}\label{eq:XIV.7.35}
\left\{u \left(t _{0}\right), y \right\} \in A.
\end{equation}

The proof of \eqref{eq:XIV.7.35} reads as follows. Let $0 < \lambda < t_{0}$. Then, by the strong differentiability of $u(t)$ at $t = t_{0}$ and by the hyperdissipative property of $A$, there exists a uniquely determined point $\{x_{\lambda}, y_{\lambda}\} \in A$ such that 
\begin{equation}\label{eq:XIV.7.36}
x _{\lambda} - \lambda y _{\lambda} = u (t _{0} - \lambda) = u (t _{0}) - \lambda y + o (\lambda).
\end{equation}

We obtain \eqref{eq:XIV.7.35} if we can show 
\begin{equation}\label{eq:XIV.7.37}
\slim _{\lambda \downarrow 0} x _{\lambda} = u (t _{0}) \quad \text{ and } \quad \slim _{\lambda \downarrow 0} y _{\lambda} = y,
\end{equation}
because $A$ is a closed set in $X \times X$ by the assumption. For the proof of \eqref{eq:XIV.7.37}, we put 
\begin{equation}\label{eq:XIV.7.38}
\parbox{0.92\linewidth}{\everymath{}\let\limits\nolimits $\hat{x}_{\lambda} = \hat{x} - \lambda \hat{y}$, where $\{\hat{x}, \hat{y}\}$ is a point of $A$ to be determined later.}
\end{equation} 


Then, by \eqref{eq:XIV.7.13} and \eqref{eq:XIV.7.25}, we obtain a certain $f \in F(u^{(\lambda)}(t) - \hat{x}_{\lambda})$ such that 
\begin{align}
2 ^{-1} \frac {d}{\dd t} \| u ^{(\lambda)} (t) - \hat {x} _{\lambda} \| ^{2}& = \operatorname{Re} \left\langle A _{\lambda} u ^{(\lambda)} (t), f \right\rangle = \operatorname{Re} \left\langle A _{\lambda} u ^{(\lambda)} (t) - A _{\lambda} \hat {x} _{\lambda}, f \right\rangle + \notag \\
+ \operatorname{Re} \left\langle A _{\lambda} \hat {x} _{\lambda}, f \right\rangle& \leqq \operatorname{Re} \left\langle \hat {y}, f \right\rangle , \notag
\end{align}
because $A_{\lambda} \hat{x}_{\lambda} = \hat{y}$ and, moreover, $A_{\lambda}$ is dissipative by Proposition~\hyperref[prop:XIV.6.Proposition.1]{$1$} in the preceding Section~\hyperref[sec:XIV.6]{$6$}. Thus 
\begin{equation}\label{eq:XIV.7.39}
\left\| u ^{(\lambda)} \left(t _{0} + h\right) - \hat {x} _{\lambda} \right\| ^{2} - \left\| u ^{(\lambda)} \left(t _{0}\right) - \hat {x} _{\lambda} \right\| ^{2} \leqq 2 \int_{0} ^{h} \left\langle \hat {y}, u ^{(\lambda)} \left(t _{0} + \tau\right) - \hat {x} _{\lambda} \right\rangle_{s} \cdot \dd \tau ,
\end{equation}
where 
\begin{equation}\label{eq:XIV.7.40}
\left\langle w, z \right\rangle_{s} = \sup\limits _{k \in F (z)} \operatorname{Re} \left\langle w, k \right\rangle .
\end{equation}

It will be proved later that 
\begin{equation}
\label{eq:XIV.7.41}
\begin{minipage}{0.92\linewidth}
\everymath{}\let\limits\nolimits
$\langle w,z\rangle_s $ is upper semi-continuous, i.e., $\slim_{n\to\infty}w_n=w_\infty$ and $\slim_{n\to\infty}z_n =z_\infty$ implies that $\limsup\limits_{n\to\infty}\langle w_n,z_n\rangle_s\leqq\langle w_\infty,z_\infty\rangle_s.$
\end{minipage}
\end{equation} 

Hence, by $\slim\limits_{\lambda \downarrow \infty} u^{(\lambda)}(t) = u(t) = \slim\limits_{n \to \infty} (I - n^{-1}tA)^{-n} x_0$ proved in \eqref{eq:XIV.7.26} and by $\slim\limits_{n \to \infty} \hat{x}_\lambda = \hat{x}$, we obtain 
\begin{equation}\label{eq:XIV.7.42}
\left\| u \left(t _{0} + h\right) - \hat {x} \right\| ^{2} - \left\| u \left(t _{0}\right) - \hat {x} \right\| ^{2} \leqq 2 \int_{0} ^{h} \left\langle \hat {y}, u \left(t _{0} + \tau\right) - \hat {x} \right\rangle_{s} \cdot \dd \tau .
\end{equation}

For every $f \in F(u(t_0) - \hat{x})$, we have 
\begin{gather}\operatorname{Re} \left\langle u (t _{0} + h) - u (t _{0}), f \right\rangle - \operatorname{Re} \left\langle \hat {x} - u (t _{0}), f \right\rangle = {} \operatorname{Re} \left\langle u (t _{0} + h) - \hat {x}, f \right\rangle \notag \\\leqq {} \left\| u \left(t _{0} + h\right) - \hat {x} \right\| \cdot \left\| \hat {x} - u \left(t _{0}\right) \right\| \leqq {} 2 ^{-1} \left\{\left\| u \left(t _{0} + h\right) - \hat {x} \right\| ^{2} + \left\| \hat {x} - u \left(t _{0}\right) \right\| ^{2} \right\} \notag
\end{gather}

Therefore, by $\operatorname{Re}\langle \hat{x} - u(t_0), f \rangle = -\|\hat{x} - u(t_0)\|^2$ and \eqref{eq:XIV.7.42}, we obtain 
\[
2 \cdot \operatorname{Re} \left\langle \frac {u (t _{0} + h) - u (t _{0})}{h}, f \right\rangle \leqq 2 \int_{0} ^{1} \left\langle \hat {y}, u (t _{0} + h \tau) - \hat {x} \right\rangle_{s} \cdot \dd \tau .
\]

Hence, by the strong differentiability of $u(t)$ at $t = t_0$ and \eqref{eq:XIV.7.41}, we have 
\begin{equation}\tag{42'}\label{eq:XIV.7.42'}
\operatorname{Re} \langle y, f \rangle \leqq \langle \hat {y}, u (t _{0}) - \hat {x} \rangle_{s}.
\end{equation}

We may replace the supremum on the right hand side of \eqref{eq:XIV.7.42'} by the maximum, because the set $F(u(t_0) - \hat{x})$ contained in the strongly closed sphere of $X'$ is $w^*$-compact in $X'$ (see Theorem~\hyperref[thm:AV.1.Theorem.1]{$1$} in Appendix to Chapter~\ref{chap:V}, \ref{sec:AV.1}). Thus there exists an $\hat{f} \in F(u(t_0) - \hat{x})$ such that 
\[
\operatorname{Re} \langle y, f _{1} \rangle \leqq \operatorname{Re} \langle \hat {y}, \hat {f} \rangle \quad \text{ for all } \quad f _{1} \in F (u (t _{0}) - \hat {x}).
\]

By taking $\hat{x} = x_{\lambda},\hat{y} = y_{\lambda}$ and $f_{1} = f$, we obtain 
\begin{equation}\label{eq:XIV.7.43}
\operatorname{Re} \langle y, f \rangle \leqq \operatorname{Re} \langle y _{\lambda}, f \rangle , \text{ that is, } \operatorname{Re} \langle y _{\lambda} - y, f \rangle \geqq 0.
\end{equation}

On the other hand, we have, from \eqref{eq:XIV.7.36}, 
\begin{equation}\label{eq:XIV.7.44}
\frac {u \left(t _{o}\right) - x _{\lambda}}{\lambda} + \frac {o (\lambda)}{\lambda} = y - y _{\lambda},
\end{equation}
and hence, by \eqref{eq:XIV.7.43}, we obtain 
\[
\left\| u \left(t _{0}\right) - x _{\lambda} \right\| ^{2} = \operatorname{Re} \left\langle u \left(t _{0}\right) - x _{\lambda}, \hat {f} \right\rangle \leqq \operatorname{Re} \left\langle o (\lambda), \hat {f} \right\rangle \leqq o (\lambda) \cdot \left\| u \left(t _{0}\right) - x _{\lambda} \right\|.
\]

Hence $\|u(t_0) - x_{\lambda}\| \leqq o(\lambda)$ and so we obtain \eqref{eq:XIV.7.37} by \eqref{eq:XIV.7.44}. 

Finally, the proof of \eqref{eq:XIV.7.41} is given as follows. Since the set $\{f\in X';\|f\|\leqq \|z_m\|\}$ is $w^*$-compact (Theorem~\ref{thm:AV.1.Theorem.1} in Appendix to Chapter~\ref{chap:V}, \ref{sec:AV.1}), we can find $f_{m}\in X'$ such that 
\[
\left\langle w _{m}, z _{m} \right\rangle_{s} = \operatorname{Re} \left\langle w _{m}, f _{m} \right\rangle \quad \text{ and } \quad \| f _{m} \| \leqq \| z _{m} \|.
\]

We may assume, without loss of generality, that $\lim\limits_{m\to \infty}\langle w_m,z_m\rangle_s$ exists. 

Let $f_{\infty} \in X'$ be any one of the $w^*$-accumulation points of $\{f_m\}$. Then we have 
\[
\lim\limits _{m \rightarrow \infty} \left\langle w _{m}, z _{m} \right\rangle_{s} = \lim\limits _{m \rightarrow \infty} \operatorname{Re} \left\langle w _{m}, f _{m} \right\rangle = \operatorname{Re} \left\langle w _{\infty}, f _{\infty} \right\rangle , \| f _{\infty} \| \leqq \| z _{\infty} \|,
\]
because $s$ - $\lim\limits_{m \to \infty} w_m = w_\infty$ and $w^*$ - $\lim\limits_{m' \to \infty} f_{m'} = f_\infty$ for a suitable subsequence 
$\{m'\}$ of the sequence $\{m\}$. Therefore we must have \eqref{eq:XIV.7.41}. 
\end{proof}

\begin{remarku} It is an unfortunate and a fortiori an interesting situation that, although $\slim\limits_{n\to \infty}(I - n^{-1}tA)^{-n}x_0$ and $\slim\limits_{\lambda \downarrow 0}u^{(\lambda)}(t)$ both exist and coincide each other under a general condition, this limit need not be strongly differentiable in $t$ anywhere. Such an example is given in M. \textsc{Crandall}-T. \textsc{Liggett} \cite{clarkson37}. Recently, P. \textsc{Bénilan} \cite{benilan13} devised the following remedy for the above situation. Let $u(t)$ be a strong solution of initial value problem \eqref{eq:XIV.7.11}. Let, further, $\{z,w\}$ be any fixed point of the hyperdissipative subset $A$ of $X\times X$. Then, as in the case of \eqref{eq:XIV.7.42}, we obtain 
\[
\|u(t)-z\|^2\leqq\|u(s)-z\|^2+2\int_s^t\langle w,u(\tau)-z\rangle_s\cdot\dd\tau
\]
for all $0\leqq s\leqq t$ and with a fixed initial condition $u(0)=x_0\in D(A)$.
\end{remarku}

P. \textsc{Bénilan} \cite{benilan13} called a strongly continuous $X$-valued function $u(t)$ as an integral solution of initial value problem \eqref{eq:XIV.7.11}, if $u(t)$ satisfies the above condition for all $\{z, w\} \in A$ simultaneously, and he proved 

\begin{theorem}\label{thm:XIV.7.Theorem.4} Let $A$ be a hyperdissipative set $\subseteq X \times X$. Then $u(t) = s\text{- }\lim\limits_{\lambda \downarrow 0} u^{(\lambda)}(t) = s\text{- }\lim\limits_{n \to \infty} (I - n^{-1}tA)^{-n} \cdot x_0$ is the uniquely determined integral solution of \eqref{eq:XIV.7.11}. 
\end{theorem}

For the proof of this theorem and further interesting extensions as well as comments, we refer the reader to P. \textsc{Bénilan} \cite{benilan14}. 

\begin{remarku} The scope of applications of Theorems in the present and the preceding Sections may well be seen, e.g., by reading S. \textsc{Aizawa} \cite{aizawa4}, M. G. \textsc{Crandall} \cite{clarkson38}, Y. \textsc{Konishi} \cite{konishi160} and \cite{konishi161}, and B. K. \textsc{Quinn} \cite{quinn265}. 
\end{remarku}

\yoschapterstar{Supplementary Notes}\label{chap:SupplementaryNotes}

\yossub{Chapter I and Chapter VI}\phantomsection\label{notes:I-VI}

\begin{enumerate}[1.]
\item For the distributions or the generalized functions due to S. L. \textsc{Sobolev}, L. \textsc{Schwartz} and I. M. \textsc{Gelfand}, see the new and enlarged edition L. \textsc{Schwartz} \cite{schwartz299} and the English edition I. M. \textsc{Gelfand} \cite{gelfand80}--\cite{gelfand84}.
\item For the recent development of the hyperfunctions due to M. \textsc{Sato}, see M. \textsc{Sato} \cite{sato285}--\cite{sato286} and P. \textsc{Schapira} \cite{schapira287}.
\end{enumerate}

\yossub{Chapter VI}\phantomsection\label{notes:VI}

\begin{enumerate}[1.]
\item (Section~\hyperref[sec:VI.7]{$7$}) Sobolev Spaces. The theorem (Sobolev's Lemma) given on p.~\pageref{thm:VI.7.Sobolev} is a very special case of the so-called ``Sobolev Imbedding Theorems'' in S. L. \textsc{Sobolev} \cite{sobolev306}--\cite{sobolev307}. A comprehensive description of the development of this imbedding theorems is given in R. A. \textsc{Adams} \cite{adams2}. See also E. M. \textsc{Stein} \cite{stein310}.
\end{enumerate}

\yossub{Chapter X}\phantomsection\label{notes:X}

\begin{enumerate}[1.]
\item (Section~\hyperref[sec:X.1]{$1$}) The trace operator or the generalized boundary values. We can prove 

\begin{theoremu} Let $\Omega$ be a bounded open domain in $\mathbb{R}^{n}$ such that its boundary hypersurface $\partial\Omega$ is $C^2$. Then, for any $f(x) \in W^{1,2}(\Omega)$ and for almost all $\xi \in \partial\Omega$, the limiting value $\varphi(\xi) = \lim\limits_{x \to \xi} f(x)$ exists if we let $x$ tend to $\xi$ along the normal at the point $\xi \in \partial\Omega$. Moreover, this boundary value $\varphi(\xi)$ of $f(x)$ belongs to $L^2(\partial\Omega)$ and satisfies the inequality: 
\[
\left\| \varphi (\xi) \right\| _ {L^{2} (\partial \Omega)} \leqq C \left\{\left\| f (x) \right\| _ {L^{2} (\Omega)} + \sum\limits_{j=1} ^{n} \left\| \partial f / \partial x _{j} \right\| _ {L^{2} (\Omega)} \right\},
\]
where the positive constant $C$ is independent of the individual choice of $f$ and the notation $\partial f/\partial x_{j}$ means differentiation in the sense of the ``distribution''. 
\end{theoremu}

The linear map $f \to \varphi$ is called the trace operator. For the proof as well as generalizations of the above theorem, see, e.g., S. \textsc{Mizohata} \cite{mizohata213}, R. A. \textsc{Adams} \cite{adams2} and F. \textsc{Trèves} \cite{trotter326}. 

\item \phantomsection\label{notes:AX} (Appendix to Chapter~\ref{chap:X}) For extensions with a unified treatment of the result of R. A. \textsc{Minlos} \cite{minlos202}, see L. \textsc{Schwartz} \cite{schwartz300} and the bibliography cited there.
\end{enumerate}

\yossub{Chapter XI}\phantomsection\label{notes:XI}

\begin{enumerate}[1.]
\item For recent development of the theory of contraction operators in Hilbert spaces, see B. Sz. \textsc{Nagy}-C. \textsc{Foias} \cite{nagy220}.
\end{enumerate}

\yossub{Chapter XII}\phantomsection\label{notes:XII}

\begin{enumerate}[1.]
\item (Section~\hyperref[sec:XII.1]{$1$}) For Choquet's refinement of the Krein-Milman Theorem, see R. R. \textsc{Phelps} \cite{phelps250}.
\end{enumerate}

\yossub{Chapter XIII}\phantomsection\label{notes:XIII}
\setcounter{equation}{18}

\begin{enumerate}[1.]
\item (Section~\hyperref[sec:XIII.1]{$1$}) We can prove (K. \textsc{Yosida} \cite{yosida360}) 

\begin{theoremn}[Theorem~\ref{thm:XIII.1.Theorem.1}']\yosnamedtag{1'}\label{thm:notes.XIII.1.1prime} Let $\mathrm{P}(t, x, E)$ be a Markov process on the phase space $(S, \mathfrak{B})$ with an invariant measure $m$ such that $m(S) = 1$. Then, for any $f \in L^1(S, \mathfrak{B}, m)$ and $t > 0$, there corresponds an $f_{-t} \in L^1(S, \mathfrak{B}, m)$ with $\|f_{-t}\|_1 \leqq \|f\|_1$ in such a way that 
\begin{equation}\label{eq:XIV.7.19a}
\int_{S} m (\dd x) f (x) P (t, x, E) = \int_ {\boldsymbol {E}} f _{-t} (x) m (\dd x) \quad \text{ for all } \quad E \in \mathfrak {B}.
\end{equation}
\end{theoremn}

\begin{proof}
 Since $P(t, x, E)$ is a bounded function, $\int_{S} m(\dd x) f(x)$ $P(t, x, E)$ does exist for any $f \in L^{1}(S, \mathfrak{B}, m)$ and we have 
\begin{equation}\label{eq:XIV.7.20a}
\sup\limits _ {\boldsymbol {E}} | \int_{S} m (\dd x) f (x) P (t, x, E) | \leqq \sup\limits _ {\boldsymbol {E}} \int_{S} m (\dd x) | f (x) | P (t, x, E) \leqq \| f \| _ {\mathbf {1}}. \end{equation}

If $f \in L^{1}(S, \mathfrak{B}, m)$ is $\in L^{\infty}(S, \mathfrak{B}, m)$, then, by \eqref{eq:XIII.1.8}, the $\sigma$-additive function $\int_{S} m(\dd x)f(x)P(t, x, E)$ of $E \in \mathfrak{B}$ is $m$-absolutely continuous such that $|\int_{S} m(\dd x)f(x)P(t, x, E)| \leqq \|f\|_{\mathbb{L}^{\infty}} \cdot m(E)$, where $\|f\|_{\mathbb{L}^{\infty}} = \text{ essential sup } |f(x)|$. Next, let $f \in L^{1}(S, \mathfrak{B}, m)$ be real-valued non-negative and put $f^{(n)}(x) = \min(f(x), n)$. Since 
\[
\int_{S} m (\dd x) f (x) P (t, x, E) = \lim\limits _{n \rightarrow \infty} \int_{S} m (\dd x) f ^{(n)} (x) P (t, x, E) \quad \text{ for every } E \in \mathfrak {B}
\]
the left hand side above is $m$-absolutely continuous by the Vitali-Hahn-Saks theorem on p.~\pageref{thm:II.2.VitaliHahnSaks}. Thus, in this special case for $f$, we can define $f_{-t} \in L^{1}(S, \mathfrak{B}, m)$ by \eqref{eq:XIV.7.19a}, and we have $\|f_{-t}\|_{1} \leqq \|f\|_{1}$ by \eqref{eq:XIV.7.20a}. This proves our Theorem~\ref{thm:notes.XIII.1.1prime}. 
\end{proof}

\begin{remarku} Consider the deterministic case given by $P(t, x, E) = C_E(y_t(x))$, defined through a one-parameter family $\{y_t(x); -\infty < t < \infty\}$ of one-one transformation $x \to y_t(x)$ of $S$ onto $S$ such that $m(E) = m(y_t \cdot E)$. Then we see that the mapping $f \to f_t$ in Theorem~\ref{thm:XIII.1.Theorem.1} (p.~\pageref{thm:XIII.1.Theorem.1}) and the mapping $f \to f_{-t}$ in Theorem~\ref{thm:notes.XIII.1.1prime} are respectively given by $f_t(x) = f(y_t(x))$. 
and $f_{-t}(x) = f(y_{-t}(x))$. In this sense, a Markov process with an invariant measure $m$ is ``time reversible''.
\end{remarku}
\item (Section~\hyperref[sec:XIII.1]{$1$}) As a corollary of Theorem~\ref{thm:notes.XIII.1.1prime} above, we obtain, as in the case of Theorem~\ref{thm:XIII.1.Theorem.2} (p.~\pageref{thm:XIII.1.Theorem.2}), the following 

\begin{theoremn}[Theorem~\ref{thm:XIII.1.Theorem.2}' (K. \textsc{Yosida} {\cite{yosida360}})]\yosnamedtag{2'}\label{thm:notes.XIII.1.2prime} Under the assumption as in Theorem~\ref{thm:notes.XIII.1.1prime}, the following mean ergodic theorem holds good: 
\begin{align}
\underset {n \rightarrow \infty} {\slim} n ^{-1} \sum\limits_{k=1} ^{n} f _{-k}& = f ^{-*} \text{ exists in } L ^{1} (S, \mathfrak {B}, m), \label{eq:XIV.7.21a} \\
&\int_{S} f (s) m (\dd s) = \int_{S} f ^{-*} (s) m (\dd s). \label{eq:XIV.7.22b}
\end{align}
\end{theoremn}
\end{enumerate}

\yossub{Chapter XIV}\phantomsection\label{notes:XIV}

\begin{enumerate}[1.]
\item (Section~\hyperref[sec:XIV.4]{$4$}) We can prove 

\begin{theoremn}[Theorem~\ref{thm:XIV.4.Theorem.1}']\yosnamedtag{1'}\label{thm:notes.XIV.4.1prime} Under the assumption $x(0) = y \in D(A(0))$ and \eqref{eq:XIV.4.3} (p.~\pageref{eq:XIV.4.3}), the solution of \eqref{eq:XIV.4.1'} (p.~\pageref{eq:XIV.4.1'}), if it exists, is uniquely determined.
\end{theoremn}

\begin{proof}
 We follow an argument given in T. \textsc{Kato} \cite{kato136}. Thus we have, for $\delta > 0$, 
\[ \begin{aligned}
&x (t + \delta)  = x (t) + \delta A (t) x (t) + o (\delta) = (I + \delta A (t)) (I - \delta A (t)) (I - \delta A (t)) ^{-1} x (t) + o (\delta) \\
&= (I - \delta A (t)) ^{-1} x (t) - \delta^{2} A (t) (I - \delta A (t)) ^{-1} A (t) x (t) + o (\delta) \\
&= (I - \delta A (t)) ^{-1} x (t) - \delta ((I - \delta A (t)) ^{-1} - I) A (t) x (t) + o (\delta)
\end{aligned}
\]

Because of \eqref{eq:XIV.4.3}, we have $s\text{- }\lim\limits_{\delta \to 0} (I - \delta A(t))^{-1} z = z$ for any $z \in X$ [see \eqref{eq:IX.7.2} in Chapter~\ref{chap:IX}, \ref{sec:IX.7}]. Hence, again by \eqref{eq:XIV.4.3}, we have 
\[
\left\| x (t + \delta) \right\| \leqq \left\| x (t) \right\| + o (\delta)
\]
and so $d^{+}\|x(t)\| / \dd t \leqq 0$ which implies that $\|x(t)\| \leqq \|x(0)\|$. 

Therefore, the difference between two solutions of \eqref{eq:XIV.4.1'} each with the same initial condition $\in D(A(0))$ must be $0$. 
\end{proof}
\end{enumerate}

\yossub{Chapter IX and Chapter XIV}\phantomsection\label{notes:IX-XIV}

\begin{enumerate}[1.]
\item For Linear evolution equations, see S. \textsc{Krein} \cite{krein168}, P. \textsc{Lax}-R. S. \textsc{Phillips} \cite{lewy174}, J. L. \textsc{Lions} \cite{liggett182} and F. \textsc{Trèves} \cite{trotter326}.
\item For Linear as well as nonlinear evolution equations, see: V. \textsc{Barbu} \cite{barbu11}, F. \textsc{Browder} \cite{browder31}, J. L. \textsc{Lions} \cite{liggett182}, R. H. \textsc{Martin} \cite{martin188}, K. \textsc{Masuda} \cite{masuda190}, I. \textsc{Miyadera} \cite{miyadera207} and H. \textsc{Tanabe} \cite{tanabe319}. For the last cited book of Tanabe, an English translation is in preparation.
\item Currently, convex analysis is playing an important rôle in nonlinear evolution equations. For this analysis, see, e.g., R. S. \textsc{Rockafellar} \cite{rockafeller278} and I. \textsc{Ekeland}-R. \textsc{Temam} \cite{ekeland54}.
\end{enumerate}

